% Lesson 10 — Grammar: Conditionals (1, 2, 3) — Anglais 1ère A — Cours signé OnBuch+
% Programme officiel MINESEC. Compiler avec : tectonic ENA10-grammar-conditionals-1-2-3.tex
\newcommand{\DOCMATIERE}{Anglais}
\newcommand{\DOCNIVEAU}{1ère A}
\newcommand{\DOCMODULE}{Chapter 1 : Family and Social Life}
\newcommand{\DOCLECON}{ENA10}
\newcommand{\DOCTITRE}{Lesson 10 — Grammar: Conditionals (1, 2, 3)}
% OnBuch+ — préambule « cours signés » ENRICHI (compilé avec tectonic / XeLaTeX)
% Système visuel « Pop » d'OnBuch+ : fond crème (jamais blanc), bordures encre
% épaisses, ombres « dures » décalées, titres Archivo Black en capitales.
% Version MATHÉMATIQUES : graphes (pgfplots/tikz), boîtes pédagogiques
% (définition, propriété, méthode, attention, le savais-tu, exercice résolu,
% exercice, TP, bilan), page de garde et signature OnBuch+ ancrée en bas.
%
% Macros attendues dans main.tex AVANT \input{preamble} :
%   \DOCMATIERE  \DOCNIVEAU  \DOCMODULE  \DOCLECON (numéro)  \DOCTITRE
\documentclass[11pt]{article}

\usepackage{fontspec}
\usepackage[english,french]{babel} % français = langue principale ; l'anglais via \eng{..} / textebox
\usepackage[a4paper,margin=2cm,top=3.4cm,bottom=2.8cm,headheight=1.4cm,footskip=1.3cm]{geometry}
\usepackage{amsmath,amssymb,amsfonts}
\usepackage{mathtools} % \xmapsto, \coloneqq... souvent utilisés par les modèles
\usepackage{cancel} % simplifications barrées (bilans de forces)
% \abs{z} (module, valeur absolue) et \norm{u} (norme), utilisés par les modèles
\providecommand{\abs}{}\let\abs\relax\DeclarePairedDelimiter{\abs}{\lvert}{\rvert}
\providecommand{\norm}{}\let\norm\relax\DeclarePairedDelimiter{\norm}{\lVert}{\rVert}
\usepackage[table]{xcolor} % [table] : fournit aussi \rowcolors (alternance de lignes)
\usepackage{pagecolor}
\usepackage{tikz}
\usetikzlibrary{arrows.meta,positioning,calc,shapes.geometric,shapes.misc,shapes.arrows,decorations.pathmorphing,decorations.pathreplacing,decorations.markings,patterns,shadows,mindmap,trees,fit,backgrounds,matrix,intersections,angles,quotes}
\usepackage{pgfplots}
\usepackage[version=4]{mhchem} % équations nucléaires \ce{^{235}_{92}U}
\usepackage[european,siunitx]{circuitikz} % schémas électriques
\pgfplotsset{compat=1.18}
\usepgfplotslibrary{fillbetween,groupplots} % aires entre courbes (calcul intégral)
\usepackage{enumitem}
\usepackage{fancyhdr}
% [most] charge toutes les bibliothèques tcolorbox (listings, minted, raster,
% documentation...) et gonfle inutilement la mémoire TeX sur un document long
% (~300 boîtes) : on ne charge que ce qui est réellement utilisé.
\usepackage[skins,breakable]{tcolorbox}
\usepackage{graphicx}
\usepackage{colortbl}
\usepackage{booktabs}
\usepackage{tabularx}
\usepackage{diagbox} % case d'en-tête diagonale des tableaux à double entrée
% Colonnes extensibles centrée / gauche / droite (C, L, R), souvent utilisées par les modèles
\newcolumntype{C}{>{\centering\arraybackslash}X}
\newcolumntype{L}{>{\raggedright\arraybackslash}X}
\newcolumntype{R}{>{\raggedleft\arraybackslash}X}
\usepackage{array}
\usepackage{multicol}
\usepackage{multirow}
\usepackage{siunitx}
% parse-numbers=false : accepte aussi \SI{2,5 \times 10^{-3}}{...}
\sisetup{locale=FR,output-decimal-marker={,},group-minimum-digits=4,parse-numbers=false}
\DeclareSIUnit\ohm{\ensuremath{\symup{\Omega}}} % U+2126 absent de la police
\DeclareSIUnit\division{div} % réglages d'oscilloscope (ms/div, V/div)
\DeclareSIUnit\tr{tr} % tours (vitesse de rotation, ex. \SI{1500}{\tr\per\minute})
\DeclareSIUnit\molar{M} % concentration molaire (ex. \SI{3}{\molar}, \SI{1}{\micro\molar})
\DeclareSIUnit\an{an} % année (le modèle confond parfois avec \year, un compteur TeX) — ex. \SI{5}{\kilo\meter\per\an}
\DeclareSIUnit\FCFA{FCFA} % franc CFA (ex. \SI{500}{\FCFA\per\kilo\gram})
\DeclareSIUnit\degreeBrix{\textdegree Brix} % teneur en sucre d'un sirop/jus (réfractométrie)
\DeclareSIUnit\month{mois} % le modèle utilise parfois \month, un compteur TeX, comme unité de durée
\usepackage{qrcode}
% \degree : commande fréquemment utilisée par les modèles pour un angle ou une
% température en dehors de \SI{}{} (non fournie par les paquets chargés).
\providecommand{\degree}{^{\circ}}
% \qed (fin de démonstration), utilisé par les modèles sans amsthm
\providecommand{\qed}{\hfill$\square$}
% \barre : double barre des valeurs interdites dans un tableau de variations (tkz-tab)
\providecommand{\barre}{\ensuremath{\|}}
% environnement proof (amsthm non chargé) utilisé par les modèles
\ifdefined\proof\else\newenvironment{proof}[1][Démonstration]{\par\noindent\textit{#1.}\ }{\qed\par}\fi
\usepackage{lastpage}
\usepackage{hyperref}
\hypersetup{hidelinks}

% ============================================================
% Palette « Pop » OnBuch+ — mêmes hex que la classe P (plus_ui.dart)
% ============================================================
\definecolor{poporange}{HTML}{F59321}
\definecolor{popdark}{HTML}{E07A0C}
\definecolor{popcream}{HTML}{FFF6E8}
\definecolor{popink}{HTML}{1C1714}
\definecolor{poppurple}{HTML}{7A5AE0}
\definecolor{popgreen}{HTML}{2E9E6A}
\definecolor{poppink}{HTML}{E0457B}
\definecolor{popblue}{HTML}{2F7DE1}
\definecolor{popgold}{HTML}{C98A12}
\definecolor{popmuted}{HTML}{8A7F76}
\definecolor{popline}{HTML}{EADFD2}
\definecolor{popgray}{HTML}{8A7F76}
\colorlet{popgrey}{popgray}
% Alias tolérants (le modèle nomme parfois les couleurs différemment)
\colorlet{popwhite}{white}
\colorlet{popblack}{popink}
\colorlet{popred}{poppink}
\colorlet{popyellow}{popgold}
% Teintes claires (fonds de tableaux/encadrés) — le modèle nomme parfois
% les couleurs avec un suffixe L (ex. popblueL) sans les avoir définies.
\colorlet{poporangeL}{poporange!20}
\colorlet{popdarkL}{popdark!20}
\colorlet{poppurpleL}{poppurple!20}
\colorlet{popgreenL}{popgreen!20}
\colorlet{poppinkL}{poppink!20}
\colorlet{popblueL}{popblue!20}
\colorlet{popgoldL}{popgold!20}
\colorlet{popredL}{popred!20}
\colorlet{popgrayL}{popgray!20}
% Teintes claires pour fonds de tableaux / zones
\colorlet{popblueL}{popblue!10!white}
\colorlet{poporangeL}{poporange!14!white}
\colorlet{popgreenL}{popgreen!12!white}
\colorlet{poppurpleL}{poppurple!12!white}
\colorlet{poppinkL}{poppink!12!white}
\colorlet{popgoldL}{popgold!14!white}

\pagecolor{popcream}

% ============================================================
% Typographie
% ============================================================
\defaultfontfeatures{Ligatures=TeX}
\setmainfont{PlusJakartaSans-Medium.ttf}[Path=../../../fonts/,BoldFont=PlusJakartaSans-Bold.ttf]
% Symboles, API et emojis absents de la police principale : repli sur DejaVu Sans (caractères actifs)
\newfontface\symfont{DejaVuSans.ttf}[Path=../../../fonts/]
\makeatletter
\newcommand\symactive[1]{\catcode#1=13 \begingroup\lccode`\~=#1 \lowercase{\endgroup\def~}{{\symfont\char#1}}}
\newcommand\symblank[1]{\catcode#1=13 \begingroup\lccode`\~=#1 \lowercase{\endgroup\def~}{}}
\newcommand\symhyph[1]{\catcode#1=13 \begingroup\lccode`\~=#1 \lowercase{\endgroup\def~}{\mbox{-}}}
\newcount\symc
\newcommand\symrange[2]{\symc=#1 \@whilenum\symc<#2 \do{\expandafter\symactive\expandafter{\the\symc}\advance\symc by 1 }}
\newcommand\symblankrange[2]{\symc=#1 \@whilenum\symc<#2 \do{\expandafter\symblank\expandafter{\the\symc}\advance\symc by 1 }}
\makeatother
\symrange{"0250}{"02FF}\symactive{"2713}\symactive{"2714}\symactive{"2717}\symactive{"2718}\symactive{"2610}\symactive{"2611}\symactive{"2612}
\symhyph{"2011}\symhyph{"2E1A}\symblankrange{"1F300}{"1FAFF}\symblank{"FE0F}
% Maths en sans-serif (Fira Math) avec chiffres et lettres droites de Plus
% Jakarta, pour que les formules se fondent dans le texte.
\usepackage[math-style=ISO,bold-style=ISO]{unicode-math}
\setmathfont{FiraMath-Regular.otf}[Path=../../../fonts/]
\setmathfont{PlusJakartaSans-Medium.ttf}[Path=../../../fonts/,range={up/num,up/{Latin,latin},\mathup}]
\usepackage{icomma} % virgule décimale sans espace en mode math
% Caractères mathématiques Unicode absents des polices de texte (Plus Jakarta,
% Archivo Black) : rendus par leur équivalent mathématique, en texte comme en
% formule — sinon ils s'affichent en carré vide (ex. « Arithmétique dans ℤ »).
\usepackage{newunicodechar}
\newunicodechar{ℝ}{\ensuremath{\mathbb{R}}}
\newunicodechar{ℤ}{\ensuremath{\mathbb{Z}}}
\newunicodechar{ℂ}{\ensuremath{\mathbb{C}}}
\newunicodechar{ℕ}{\ensuremath{\mathbb{N}}}
\newunicodechar{ℚ}{\ensuremath{\mathbb{Q}}}
\newunicodechar{𝔻}{\ensuremath{\mathbb{D}}}
\newunicodechar{α}{\ensuremath{\alpha}}
\newunicodechar{β}{\ensuremath{\beta}}
\newunicodechar{γ}{\ensuremath{\gamma}}
\newunicodechar{δ}{\ensuremath{\delta}}
\newunicodechar{ε}{\ensuremath{\varepsilon}}
\newunicodechar{θ}{\ensuremath{\theta}}
\newunicodechar{λ}{\ensuremath{\lambda}}
\newunicodechar{μ}{\ensuremath{\mu}}
\newunicodechar{σ}{\ensuremath{\sigma}}
\newunicodechar{τ}{\ensuremath{\tau}}
\newunicodechar{φ}{\ensuremath{\varphi}}
\newunicodechar{ψ}{\ensuremath{\psi}}
\newunicodechar{ω}{\ensuremath{\omega}}
\newunicodechar{Σ}{\ensuremath{\Sigma}}
% majuscules produites par \MakeUppercase dans les titres (α-aminés -> Α)
\newunicodechar{Α}{\ensuremath{\alpha}}
\newunicodechar{Β}{\ensuremath{\beta}}
\newunicodechar{Γ}{\ensuremath{\Gamma}}
\newunicodechar{Δ}{\ensuremath{\Delta}}
\newunicodechar{Ε}{\ensuremath{\varepsilon}}
\newunicodechar{Θ}{\ensuremath{\Theta}}
\newunicodechar{Λ}{\ensuremath{\Lambda}}
\newunicodechar{Μ}{\ensuremath{\mu}}
\newunicodechar{Τ}{\ensuremath{\tau}}
\newunicodechar{Φ}{\ensuremath{\Phi}}
\newunicodechar{Ψ}{\ensuremath{\Psi}}
\newunicodechar{Ω}{\ensuremath{\Omega}}
\newunicodechar{↦}{\ensuremath{\mapsto}}
\newunicodechar{∈}{\ensuremath{\in}}
\newunicodechar{∉}{\ensuremath{\notin}}
\newunicodechar{∧}{\ensuremath{\wedge}}
\newunicodechar{⇔}{\ensuremath{\Leftrightarrow}}
\newunicodechar{⟺}{\ensuremath{\Longleftrightarrow}}
\newunicodechar{⇒}{\ensuremath{\Rightarrow}}
\newunicodechar{⟹}{\ensuremath{\Longrightarrow}}
\newunicodechar{∘}{\ensuremath{\circ}}
\newunicodechar{∪}{\ensuremath{\cup}}
\newunicodechar{∩}{\ensuremath{\cap}}
\newunicodechar{≡}{\ensuremath{\equiv}}
\newunicodechar{⊂}{\ensuremath{\subset}}
\newunicodechar{⊥}{\ensuremath{\perp}}
\newunicodechar{∃}{\ensuremath{\exists}}
\newunicodechar{∀}{\ensuremath{\forall}}
\newunicodechar{∛}{\ensuremath{\sqrt[3]{\,}}}
\newunicodechar{∜}{\ensuremath{\sqrt[4]{\,}}}
\newunicodechar{⃗}{\ensuremath{{}^{\to}}}
\newunicodechar{ⁿ}{\ifmmode^{n}\else\textsuperscript{\ensuremath{n}}\fi}
\newunicodechar{ˣ}{\ifmmode^{x}\else\textsuperscript{\ensuremath{x}}\fi}
\newunicodechar{ᵇ}{\ifmmode^{b}\else\textsuperscript{\ensuremath{b}}\fi}
\newunicodechar{ʳ}{\ifmmode^{r}\else\textsuperscript{\ensuremath{r}}\fi}
\newunicodechar{ᵘ}{\ifmmode^{u}\else\textsuperscript{\ensuremath{u}}\fi}
\newunicodechar{ʸ}{\ifmmode^{y}\else\textsuperscript{\ensuremath{y}}\fi}
\newunicodechar{ᵃ}{\ifmmode^{a}\else\textsuperscript{\ensuremath{a}}\fi}
\newunicodechar{ᵉ}{\ifmmode^{e}\else\textsuperscript{\ensuremath{e}}\fi}
\newunicodechar{ᵏ}{\ifmmode^{k}\else\textsuperscript{\ensuremath{k}}\fi}
\newunicodechar{ᵖ}{\ifmmode^{p}\else\textsuperscript{\ensuremath{p}}\fi}
\newunicodechar{ᴺ}{\ifmmode^{N}\else\textsuperscript{\ensuremath{N}}\fi}
\newunicodechar{ᶜ}{\ifmmode^{c}\else\textsuperscript{\ensuremath{c}}\fi}
\newunicodechar{ʰ}{\ifmmode^{h}\else\textsuperscript{\ensuremath{h}}\fi}
\newunicodechar{ᵠ}{\ifmmode^{q}\else\textsuperscript{\ensuremath{q}}\fi}
\newunicodechar{ₙ}{\ifmmode_{n}\else\textsubscript{\ensuremath{n}}\fi}
\newunicodechar{ₐ}{\ifmmode_{a}\else\textsubscript{\ensuremath{a}}\fi}
\newunicodechar{ᵢ}{\ifmmode_{i}\else\textsubscript{\ensuremath{i}}\fi}
\newunicodechar{ᵣ}{\ifmmode_{r}\else\textsubscript{\ensuremath{r}}\fi}
\newunicodechar{ₖ}{\ifmmode_{k}\else\textsubscript{\ensuremath{k}}\fi}
\newunicodechar{ᵦ}{\ifmmode_{\beta}\else\textsubscript{\ensuremath{\beta}}\fi}
\newunicodechar{ₓ}{\ifmmode_{x}\else\textsubscript{\ensuremath{x}}\fi}
\newfontfamily\popdisplay{ArchivoBlack-Regular.ttf}[Path=../../../fonts/]
\newfontfamily\poplogo{SpaceGrotesk-Bold.ttf}[Path=../../../fonts/]
\newfontfamily\popsemi{PlusJakartaSans-SemiBold.ttf}[Path=../../../fonts/]
\newfontfamily\popxbold{PlusJakartaSans-ExtraBold.ttf}[Path=../../../fonts/]
\newfontfamily\popxboldls{PlusJakartaSans-ExtraBold.ttf}[Path=../../../fonts/,LetterSpace=3.0]

\setlist[enumerate]{leftmargin=1.7em,itemsep=5pt,topsep=4pt}
\setlist[itemize]{leftmargin=1.7em,itemsep=4pt,topsep=4pt,label={\color{poporange}\textbullet}}
\setlength{\parindent}{0pt}
\setlength{\parskip}{5pt}
\renewcommand{\arraystretch}{1.25}

% pgfplots : style Pop par défaut
\pgfplotsset{
  every axis/.append style={axis line style={popink,line width=1pt},tick style={popink},
    grid style={popline,line width=0.5pt},label style={font=\small},tick label style={font=\footnotesize}},
  popcurve/.style={poporange,line width=1.8pt,no markers,smooth},
}
% Même style disponible dans un \draw TikZ simple (hors pgfplots)
\tikzset{popcurve/.style={poporange,line width=1.8pt}}
\tikzset{popgrid/.style={popline,very thin}}
\colorlet{popcurve}{poporange} % le nom du style est parfois utilisé comme couleur

% ============================================================
% Pill / squiggle
% ============================================================
\newcommand{\poppill}[3]{%
  \tikz[baseline=-0.55ex]\node[draw=popink,line width=1.2pt,rounded corners=2.6pt,%
    fill=#2,inner xsep=6.5pt,inner ysep=3pt]{%
    \popxboldls\fontsize{7}{7}\selectfont\color{#3}\MakeUppercase{#1}};%
}
\newcommand{\popsquiggle}[1]{%
  \tikz[baseline=-0.4ex]{
    \draw[#1,line width=2.6pt,line cap=round]
      plot [smooth] coordinates {(0,0.09) (0.34,0.01) (0.68,0.21) (1.02,0.01) (1.36,0.10)};
  }%
}
% Mot-clé surligné (marqueur orange sous le texte)
\newcommand{\cle}[1]{{\popxbold\color{popdark}#1}}

% ============================================================
% En-tête / pied
% ============================================================
\pagestyle{fancy}
\fancyhf{}
\renewcommand{\headrulewidth}{0pt}
\renewcommand{\footrulewidth}{0pt}
\lhead{\raisebox{-0.15ex}{\poplogo\fontsize{13}{13}\selectfont\color{popink}OnBuch\color{poporange}+}}
\rhead{\poppill{\DOCMATIERE\ \textbullet\ \DOCNIVEAU\ \textbullet\ Leçon \DOCLECON}{white}{popink}}
\lfoot{\popxbold\fontsize{7.6}{7.6}\selectfont\color{popmuted}\MakeUppercase{Cours signé OnBuch+}\ \textbullet\ {\popsemi\DOCTITRE}}
\rfoot{\tikz[baseline=-0.6ex]\node[rounded corners=8.5pt,fill=popink,minimum height=17pt,minimum width=17pt,inner xsep=5pt,inner sep=0]{\popxbold\fontsize{8}{8}\selectfont\color{popcream}\thepage\,/\,\pageref*{LastPage}};}

% ============================================================
% Titres
% ============================================================
\newcommand{\chaptitre}[2]{%
  \begin{tcolorbox}[enhanced,colback=poporange,colframe=popink,boxrule=2.6pt,arc=11pt,
      drop shadow=popink,left=15pt,right=15pt,top=12pt,bottom=11pt,before skip=3pt,after skip=18pt]
    \poppill{#2}{popink}{popcream}\par\vspace{9pt}
    {\raggedright\hyphenpenalty=10000\exhyphenpenalty=10000\popdisplay\fontsize{22}{24}\selectfont\color{popcream}\MakeUppercase{#1}\par}
  \end{tcolorbox}
}
% Section numérotée : pastille orange avec le numéro + titre Archivo
\newcounter{popsec}
\newcommand{\coursec}[1]{%
  \stepcounter{popsec}\setcounter{popsub}{0}%
  \par\vspace{16pt}\noindent
  \begin{minipage}{\linewidth}
  \tikz[baseline=-0.7ex]\node[draw=popink,line width=1.6pt,fill=poporange,rounded corners=4pt,minimum size=22pt,inner sep=0]{\popdisplay\fontsize{12}{12}\selectfont\color{popcream}\thepopsec};%
  \hspace{8pt}{\popdisplay\fontsize{15}{17}\selectfont\color{popink}\MakeUppercase{#1}}\par
  \vspace{4pt}\noindent{\color{poporange}\rule{36pt}{3.6pt}}
  \end{minipage}\par\nopagebreak\vspace{8pt}
}
\newcounter{popsub}
\newcommand{\courssub}[1]{%
  \stepcounter{popsub}%
  \par\vspace{9pt}\noindent{\popxbold\fontsize{11.5}{13}\selectfont\color{popink}{\color{poporange}\thepopsec.\thepopsub}\hspace{6pt}#1}\par\nopagebreak\vspace{4pt}
}

% ============================================================
% Boîtes pédagogiques — même recette : fond blanc, bordure épaisse colorée,
% ombre dure encre, badge plein. Titre optionnel [#1].
% ============================================================
\tcbset{popbase/.style={breakable,enhanced,colback=white,boxrule=2.3pt,arc=9pt,
  shadow={3pt}{-3pt}{0pt}{popink},left=14pt,right=14pt,top=11pt,bottom=11pt,before skip=11pt,after skip=15pt,
  fontupper=\popsemi\fontsize{10.6}{14.5}\selectfont\color{popink},
  fontlower=\popsemi\fontsize{10.6}{14.5}\selectfont\color{popink}}}
% Coche dessinée (le glyphe ✓ n'existe pas dans les polices du document)
\newcommand{\popcheck}[1]{\tikz[baseline=-0.5ex]\draw[#1,line width=1.6pt,line cap=round,line join=round](-0.13,0.0)--(-0.04,-0.09)--(0.13,0.1);}
\providecommand{\coche}{\popcheck{popgreen}}
\providecommand{\resultat}[1]{\par\smallskip\noindent\fcolorbox{popgreen}{popgreenL}{\parbox{\dimexpr\linewidth-2\fboxsep-2\fboxrule\relax}{\textbf{Résultat.} #1}}\par\smallskip}
\newcommand{\popboxhead}[3]{\poppill{#1}{#2}{white}\ifx\relax#3\relax\else\hspace{6pt}{\popxbold\fontsize{10.6}{12}\selectfont\color{popink}#3}\fi\par\vspace{7pt}}

\newtcolorbox{aretenir}[1][]{popbase,colframe=popgreen,before upper={\popboxhead{À retenir}{popgreen}{#1}}}
% Texte en anglais (document de lecture, dialogue, extrait) : cadre
% d'étude. Un titre optionnel (source, auteur) s'affiche dans l'en-tête.
\newtcolorbox{textebox}[1][]{popbase,colframe=popblue,before upper={\popboxhead{Text}{popblue}{#1}\begin{otherlanguage}{english}},after upper={\end{otherlanguage}}}
% Marques ✓ / ✗ (forme correcte / fautive) souvent utilisées par les modèles
\providecommand{\cross}{\textcolor{poppink}{\ensuremath{\times}}}
\providecommand{\xmark}{\cross}\providecommand{\crossmark}{\cross}\providecommand{\cmark}{\ensuremath{\checkmark}}
% Ligne de réponse / justification à compléter (inventée par les modèles)
\providecommand{\justification}[1]{\par\noindent\textit{Justification~:} \dotfill\par}
% \textipa (package tipa non chargé) : la police ne couvre pas l'API -> simple passage
\providecommand{\textipa}[1]{#1}
% Soulignement (ulem non chargé) : \uline, \ul
\providecommand{\uline}[1]{\underline{#1}}\providecommand{\ul}[1]{\underline{#1}}
% \glossaire{\item ...} inventée par les modèles : liste de glossaire
\providecommand{\glossaire}[1]{\begin{itemize}#1\end{itemize}}
% \alert (beamer) inventée par les modèles : texte mis en évidence
\providecommand{\alert}[1]{\textbf{\textcolor{poppink}{#1}}}
% variantes ASCII d'ordinaux écrites en macro
\providecommand{\iere}{\textsuperscript{re}}\providecommand{\ieme}{\textsuperscript{e}}\providecommand{\ere}{\textsuperscript{re}}\providecommand{\eme}{\textsuperscript{e}}\providecommand{\er}{\textsuperscript{er}}\providecommand{\ie}{\textsuperscript{e}}
% Ordinaux français écrits en macro par les modèles : 1\ière, 3\ième
\newcommand{\ière}{\textsuperscript{re}}\newcommand{\ième}{\textsuperscript{e}}
% \exemple (inventée par les modèles) : étiquette « Exemple : »
\providecommand{\exemple}{\textbf{Exemple~:}\ }
% \square utilisable aussi hors mode math (cases à cocher)
\AtBeginDocument{\let\squareMath\square\renewcommand{\square}{\ensuremath{\squareMath}}}
% Mot ou phrase en anglais à l'intérieur d'un paragraphe français
\newcommand{\eng}[1]{\ifmmode\text{\textit{#1}}\else\begingroup\selectlanguage{english}\itshape #1\endgroup\fi}
\let\esp\eng % alias toléré
% flèches d'intonation inventées par les modèles
\providecommand{\textsearrow}{}\renewcommand{\textsearrow}{\ensuremath{\searrow}}\providecommand{\textnearrow}{}\renewcommand{\textnearrow}{\ensuremath{\nearrow}}
% \fr{..} : mot français (inventé par les modèles)
\providecommand{\fr}[1]{\textit{#1}}
\newtcolorbox{exemplebox}[1][]{popbase,colframe=poporange,before upper={\popboxhead{Exemple}{poporange}{#1}}}
\newtcolorbox{definition}[1][]{popbase,colframe=popblue,before upper={\popboxhead{Définition}{popblue}{#1}}}
\newtcolorbox{propriete}[1][]{popbase,colframe=poppurple,before upper={\popboxhead{Propriété}{poppurple}{#1}}}
\newtcolorbox{methode}[1][]{popbase,colframe=popgold,before upper={\popboxhead{Méthode}{popgold}{#1}}}
\newtcolorbox{attention}[1][]{popbase,colframe=poppink,before upper={\popboxhead{Attention}{poppink}{#1}}}
\newtcolorbox{experience}[1][]{popbase,colframe=popblue,colback=popblueL,before upper={\popboxhead{Expérience}{popblue}{#1}}}
% « Le savais-tu ? » : carte violette pleine (contexte, culture, Cameroun)
\newtcolorbox{savaistu}[1][]{popbase,colback=poppurple,colframe=popink,
  fontupper=\popsemi\fontsize{10.4}{14.2}\selectfont\color{white},
  before upper={\poppill{Le savais-tu\,?}{white}{poppurple}\ifx\relax#1\relax\else\hspace{6pt}{\popxbold\color{white}#1}\fi\par\vspace{7pt}}}
% Exercice résolu : énoncé en haut, \tcblower puis la solution sur fond crème
\newcounter{popexo}\newcounter{popexr}
\newtcolorbox{exoresolu}[1][]{popbase,colframe=poporange,colbacklower=popcream,
  segmentation style={popink,line width=1.2pt,dashed},
  before upper={\stepcounter{popexr}\popboxhead{Exercice résolu \thepopexr}{poporange}{#1}},
  before lower={\poppill{Solution}{popink}{popcream}\par\vspace{6pt}}}
% Exercice (non résolu en place) — la correction est dans la partie Corrigés
\newtcolorbox{exercice}[1][]{popbase,colframe=popink,
  before upper={\stepcounter{popexo}\popboxhead{Exercice \thepopexo}{popink}{#1}}}
% Corrigé
\newtcolorbox{corrige}[1][]{popbase,colframe=popgreen,colback=popgreenL,
  before upper={\popboxhead{Corrigé}{popgreen}{#1}}}
% Objectifs / prérequis / bilan
\newtcolorbox{objectifs}{popbase,colframe=popink,colback=poporangeL,
  before upper={\popboxhead{Objectifs de la leçon}{popink}{}}}
\newtcolorbox{prerequis}{popbase,colframe=popink,colback=popblueL,
  before upper={\popboxhead{Prérequis}{popblue}{}}}
\newtcolorbox{bilan}{popbase,colframe=popink,colback=white,boxrule=2.8pt,
  before upper={\popboxhead{Fiche bilan}{poporange}{L'essentiel à connaître pour l'examen}}}
% Figure encadrée avec légende
\newtcolorbox{popfigure}[1][]{enhanced,breakable=false,colback=white,colframe=popline,boxrule=1.4pt,arc=8pt,
  left=10pt,right=10pt,top=10pt,bottom=8pt,before skip=10pt,after skip=12pt,halign=center,
  lower separated=false,halign lower=center,
  fontlower=\popsemi\fontsize{9}{11.5}\selectfont\color{popmuted},
  #1}
\newcommand{\legende}[1]{\par\vspace{6pt}{\popsemi\fontsize{9}{11.5}\selectfont\color{popmuted}#1}}

% ============================================================
% Signature OnBuch+
% ============================================================
\newcommand{\poplogoinline}[1]{{\poplogo\fontsize{#1}{#1}\selectfont\color{popink}OnBuch\color{poporange}+}}
\newcommand{\signatureonbuch}{%
  \begin{tcolorbox}[enhanced,colback=popink,colframe=popink,boxrule=2.2pt,arc=10pt,
      drop shadow=poporange,left=14pt,right=14pt,top=10pt,bottom=10pt,before skip=0pt,after skip=0pt]
    \tikz[baseline=-0.7ex]\node[circle,fill=popgreen,draw=popcream,line width=1.3pt,minimum size=22pt,inner sep=0]{\popcheck{white}};%
    \hspace{9pt}\begin{minipage}[c]{0.62\linewidth}
      {\popxbold\fontsize{12}{13}\selectfont\color{popcream}Signé\ \textbullet\ {\poplogo OnBuch\color{poporange}+}}\par\vspace{2pt}
      {\raggedright\popsemi\fontsize{8}{10.5}\selectfont\color{popline}\MakeUppercase{Équipe pédagogique OnBuch+}\\\MakeUppercase{Conforme au programme officiel MINESEC}\par}
    \end{minipage}\hfill
    {\poplogo\fontsize{22}{22}\selectfont\color{popcream}OnBuch\color{poporange}+}
  \end{tcolorbox}%
}

% ============================================================
% Page de garde
% #1 titre · #2 matière · #3 niveau · #4 module · #5 n° leçon · #6 durée ·
% #7 description / objectifs (texte ou itemize)
% ============================================================
\newcommand{\pagegarde}[7]{%
  \thispagestyle{empty}
  \newgeometry{a4paper,margin=1.7cm,top=1.6cm,bottom=1.4cm}%
  \begin{tikzpicture}[remember picture,overlay]
    \fill[poporange!18] ([xshift=-3.2cm,yshift=-3.6cm]current page.north east) circle (4.6cm);
    \fill[poppurple!14] ([xshift=2.4cm,yshift=7.6cm]current page.south west) circle (3.8cm);
  \end{tikzpicture}%
  {\poplogo\fontsize{30}{30}\selectfont\color{popink}OnBuch\color{poporange}+}\hfill
  \raisebox{0.6ex}{\poppill{Cours signé \textbullet\ Premium}{popink}{popcream}}\par
  \vspace{1pt}\popsquiggle{poppurple}\par
  \vspace{8pt}
  \begin{tcolorbox}[enhanced,colback=poporange,colframe=popink,boxrule=2.8pt,arc=13pt,
      drop shadow=popink,left=18pt,right=18pt,top=12pt,bottom=13pt,after skip=10pt]
    \poppill{#2 \textbullet\ #3}{popink}{popcream}\hspace{5pt}\poppill{#4}{white}{popink}\par\vspace{8pt}
    {\popdisplay\fontsize{14}{14}\selectfont\color{popink}LEÇON #5}\par\vspace{4pt}
    {\raggedright\hyphenpenalty=10000\exhyphenpenalty=10000\popdisplay\fontsize{25}{27}\selectfont\color{popcream}\MakeUppercase{#1}\par}
    \vspace{8pt}
    {\popxbold\fontsize{9.5}{9.5}\selectfont\color{popink}\MakeUppercase{Durée conseillée : #6}}
  \end{tcolorbox}
  \begin{tcolorbox}[enhanced,colback=white,colframe=popink,boxrule=2.2pt,arc=10pt,
      drop shadow=popink,left=16pt,right=16pt,top=10pt,bottom=10pt,after skip=8pt]
    \poppill{Dans cette leçon}{poppurple}{white}\par\vspace{5pt}
    \popsemi\fontsize{9.3}{12.6}\selectfont\color{popink}#7
  \end{tcolorbox}
  \noindent\begin{minipage}[t]{0.487\linewidth}
    \begin{tcolorbox}[enhanced,colback=popgreen,colframe=popink,boxrule=2.2pt,arc=10pt,
        drop shadow=popink,left=11pt,right=11pt,top=10pt,bottom=10pt,height=3.1cm,valign=center]
      \tcbox[colback=white,colframe=popink,boxrule=1.3pt,arc=5pt,boxsep=0pt,nobeforeafter,
        left=3pt,right=3pt,top=3pt,bottom=3pt,valign=center]{\qrcode[height=1.9cm]{https://play.google.com/store/apps/details?id=com.luvvix.onbuch&hl=fr}}%
      \hspace{8pt}\begin{minipage}[c]{0.5\linewidth}
        {\popdisplay\fontsize{11}{12.5}\selectfont\color{white}\MakeUppercase{Télécharge OnBuch}}\par\vspace{4pt}
        {\raggedright\popsemi\fontsize{8.4}{11}\selectfont\color{white}Gratuit \textbullet\ Google Play\\Tuteur IA, annales, concours, résultats\par}
      \end{minipage}
    \end{tcolorbox}
  \end{minipage}\hfill
  \begin{minipage}[t]{0.487\linewidth}
    \begin{tcolorbox}[enhanced,colback=poppurple,colframe=popink,boxrule=2.2pt,arc=10pt,
        drop shadow=popink,left=11pt,right=11pt,top=10pt,bottom=10pt,height=3.1cm,valign=center]
      \tikz[baseline=-0.7ex]\node[circle,fill=white,draw=popink,line width=1.3pt,minimum size=1.6cm,inner sep=0]{\popdisplay\fontsize{24}{24}\selectfont\color{poppurple}+};%
      \hspace{8pt}\begin{minipage}[c]{0.6\linewidth}
        {\popdisplay\fontsize{11}{12.5}\selectfont\color{white}\MakeUppercase{Passe à OnBuch+}}\par\vspace{4pt}
        {\raggedright\popsemi\fontsize{8.4}{11}\selectfont\color{white}Tous les cours signés, Tuteur IA illimité, fiches PDF — onglet OnBuch+ de l'app\par}
      \end{minipage}
    \end{tcolorbox}
  \end{minipage}
  \vspace{8pt}
  \popcontenu
  \vfill
  \signatureonbuch
  \restoregeometry
  \newpage
}

% Rangée « Ce cours contient » (page de garde)
\newcommand{\poptile}[3]{% #1 couleur #2 gros texte #3 légende
  \begin{tcolorbox}[enhanced,colback=#1,colframe=popink,boxrule=2pt,arc=9pt,shadow={2.5pt}{-2.5pt}{0pt}{popink},
      left=6pt,right=6pt,top=9pt,bottom=9pt,halign=center,valign=center,height=1.75cm,width=\linewidth,nobeforeafter]
    {\popdisplay\fontsize{12.5}{13}\selectfont\color{white}\MakeUppercase{#2}}\par\vspace{3pt}
    {\centering\popsemi\fontsize{7.8}{9.5}\selectfont\color{white}#3\par}
  \end{tcolorbox}}
\newcommand{\popcontenu}{%
  \par\noindent{\popxboldls\fontsize{7.5}{7.5}\selectfont\color{popmuted}CE COURS SIGNÉ CONTIENT}\par\vspace{7pt}
  \noindent\begin{minipage}[t]{0.235\linewidth}\poptile{popblue}{Cours}{complet, illustré, pas à pas}\end{minipage}\hfill
  \begin{minipage}[t]{0.235\linewidth}\poptile{poporange}{Méthodes}{et exercices résolus}\end{minipage}\hfill
  \begin{minipage}[t]{0.235\linewidth}\poptile{poppink}{Exercices}{type examen + corrigés}\end{minipage}\hfill
  \begin{minipage}[t]{0.235\linewidth}\poptile{popgreen}{Fiche bilan}{l'essentiel à retenir}\end{minipage}\par}

% Sommaire
\newtcolorbox{sommaire}{enhanced,colback=white,colframe=popink,boxrule=2.2pt,arc=10pt,
  drop shadow=popink,left=18pt,right=18pt,top=13pt,bottom=13pt,before skip=6pt,after skip=20pt,
  before upper={\poppill{Sommaire}{poppurple}{white}\par\vspace{9pt}\popsemi\fontsize{10.6}{14.5}\selectfont\color{popink}}}

% Signature de fin de cours — poussée tout en bas de la dernière page
\newcommand{\findecours}{%
  \par\vfill\nopagebreak
  \begin{center}\popsquiggle{poporange}\end{center}\vspace{4pt}
  \signatureonbuch
}
\AtBeginDocument{\def\llbracket{\mathopen{[\mkern-3.5mu[}}\def\rrbracket{\mathclose{]\mkern-3.5mu]}}} % intervalles d'entiers (glyphe ⟦ absent de la police)
% Glyphes absents de Fira Math : redéfinis à partir de glyphes disponibles
\AtBeginDocument{%
  \def\setminus{\mathbin{\backslash}}\let\smallsetminus\setminus
  \def\vdots{\mathord{\vbox{\baselineskip3.2pt\lineskiplimit0pt\kern4pt\hbox{.}\hbox{.}\hbox{.}}}}%
  \def\ddots{\mathinner{\mkern1mu\raise6.5pt\hbox{.}\mkern2mu\raise3.6pt\hbox{.}\mkern2mu\raise0.7pt\hbox{.}\mkern1mu}}%
  \def\triangle{\mathord{\scalebox{0.9}{$\symup{\Delta}$}}}%
  \def\nabla{\mathord{\raisebox{\depth}{\scalebox{1}[-1]{$\symup{\Delta}$}}}}%
  \def\checkmark{\popcheck{popdark}}%
  \def\star{\mathbin{\ast}}\def\bigstar{\mathord{\ast}}%
  \def\diamond{\mathbin{\scalebox{0.6}{$\Diamond$}}}%
  \def\bigcup{\mathop{\mathchoice{\raisebox{-0.3em}{\scalebox{1.7}{$\cup$}}}{\scalebox{1.25}{$\cup$}}{\cup}{\cup}}\displaylimits}%
  \def\bigcap{\mathop{\mathchoice{\raisebox{-0.3em}{\scalebox{1.7}{$\cap$}}}{\scalebox{1.25}{$\cap$}}{\cap}{\cap}}\displaylimits}%
  \def\lesseqqgtr{\mathrel{\lessgtr}}\def\bigoplus{\mathop{\oplus}\displaylimits}%
}
\providecommand{\ssi}{si et seulement si} % abréviation fréquente
\AtBeginDocument{\providecommand{\wideparen}{\overparen}} % arc de cercle

%% FIGFIX-BEGIN v4 — correctifs globaux des figures (docs/onbuch-generation/tools/figfix)
\usepackage{adjustbox}\usepackage{etoolbox}
% toute figure TikZ/pgfplots plus large que la ligne est réduite à la largeur disponible
\BeforeBeginEnvironment{tikzpicture}{\begin{adjustbox}{max width=\linewidth}}
\AfterEndEnvironment{tikzpicture}{\end{adjustbox}}
% légendes pgfplots lisibles par-dessus les courbes
\pgfplotsset{every axis/.append style={legend style={fill=white,fill opacity=0.92,text opacity=1,draw=black!15,inner sep=3pt}}}
% étiquettes d'arêtes (syntaxe edge["..."]) avec fond blanc
\tikzset{every edge quotes/.append style={fill=white,inner sep=1.2pt}}
%% FIGFIX-END
\begin{document}

\pagegarde{Lesson 10 — Grammar: Conditionals (1, 2, 3)}{Anglais}{1ère A}{Chapter 1 : Family and Social Life}{ENA10}{2 h}{Cette leçon de grammaire présente les trois conditionnels de l'anglais (first, second, third conditional) : leur formation, leurs emplois et leurs différences de sens. À travers l'observation d'exemples contextualisés, des tableaux de formes et une comparaison systématique avec le français, les élèves apprennent à exprimer l'hypothèse réelle, irréelle et le regret, en évitant les erreurs typiques des francophones.
\vspace{4pt}
\begin{multicols}{2}\small
\begin{itemize}
\item À la fin de cette leçon, je sais identifier les trois types de conditionnels dans un texte ou un énoncé
\item À la fin de cette leçon, je sais construire correctement le first conditional (If + present simple, will + base verbale)
\item À la fin de cette leçon, je sais construire correctement le second conditional (If + past simple, would + base verbale)
\item À la fin de cette leçon, je sais construire correctement le third conditional (If + past perfect, would have + participe passé)
\item À la fin de cette leçon, je sais choisir le bon conditionnel selon le degré de réalité de l'hypothèse
\item À la fin de cette leçon, je sais utiliser « were » avec toutes les personnes dans le second conditional formel (If I were you...)
\end{itemize}\end{multicols}}

\begin{sommaire}
\begin{multicols}{2}
\begin{enumerate}
\item Observation et découverte
\item First conditional : l'hypothèse réelle
\item Second conditional : l'irréel du présent
\item Third conditional : le regret du passé
\item Synthèse, comparaison avec le français et pièges
\item Entraînement guidé et production
\item Activité d'intégration
\item Méthodes et exercices résolus
\item Exercices
\item Corrigés des exercices
\item Fiche bilan
\end{enumerate}
\end{multicols}
\end{sommaire}

\begin{objectifs}
\begin{itemize}
\item À la fin de cette leçon, je sais identifier les trois types de conditionnels dans un texte ou un énoncé
\item À la fin de cette leçon, je sais construire correctement le first conditional (If + present simple, will + base verbale)
\item À la fin de cette leçon, je sais construire correctement le second conditional (If + past simple, would + base verbale)
\item À la fin de cette leçon, je sais construire correctement le third conditional (If + past perfect, would have + participe passé)
\item À la fin de cette leçon, je sais choisir le bon conditionnel selon le degré de réalité de l'hypothèse
\item À la fin de cette leçon, je sais utiliser « were » avec toutes les personnes dans le second conditional formel (If I were you...)
\item À la fin de cette leçon, je sais éviter les erreurs fréquentes des francophones (will ou would dans la proposition avec if)
\item À la fin de cette leçon, je sais employer les conditionnels dans des phrases et de courts textes sur la vie familiale et sociale
\item À la fin de cette leçon, je sais exprimer un regret ou un reproche sur le passé avec le third conditional
\end{itemize}
\end{objectifs}

\begin{prerequis}
\begin{itemize}
\item Le present simple et le futur avec will (Seconde)
\item Le past simple et les verbes irréguliers (Seconde)
\item Le present perfect et la formation du participe passé (Seconde)
\item La modalité : would comme conditionnel (début de 1ère)
\item La structure de la phrase anglaise : sujet + verbe + complément
\end{itemize}
\end{prerequis}

\newpage

\coursec{Observation et découverte}

\courssub{Situation d'accroche : trois phrases, trois mondes}

Pendant les grandes vacances, à Yaoundé, trois membres de la même famille parlent de l'avenir, de leurs rêves et de leurs regrets. Divine, qui vient de passer son GCE Advanced Level, dit à son ami :

\eng{“If I pass my GCE Advanced Level, my father will buy me a laptop.”}

« Si je réussis mon GCE Advanced Level, mon père m'achètera un ordinateur portable. »

Son ami, lui, rêve tout haut :

\eng{“If I were rich, I would travel to London.”}

« Si j'étais riche, je voyagerais à Londres. »

Et leur grand frère, qui a échoué l'année précédente, soupire :

\eng{“If I had studied harder, I would have succeeded last year.”}

« Si j'avais travaillé plus dur, j'aurais réussi l'année dernière. »

Trois phrases, trois situations complètement différentes : une \cle{possibilité réelle} (Divine peut vraiment réussir), un \cle{rêve irréel} (son ami n'est pas riche), un \cle{regret sur le passé} (le frère n'a pas assez travaillé, et c'est trop tard). Pourtant, les trois phrases commencent par le même petit mot : \eng{if}.

\textbf{Problématique :} comment l'anglais exprime-t-il ces trois types d'hypothèses ? Pourquoi les temps des verbes changent-ils d'une phrase à l'autre ? C'est toute la logique des \cle{conditionnels} que nous allons découvrir aujourd'hui, en observant d'abord, avant d'énoncer les règles.

\courssub{Lecture d'un mini-dialogue : la famille de Manka}

Lisons maintenant un dialogue complet. La scène se passe à Buea, un samedi matin. La pluie tombe, et la famille de Manka voulait aller à la ferme.

\begin{textebox}[A rainy Saturday in Buea — a family dialogue]
It is Saturday morning in Buea, and heavy rain is falling on the town. Manka, her brother Tabi and their mother are sitting in the living room, looking through the window.

“If the rain stops, we will go to the farm,” says Manka. “We need to harvest the plantains before Monday.”

“If I had a car, I would drive you there,” replies Tabi. “But I don't, so we will have to wait for a bendskin or a taxi.”

Their mother sighs. “If we had left earlier, we would have arrived before the storm,” she adds. “I told you yesterday evening that the sky was dark.”

“Don't worry, Mum,” says Manka. “If the weather improves this afternoon, we will still have time. And if Tabi passes his driving test next month, he will buy a second-hand car, won't you, Tabi?”

Tabi laughs. “If I were a rich businessman, I would buy a new car, not a second-hand one! But first, I need to pass the test. If I had practised more last year, I would have passed it already.”

“Well,” says their mother with a smile, “if you help me in the kitchen now, I will teach you how to drive carefully — with words, not with a car!”
\end{textebox}

\begin{definition}[Glossaire du dialogue]
\begin{itemize}
\item \eng{to harvest} (verbe) : récolter
\item \eng{plantains} (nom pluriel) : les plantains
\item \eng{a bendskin} (nom, camerounianisme) : une moto-taxi
\item \eng{to sigh} (verbe) : soupirer
\item \eng{second-hand} (adjectif) : d'occasion
\item \eng{a driving test} (nom) : un examen de conduite (permis)
\item \eng{to practise} (verbe) : s'entraîner
\item \eng{carefully} (adverbe) : prudemment
\end{itemize}
\end{definition}

\textbf{Questions de compréhension (reading) :}
\begin{enumerate}
\item Why can't the family go to the farm immediately?
\item What does Tabi dream of buying?
\item What does the mother regret?
\item What will happen if Tabi passes his driving test?
\end{enumerate}

\courssub{Repérage : trouvons toutes les phrases en « if »}

\begin{methode}[Comment observer une phrase conditionnelle en quatre étapes]
Pour analyser n'importe quelle phrase conditionnelle, suis toujours la même démarche :
\begin{enumerate}
\item \textbf{Souligne} le mot \eng{if} et repère la proposition qu'il introduit (la \cle{if-clause}).
\item \textbf{Identifie le temps du verbe} dans la if-clause : present simple ? past simple ? past perfect ?
\item \textbf{Identifie la forme du verbe} dans l'autre proposition (la \cle{main clause}) : \eng{will} + infinitif ? \eng{would} + infinitif ? \eng{would have} + participe passé ?
\item \textbf{Déduis le type} : hypothèse possible (first), imaginaire (second) ou passée avec regret (third).
\end{enumerate}
\end{methode}

Appliquons cette méthode au dialogue. Voici toutes les phrases en \eng{if} qu'il contient, classées selon les temps utilisés :

\begin{center}
\begin{tabularx}{\linewidth}{|X|X|X|}
\hline
\rowcolor{popblueL} \textbf{Exemple du dialogue} & \textbf{Temps utilisés} & \textbf{Degré de réalité} \\
\hline
\eng{If the rain stops, we will go to the farm.} & present simple + \eng{will} + infinitif & hypothèse possible (la pluie peut s'arrêter) \\
\hline
\eng{If Tabi passes his driving test, he will buy a car.} & present simple + \eng{will} + infinitif & hypothèse possible (il peut réussir) \\
\hline
\eng{If you help me, I will teach you.} & present simple + \eng{will} + infinitif & hypothèse possible (marchandage réel) \\
\hline
\eng{If I had a car, I would drive you there.} & past simple + \eng{would} + infinitif & hypothèse imaginaire (Tabi n'a pas de voiture) \\
\hline
\eng{If I were rich, I would buy a new car.} & past simple + \eng{would} + infinitif & hypothèse imaginaire (il n'est pas riche) \\
\hline
\eng{If we had left earlier, we would have arrived before the storm.} & past perfect + \eng{would have} + participe passé & hypothèse passée, regret (ils ne sont pas partis tôt) \\
\hline
\eng{If I had practised more, I would have passed already.} & past perfect + \eng{would have} + participe passé & hypothèse passée, regret (il ne s'est pas entraîné) \\
\hline
\end{tabularx}
\end{center}

\textbf{Question guide :} pour chaque phrase, demande-toi : l'hypothèse est-elle \cle{possible} (elle peut se réaliser), \cle{imaginaire} (elle est contraire à la réalité présente) ou \cle{passée} (elle est contraire à ce qui s'est réellement passé) ? La réponse détermine le type de conditionnel.

\begin{exoresolu}[Repérage guidé]
Énoncé : dans la phrase \eng{“If I had studied harder, I would have succeeded last year,”} souligne la if-clause, identifie les deux temps verbaux et précise le degré de réalité.
\tcblower
Solution détaillée :
\begin{enumerate}
\item La if-clause est \eng{If I had studied harder} ; la main clause est \eng{I would have succeeded last year}.
\item Temps après \eng{if} : \eng{had studied} = past perfect (auxiliaire \eng{had} + participe passé \eng{studied}).
\item Forme dans la main clause : \eng{would have succeeded} = \eng{would have} + participe passé.
\item Conclusion : l'hypothèse porte sur le passé et elle est contraire à la réalité (il n'a pas étudié assez) : c'est un \cle{third conditional}, qui exprime un regret.
\end{enumerate}
\end{exoresolu}

\courssub{La structure en deux propositions : if-clause et main clause}

Observons la construction de près. Toute phrase conditionnelle contient \textbf{deux propositions} :

\begin{itemize}
\item la \cle{if-clause} (proposition subordonnée) : elle introduit la \textbf{condition}, l'hypothèse ;
\item la \cle{main clause} (proposition principale) : elle exprime le \textbf{résultat}, la conséquence.
\end{itemize}

\begin{exemplebox}[Les deux propositions en action]
\eng{If I pass the exam, I will celebrate.}

« Si je réussis l'examen, je fêterai cela. »

\begin{itemize}
\item if-clause (condition) : \eng{If I pass the exam}
\item main clause (résultat) : \eng{I will celebrate}
\end{itemize}

\eng{If I were you, I would apologise.}

« Si j'étais toi, je m'excuserais. »

\begin{itemize}
\item if-clause (condition imaginaire) : \eng{If I were you}
\item main clause (conseil) : \eng{I would apologise}
\end{itemize}
\end{exemplebox}

\begin{propriete}[L'ordre des propositions et la ponctuation]
Les deux propositions peuvent \textbf{changer de place} sans changer le sens, mais la ponctuation change :
\begin{itemize}
\item Si la \textbf{if-clause est en tête}, on met une \cle{virgule} entre les deux propositions : \eng{If the rain stops, we will go to the farm.}
\item Si la \textbf{main clause est en tête}, il n'y a \cle{pas de virgule} : \eng{We will go to the farm if the rain stops.}
\end{itemize}
Le sens reste identique ; seule la ponctuation change.
\end{propriete}

\begin{attention}[Erreur fréquente des francophones]
En français, on hésite parfois avec la virgule ; en anglais, la règle est stricte : \textbf{virgule uniquement quand la if-clause ouvre la phrase}. N'écris jamais : \eng{I will celebrate, if I pass the exam.} (virgule fautive). Écris : \eng{I will celebrate if I pass the exam.}
\end{attention}

\courssub{Dégagement progressif des trois modèles}

Récapitulons ce que l'observation nous a appris. Chaque type d'hypothèse possède sa propre « recette » de temps verbaux :

\begin{center}
\begin{tabular}{|l|l|l|}
\hline
\rowcolor{popgreenL} \textbf{Type} & \textbf{If-clause} & \textbf{Main clause} \\
\hline
First conditional (possible) & \eng{if} + present simple & \eng{will} + infinitif \\
\hline
Second conditional (imaginaire) & \eng{if} + past simple & \eng{would} + infinitif \\
\hline
Third conditional (passé, regret) & \eng{if} + past perfect & \eng{would have} + participe passé \\
\hline
\end{tabular}
\end{center}

\begin{aretenir}[L'idée directrice]
Plus l'hypothèse s'éloigne de la réalité, plus les temps « reculent » dans le passé : présent pour ce qui est possible, past simple pour ce qui est imaginaire, past perfect pour ce qui est définitivement passé. C'est ce recul des temps que les grammairiens appellent le \cle{recul temporel} (\eng{backshift}).
\end{aretenir}

\begin{popfigure}
\begin{tikzpicture}[
  node distance=0.6cm,
  start/.style={rectangle, rounded corners, draw=popdark, fill=popcream, text width=4.2cm, align=center, font=\small\bfseries},
  quest/.style={diamond, draw=popblue, fill=popblueL, text width=3.4cm, align=center, font=\small, aspect=2.2},
  res/.style={rectangle, rounded corners, draw=popgreen, fill=popgreenL, text width=4.6cm, align=center, font=\small},
  arr/.style={-{Stealth[length=2.5mm]}, thick, popdark}
]
\node[start] (s) {Phrase avec \eng{if} : quelle hypothèse ?};
\node[quest, below=of s] (q1) {L'hypothèse est-elle possible, réelle ?};
\node[res, below left=0.8cm and -1.2cm of q1, align=center] (r1){\textbf{First conditional}\\ \eng{if} + present, \eng{will} + infinitif};
\node[quest, below right=0.8cm and -1.2cm of q1] (q2) {Non : est-elle imaginaire (présent) ?};
\node[res, below left=0.8cm and -1.6cm of q2, align=center] (r2){\textbf{Second conditional}\\ \eng{if} + past simple, \eng{would} + infinitif};
\node[res, below right=0.8cm and -1.6cm of q2, align=center] (r3){\textbf{Third conditional}\\ \eng{if} + past perfect, \eng{would have} + participe passé};
\draw[arr] (s) -- (q1);
\draw[arr] (q1) -- node[left, font=\small]{oui} (r1);
\draw[arr] (q1) -- node[right, font=\small]{non} (q2);
\draw[arr] (q2) -- node[left, font=\small]{oui} (r2);
\draw[arr] (q2) -- node[right, font=\small]{non : passée / regret} (r3);
\end{tikzpicture}
\legende{Organigramme de choix : quel conditionnel pour quelle hypothèse ?}
\end{popfigure}

\begin{popfigure}
\begin{tikzpicture}[font=\small]
\draw[fill=popblueL, draw=popblue] (0,0) rectangle (4.4,1.6);
\draw[fill=poporangeL, draw=poporange] (4.6,0) rectangle (9.4,1.6);
\draw[fill=popgreenL, draw=popgreen] (9.6,0) rectangle (14.2,1.6);
\node[text width=4cm, align=center] at (2.2,0.8) {\textbf{Exemple}\\ \eng{If I pass, I will celebrate.}};
\node[text width=4.4cm, align=center] at (7,0.8) {\textbf{Temps utilisés}\\ present simple + \eng{will}};
\node[text width=4.2cm, align=center] at (11.9,0.8) {\textbf{Degré de réalité}\\ possible, réel};
\draw[fill=popblueL, draw=popblue] (0,-1.9) rectangle (4.4,-0.3);
\draw[fill=poporangeL, draw=poporange] (4.6,-1.9) rectangle (9.4,-0.3);
\draw[fill=popgreenL, draw=popgreen] (9.6,-1.9) rectangle (14.2,-0.3);
\node[text width=4cm, align=center] at (2.2,-1.1) {\eng{If I were you, I would apologise.}};
\node[text width=4.4cm, align=center] at (7,-1.1) {past simple + \eng{would}};
\node[text width=4.2cm, align=center] at (11.9,-1.1) {imaginaire, irréel du présent};
\draw[fill=popblueL, draw=popblue] (0,-3.8) rectangle (4.4,-2.2);
\draw[fill=poporangeL, draw=poporange] (4.6,-3.8) rectangle (9.4,-2.2);
\draw[fill=popgreenL, draw=popgreen] (9.6,-3.8) rectangle (14.2,-2.2);
\node[text width=4cm, align=center] at (2.2,-3) {\eng{If we had left earlier, we would have arrived.}};
\node[text width=4.4cm, align=center] at (7,-3) {past perfect + \eng{would have} + participe};
\node[text width=4.2cm, align=center] at (11.9,-3) {passé, regret, irréversible};
\end{tikzpicture}
\legende{Tableau à trois colonnes : exemple, temps utilisés, degré de réalité.}
\end{popfigure}

\courssub{Comparaison avec le français : la grande différence}

\begin{center}
\begin{tabularx}{\linewidth}{|X|X|X|}
\hline
\rowcolor{popblueL} \textbf{Français} & \textbf{English} & \textbf{Remarque} \\
\hline
Si je réussis, je fêterai. & \eng{If I pass, I will celebrate.} & présent après « si » dans les deux langues \\
\hline
Si j'étais toi, je m'excuserais. & \eng{If I were you, I would apologise.} & imparfait en français = past simple en anglais \\
\hline
Si j'avais étudié, j'aurais réussi. & \eng{If I had studied, I would have succeeded.} & plus-que-parfait = past perfect \\
\hline
Si j'aurais su\ldots{} (fautif mais fréquent à l'oral) & impossible en anglais & jamais \eng{would} après \eng{if} \\
\hline
\end{tabularx}
\end{center}

\begin{savaistu}[Jamais « will » ni « would » juste après « if » !]
En anglais, on ne met \textbf{JAMAIS} \eng{will} ou \eng{would} juste après \eng{if} dans une hypothèse. C'est la différence majeure avec le français oral, où l'on entend parfois « si j'aurais su ». En anglais, \eng{If I would have known} est une faute grave : il faut dire \eng{If I had known}. Le mot \eng{if} « attire » toujours un temps simple : present, past simple ou past perfect. La seule exception : \eng{if} signifiant « si » interrogatif indirect (\eng{I don't know if he will come.} = « Je ne sais pas s'il viendra. ») — mais ce n'est plus une hypothèse !
\end{savaistu}

\begin{exercice}[À toi de jouer : classe ces phrases]
Classe chaque phrase selon le type d'hypothèse (possible, imaginaire, passée) et souligne la if-clause :
\begin{enumerate}
\item \eng{If it rains tomorrow, we will stay at home.}
\item \eng{If I lived in Douala, I would eat fresh fish every day.}
\item \eng{If she had taken the bus, she would have arrived on time.}
\item \eng{If you study hard, you will pass the GCE.}
\end{enumerate}
\end{exercice}

\begin{corrige}[Exercice : classe ces phrases]
\begin{enumerate}
\item \eng{If it rains tomorrow} : present + \eng{will} → hypothèse \textbf{possible} (first conditional).
\item \eng{If I lived in Douala} : past simple + \eng{would} → hypothèse \textbf{imaginaire} (second conditional) ; la personne n'habite pas à Douala.
\item \eng{If she had taken the bus} : past perfect + \eng{would have} → hypothèse \textbf{passée}, regret (third conditional) ; elle n'a pas pris le bus.
\item \eng{If you study hard} : present + \eng{will} → hypothèse \textbf{possible} (first conditional).
\end{enumerate}
\end{corrige}

\begin{aretenir}[Ce qu'il faut retenir de cette observation]
\begin{itemize}
\item Une phrase conditionnelle = \cle{if-clause} (condition) + \cle{main clause} (résultat).
\item Virgule seulement si la if-clause est en tête de phrase.
\item Trois modèles : possible (present + \eng{will}), imaginaire (past simple + \eng{would}), passée/regret (past perfect + \eng{would have}).
\item Jamais \eng{will} ni \eng{would} directement après \eng{if}.
\end{itemize}
Dans la prochaine section, nous étudierons en détail le \cle{first conditional} : l'hypothèse réelle.
\end{aretenir}

\coursec{First conditional : l'hypothèse réelle}

Dans la section précédente, nous avons observé le dialogue entre Mme Ngo et son fils Junior : \eng{If you pass your exams, we will celebrate.} Vous avez remarqué que la phrase se coupe en deux morceaux : une petite proposition introduite par \eng{if} au présent, et une grande proposition avec \eng{will}. Nous allons maintenant étudier cette structure en détail : c'est le \cle{first conditional}, le conditionnel de l'hypothèse \textbf{réelle} ou \textbf{probable}.

\courssub{Observer avant de généraliser}

\begin{textebox}[A family conversation in Yaoundé]
Mrs Ngo is talking to her son Junior about the end-of-year exams. She says: “If my son passes his exams, we will organise a big family party. If the guests come early, they will help us cook. If we have enough money, we will buy a goat at the market. And if it rains on Saturday, we will move the tables inside the house.” Junior answers: “Mum, if I work hard every evening, I will pass. But if the teachers give us a difficult paper, I will need your prayers!” His little sister adds: “If Junior passes, I will eat all the cake!” Everybody laughs. In this family, like in many Cameroonian families, success at school is a real possibility, and everybody is already planning the celebration.
\end{textebox}

\textbf{Glossaire :} \emph{guests} (n., invités) ; \emph{goat} (n., chèvre) ; \emph{paper} (n., ici : épreuve d'examen) ; \emph{prayers} (n., prières).

\textbf{Questions d'observation :}
\begin{enumerate}
\item Combien de phrases avec \eng{if} trouvez-vous dans ce texte ?
\item Quel temps verbal suit \eng{if} ? Quel auxiliaire apparaît dans l'autre proposition ?
\item Les événements décrits sont-ils impossibles, ou bien réellement envisageables ?
\end{enumerate}

Réponses attendues : toutes les propositions en \eng{if} sont au \textbf{present simple} (\eng{passes, come, have, rains, work, give}), et toutes les propositions principales contiennent \textbf{will + base verbale} (\eng{will organise, will help, will buy, will move, will pass, will need, will eat}). Les événements sont tout à fait possibles : Junior peut réussir, les invités peuvent venir tôt, il peut pleuvoir samedi. C'est exactement cela, le first conditional.

\courssub{La règle : forme et emploi}

\begin{propriete}[First conditional : la forme]
Le \cle{first conditional} exprime une hypothèse \textbf{réelle ou probable} concernant le \textbf{futur}. Sa structure est :
\begin{center}
\textbf{If + present simple, ... will + base verbale}
\end{center}
\begin{itemize}
\item La \textbf{if-clause} (proposition conditionnelle) est au \emph{present simple}, même si elle parle du futur.
\item La \textbf{main clause} (proposition principale) contient \emph{will} suivi de la base verbale (l'infinitif sans \eng{to}).
\item La condition est considérée comme possible : le locuteur pense qu'elle peut vraiment se réaliser.
\end{itemize}
\end{propriete}

\begin{exemplebox}[Exemples traduits]
\eng{If you study hard, you will pass the GCE.} « Si tu travailles dur, tu réussiras au GCE. »

\eng{If she doesn't hurry, she will miss the bus.} « Si elle ne se dépêche pas, elle ratera le car. »

\eng{If the rain stops, we will go to the market.} « Si la pluie s'arrête, nous irons au marché. »

\eng{If they win the match, the whole village will celebrate.} « S'ils gagnent le match, tout le village fera la fête. »

\eng{If I don't have money, I will not take a bendskin.} « Si je n'ai pas d'argent, je ne prendrai pas de bendskin. »
\end{exemplebox}

\begin{popfigure}
\begin{center}
\begin{tikzpicture}[node distance=0.6cm and 1.6cm]
\node[draw=popblue, fill=popblueL, rounded corners, thick, align=center, minimum width=4.2cm, minimum height=1.4cm] (ifc) {\textbf{If-clause}\\ If + present simple\\ \emph{If you study hard,}};
\node[draw=poporange, fill=poporangeL, rounded corners, thick, align=center, minimum width=4.2cm, minimum height=1.4cm, right=of ifc] (main) {\textbf{Main clause}\\ will + base verbale\\ \emph{you will pass.}};
\draw[-{Stealth[length=3mm]}, very thick, popdark] (ifc) -- node[above, align=center, font=\small\itshape, text=popdark]{futur probable} (main);
\end{tikzpicture}
\legende{Structure du first conditional : la condition au présent entraîne un résultat futur avec \eng{will}.}
\end{center}
\end{popfigure}

\begin{popfigure}
\begin{center}
\begin{tikzpicture}
\draw[-{Stealth[length=3mm]}, very thick, popline] (0,0) -- (10.5,0);
\foreach \x/\txt/\col in {1/{maintenant\\ (now)}/popdark, 5/{condition\\ (present simple)}/popblue, 9/{résultat futur\\ (will + BV)}/poporange}{
  \fill[\col] (\x,0) circle (0.12);
  \node[above, align=center, font=\small, text=\col] at (\x,0.15) {\txt};
}
\node[below, align=center, font=\small\itshape, text=popmuted] at (1,-0.2) {I am at home.};
\node[below, align=center, font=\small\itshape, text=popblue] at (5,-0.2) {If the rain stops,};
\node[below, align=center, font=\small\itshape, text=poporange] at (9,-0.2) {we will go out.};
\end{tikzpicture}
\legende{Frise chronologique : la condition est exprimée au présent mais porte sur l'avenir ; le résultat est futur.}
\end{center}
\end{popfigure}

\courssub{Le présent après if : une idée de futur}

Point capital, et excellente nouvelle pour les francophones : le français fonctionne exactement pareil. On dit « Si tu \textbf{viens} (présent), je \textbf{serai} (futur) content », et non « si tu viendras ». L'anglais suit la même logique : la if-clause est au présent, mais elle parle du futur.

\begin{center}
\begin{tabularx}{\linewidth}{|X|X|X|}
\hline
\rowcolor{popblueL} Français & English & Remarque \\
\hline
Si tu viens, je serai content. & If you come, I will be happy. & Présent après \eng{if} dans les deux langues. \\
\hline
S'il pleut, nous resterons à la maison. & If it rains, we will stay at home. & Jamais de futur dans la if-clause. \\
\hline
Si elle travaille, elle réussira. & If she works, she will succeed. & 3\textsuperscript{e} personne : \eng{works} avec -s. \\
\hline
Si vous ne partez pas, vous le manquerez. & If you don't leave, you will miss him. & Négation : \eng{don't / doesn't}. \\
\hline
\end{tabularx}
\end{center}

\begin{attention}[L'erreur n°1 des francophones]
On ne met \textbf{jamais} \eng{will} dans la if-clause. La phrase \eng{If you will come, I will be happy} est \textbf{fausse}. Il faut dire : \eng{If you come, I will be happy.} « Si tu viens, je serai content. »

De même : \eng{If it will rain...} est faux ; dites \eng{If it rains, we will stay at home.} Retenez : \textbf{un seul \eng{will} par phrase}, et il est dans la main clause, pas après \eng{if}.
\end{attention}

\courssub{L'ordre des propositions et la virgule}

Les deux propositions peuvent changer de place. La règle de ponctuation est simple :
\begin{itemize}
\item \textbf{If-clause d'abord} : on met une virgule entre les deux. \eng{If the rain stops, we will go out.}
\item \textbf{Main clause d'abord} : \textbf{pas de virgule}. \eng{We will go out if the rain stops.} « Nous sortirons si la pluie s'arrête. »
\end{itemize}

\begin{exemplebox}[Inversion sans virgule]
\eng{I will call you if I arrive late.} « Je t'appellerai si j'arrive en retard. »

\eng{She will be angry if you forget her birthday.} « Elle sera fâchée si tu oublies son anniversaire. »

\eng{The students will pass if they revise their lessons.} « Les élèves réussiront s'ils révisent leurs leçons. »
\end{exemplebox}

\courssub{Variantes utiles de la main clause}

\eng{Will} n'est pas le seul mot possible dans la proposition principale. Selon le sens, on peut utiliser :

\begin{center}
\begin{tabularx}{\linewidth}{|l|X|X|}
\hline
\rowcolor{popblueL} Forme & Exemple & Sens \\
\hline
\eng{will + BV} & If you work, you \textbf{will pass}. & certitude, prédiction \\
\hline
\eng{can + BV} & If you finish early, you \textbf{can go} home. & possibilité, permission \\
\hline
\eng{may / might + BV} & If it rains, we \textbf{may stay} inside. & probabilité plus faible \\
\hline
Impératif & If you see Peter, \textbf{tell} him to call me. & conseil, instruction, ordre \\
\hline
\end{tabularx}
\end{center}

\begin{exemplebox}[Exemples traduits]
\eng{If you see Peter, tell him to call me.} « Si tu vois Peter, dis-lui de m'appeler. »

\eng{If you feel tired, go to bed early.} « Si tu te sens fatigué, couche-toi tôt. »

\eng{If we save enough money, we can travel to Buea.} « Si nous économisons assez d'argent, nous pourrons voyager à Buea. »

\eng{If the weather is bad, the match might be cancelled.} « Si le temps est mauvais, le match pourrait être annulé. »
\end{exemplebox}

\courssub{Unless = if... not}

Le mot \cle{unless} signifie « si... ne... pas », c'est-à-dire « à moins que » ou « sauf si ». Il introduit une condition négative : \eng{unless + présent} équivaut à \eng{if + présent négatif}.

\begin{exemplebox}[Unless en action]
\eng{Unless you work hard, you will fail.} « Si tu ne travailles pas dur, tu échoueras. » (= \eng{If you don't work hard, you will fail.})

\eng{Unless it rains, we will play football.} « À moins qu'il (ne) pleuve, nous jouerons au football. »

\eng{You will not improve unless you practise every day.} « Tu ne progresseras pas si tu ne t'entraînes pas tous les jours. »
\end{exemplebox}

\begin{methode}[Traduire « à moins que » en anglais]
\begin{enumerate}
\item Repérez la condition négative en français : « à moins qu'il \textbf{ne} pleuve ».
\item Choisissez \eng{unless} et supprimez la négation : le « ne » français disparaît.
\item Mettez le verbe anglais au \textbf{present simple} (jamais de subjonctif, jamais de \eng{will}) : \eng{unless it rains}.
\item Ajoutez la main clause avec \eng{will} : \eng{Unless it rains, we will play.}
\end{enumerate}
Vérification : « À moins qu'il ne pleuve, nous jouerons. » = \eng{Unless it rains, we will play.} On peut toujours contrôler avec \eng{if... not} : \eng{If it doesn't rain, we will play.}
\end{methode}

\begin{attention}[Pièges avec unless]
\begin{itemize}
\item Ne doublez pas la négation : \eng{Unless you don't work...} est faux. \eng{Unless} contient déjà l'idée négative.
\item Le subjonctif français disparaît en anglais : « à moins qu'il ne \emph{pleuve} » devient simplement \eng{unless it \emph{rains}}.
\item Pas de \eng{will} après \eng{unless}, exactement comme après \eng{if}.
\end{itemize}
\end{attention}

\courssub{Emplois fréquents du first conditional}

Le first conditional s'utilise partout dans la vie quotidienne. Voici ses quatre grands emplois :

\begin{itemize}
\item \textbf{Promesses} : \eng{If you pass, I will buy you a phone.} « Si tu réussis, je t'achèterai un téléphone. »
\item \textbf{Avertissements} : \eng{If you touch that wire, you will get a shock.} « Si tu touches ce fil, tu prendras une décharge. »
\item \textbf{Prédictions} : \eng{If Cameroon plays well, we will win the cup.} « Si le Cameroun joue bien, nous gagnerons la coupe. »
\item \textbf{Conseils pratiques} : \eng{If you are ill, see a doctor.} « Si tu es malade, consulte un médecin. »
\end{itemize}

\begin{exoresolu}[Mettez les verbes à la forme correcte]
Énoncé : complétez avec le present simple ou \eng{will} + base verbale.

1. If my father (come) home early, we (eat) together.

2. She (not / catch) the bus unless she (run).

3. If they (invite) us, we (bring) some drinks.

4. I (help) you with your homework if you (ask) me politely.

\tcblower
Solution détaillée :

1. \eng{If my father \textbf{comes} home early, we \textbf{will eat} together.} « Si mon père rentre tôt, nous mangerons ensemble. » If-clause au present simple (\eng{comes}, 3\textsuperscript{e} personne avec -s), main clause avec \eng{will}.

2. \eng{She \textbf{will not catch} the bus unless she \textbf{runs}.} « Elle n'attrapera pas le bus à moins qu'elle ne coure. » Après \eng{unless}, present simple affirmatif (\eng{runs}) ; la négation est dans la main clause.

3. \eng{If they \textbf{invite} us, we \textbf{will bring} some drinks.} « S'ils nous invitent, nous apporterons des boissons. »

4. \eng{I \textbf{will help} you with your homework if you \textbf{ask} me politely.} « Je t'aiderai pour tes devoirs si tu me le demandes poliment. » Main clause d'abord : pas de virgule.
\end{exoresolu}

\begin{aretenir}[L'essentiel du first conditional]
\begin{itemize}
\item Structure : \textbf{If + present simple, will + base verbale} ; la condition est réelle ou probable.
\item Jamais de \eng{will} après \eng{if} (ni après \eng{unless}).
\item Inversion possible : main clause d'abord = pas de virgule.
\item Variantes : \eng{can}, \eng{may}, \eng{might} ou l'impératif dans la main clause.
\item \eng{Unless} = \eng{if... not} = « à moins que ».
\end{itemize}
\end{aretenir}

Maintenant que l'hypothèse réelle est bien maîtrisée, nous verrons dans la section suivante comment exprimer l'\textbf{irréel du présent} avec le second conditional : ce qui se passerait si les choses étaient différentes.

\coursec{Second conditional : l'irréel du présent}

Dans la section précédente, nous avons étudié le \cle{first conditional} : \eng{If it rains, we will stay at home.} La pluie est une possibilité réelle, le ciel est gris, cela peut vraiment arriver. Mais que se passe-t-il quand l'hypothèse n'est \emph{pas} réelle ? Quand nous rêvons, imaginons, ou donnons un conseil ? C'est exactement le rôle du \cle{second conditional}, que nous allons maintenant découvrir pas à pas.

\courssub{Observation : un dialogue au lycée}

\begin{textebox}[A conversation between two classmates in Yaoundé]
\textbf{Marlyse:} “You look sad, Junior. What's the matter?”

\textbf{Junior:} “I argued with my father yesterday. He was very angry.”

\textbf{Marlyse:} “If I were you, I would apologise to him. He loves you.”

\textbf{Junior:} “You're right. But he never listens to me.”

\textbf{Marlyse:} “If he listened to you, what would you tell him?”

\textbf{Junior:} “I would tell him that I want to study music, not medicine.”

\textbf{Marlyse:} “Music! If I had your talent, I would sing in every church in the city!”

\textbf{Junior:} “And if we were rich, we would build a big studio in Bastos!”

\textbf{Marlyse:} “Stop dreaming! If we studied more, we could pass the Probatoire first.”

\textbf{Junior:} “OK, OK. But dreaming is free, Marlyse!”
\end{textebox}

\textbf{Glossaire.} \emph{to argue} (v.) : se disputer ; \emph{to apologise} (v.) : s'excuser ; \emph{talent} (n.) : talent ; \emph{to dream} (v.) : rêver ; \emph{free} (adj.) : gratuit.

\textbf{Questions d'observation.}
\begin{enumerate}
\item Soulignez toutes les phrases qui commencent par \eng{If}. Quel temps verbal suit \eng{if} ?
\item Quel auxiliaire trouve-t-on dans l'autre partie de la phrase ?
\item Junior est-il riche ? Marlyse est-elle Junior ? Ces situations sont-elles réelles ?
\end{enumerate}

\textbf{Réponses attendues.} Après \eng{if}, on trouve le \cle{past simple} (\eng{were, had, listened, studied}) ; dans l'autre proposition, on trouve \eng{would} (ou \eng{could}) suivi de la base verbale. Et surtout : Junior n'est \emph{pas} riche, Marlyse n'est \emph{pas} Junior. Toutes ces situations sont \textbf{imaginaires}, contraires à la réalité présente.

\courssub{La règle : forme et nature du « passé d'irréel »}

\begin{propriete}[Second conditional : la structure]
Le \cle{second conditional} se construit ainsi :
\begin{center}
\textbf{If + sujet + past simple, \quad sujet + would + base verbale}
\end{center}
\eng{If I had a million francs, I would build a school in my village.}

Points essentiels :
\begin{itemize}
\item La proposition avec \eng{if} (la \emph{if-clause}) est au \textbf{past simple} : \eng{had, were, lived, knew}.
\item La proposition principale contient \textbf{would + base verbale} (l'infinitif sans \eng{to}) : \eng{would build, would talk, would go}.
\item L'ordre des deux propositions peut s'inverser, sans virgule alors : \eng{I would build a school if I had a million francs.}
\end{itemize}
\end{propriete}

\begin{definition}[Le « passé d'irréel »]
On appelle \cle{passé d'irréel} la forme du past simple utilisée \textbf{sans valeur de temps passé}, pour parler d'une situation \textbf{contraire à la réalité présente}. Dans \eng{If I were rich...}, le verbe \eng{were} est au passé, mais la phrase parle du \textbf{présent} : « maintenant, je ne suis pas riche ». Le passé est ici un \emph{marqueur de distance par rapport à la réalité}, pas un marqueur de temps. C'est exactement le rôle de l'imparfait français dans « Si j'\emph{étais} riche... » : « étais » est un imparfait, mais on parle d'aujourd'hui.
\end{definition}

\begin{aretenir}[Le parallèle français / anglais]
Bonne nouvelle : le second conditional fonctionne presque comme le français !
\begin{center}
\begin{tabularx}{\linewidth}{|X|X|X|}
\hline
\rowcolor{popblueL} Français & English & Remarque \\
\hline
Si j'\textbf{avais} de l'argent & If I \textbf{had} money & imparfait = past simple \\
\hline
je t'\textbf{aiderais} & I \textbf{would help} you & conditionnel = would + base verbale \\
\hline
Si j'\textbf{étais} toi & If I \textbf{were} you & « étais » = were (passé d'irréel) \\
\hline
\end{tabularx}
\end{center}
Le français utilise \emph{imparfait + conditionnel} ; l'anglais utilise \emph{past simple + would}. La logique est identique : il suffit de ne pas traduire mot à mot.
\end{aretenir}

\courssub{Emploi 1 : la situation imaginaire ou impossible au présent}

On utilise le second conditional pour imaginer une réalité différente de la réalité actuelle : un rêve, une hypothèse impossible ou très improbable.

\begin{exemplebox}[Situations imaginaires]
\begin{itemize}
\item \eng{If I had a million francs, I would build a school in my village.} « Si j'avais un million de francs, je construirais une école dans mon village. » (Je n'ai pas un million de francs.)
\item \eng{If we lived in Buea, we would see Mount Cameroon every day.} « Si nous habitions à Buea, nous verrions le mont Cameroun tous les jours. » (Nous n'habitons pas à Buea.)
\item \eng{If she spoke English fluently, she would get the job.} « Si elle parlait anglais couramment, elle obtiendrait le poste. » (Elle ne le parle pas couramment.)
\item \eng{If I were a bird, I would fly over the Wouri bridge.} « Si j'étais un oiseau, je survolerais le pont du Wouri. » (Impossible : je ne suis pas un oiseau.)
\end{itemize}
\end{exemplebox}

\courssub{Emploi 2 : le conseil poli — « If I were you... »}

Le second conditional sert aussi à donner un \textbf{conseil} de manière douce et polie. La structure figée est \eng{If I were you, I would + base verbale}.

\begin{exemplebox}[Donner un conseil]
\begin{itemize}
\item \eng{If I were you, I would talk to the teacher.} « Si j'étais toi, je parlerais au professeur. »
\item \eng{If I were you, I would apologise to your father.} « Si j'étais toi, je m'excuserais auprès de ton père. »
\item \eng{If I were him, I would not sell the land.} « Si j'étais lui, je ne vendrais pas le terrain. »
\end{itemize}
C'est l'équivalent exact du français « Si j'étais toi, je... ». Cette tournure est plus polie qu'un impératif direct : \eng{Apologise!} est brutal, \eng{If I were you, I would apologise} est amical.
\end{exemplebox}

\courssub{La règle spéciale : « were » pour toutes les personnes}

\begin{propriete}[If I were, if she were]
Dans la langue soignée — et donc \textbf{à l'écrit et à l'examen} — on emploie \cle{were} avec \textbf{toutes} les personnes dans la \emph{if-clause} du second conditional, y compris à la première et à la troisième personne du singulier :
\begin{center}
\begin{tabular}{|l|l|}
\hline
\rowcolor{popblueL} Personne & Forme soignée \\
\hline
I & If I \textbf{were} rich... \\
\hline
he / she / it & If she \textbf{were} here... \\
\hline
we / you / they & If they \textbf{were} ready... \\
\hline
\end{tabular}
\end{center}
\eng{If I were you} (et non \eng{If I was you}) est la forme attendue dans une rédaction ou au Probatoire. \eng{If I was you} s'entend dans la conversation familière, mais il est déconseillé à l'écrit.
\end{propriete}

\courssub{Variante : could et might à la place de would}

Dans la proposition principale, \eng{would} peut être remplacé par :
\begin{itemize}
\item \cle{could} (= « pourrais », capacité ou possibilité) : \eng{If I had time, I could help you.} « Si j'avais le temps, je pourrais t'aider. »
\item \cle{might} (= « il se pourrait que », possibilité incertaine) : \eng{If you studied harder, you might pass the exam.} « Si tu travaillais davantage, tu pourrais peut-être réussir l'examen. »
\end{itemize}
La logique reste la même : past simple dans la \emph{if-clause}, modal + base verbale dans la principale.

\courssub{First ou second conditional ? Une question de réalité}

La différence entre les deux structures n'est pas une différence de \emph{temps}, mais de \emph{degré de réalité} :

\begin{center}
\begin{tabularx}{\linewidth}{|X|X|}
\hline
\rowcolor{popblueL} First conditional (réel, possible) & Second conditional (imaginaire, improbable) \\
\hline
\eng{If it rains, we will stay at home.} (Le ciel est gris : c'est possible.) & \eng{If it rained, we would stay at home.} (Il fait beau : peu probable.) \\
\hline
\eng{If I pass the Probatoire, my father will buy me a phone.} (Je travaille : c'est envisageable.) & \eng{If I passed the Probatoire, my father would buy me a phone.} (Je doute de réussir.) \\
\hline
\end{tabularx}
\end{center}

\begin{popfigure}
\begin{tikzpicture}[
  box/.style={draw=popblue, fill=popblueL, rounded corners, align=center, text width=5.6cm, inner sep=6pt},
  lab/.style={align=center, font=\bfseries, text=popdark}
]
\node[lab] (t1) at (0,2.2) {First conditional : RÉEL};
\node[box, align=center] (f1) at (0,0.8){If it \textbf{rains},\\ we \textbf{will stay} at home.};
\node[align=center, text width=5.6cm, font=\small] at (0,-0.6) {La pluie est possible.\\ « S'il pleut, nous resterons. »};
\node[lab] (t2) at (8,2.2) {Second conditional : IMAGINAIRE};
\node[box, draw=poporange, fill=poporangeL, align=center] (f2) at (8,0.8){If it \textbf{rained},\\ we \textbf{would stay} at home.};
\node[align=center, text width=5.6cm, font=\small] at (8,-0.6) {La pluie est improbable.\\ « S'il pleuvait, nous resterions. »};
\draw[-{Stealth}, very thick, poppurple] (3.2,0.8) -- (4.8,0.8) node[midway, above, font=\small\itshape, text=poppurple] {même scène, autre réalité};
\end{tikzpicture}
\legende{Deux phrases miroirs : la même scène, mais un degré de réalité différent.}
\end{popfigure}

\begin{popfigure}
\begin{tikzpicture}[mindmap, grow cyclic,
  every node/.style={concept, align=center},
  root concept/.append style={concept color=popblue, font=\bfseries, text=white},
  level 1/.append style={sibling angle=90, level distance=3.4cm},
  level 1 concept/.append style={concept color=poporange, text=popdark, font=\small}]
\node[root concept, align=center]{Second\\ conditional}
  child { node[align=center]{Rêve\\ \emph{If I had money, I would travel.}} }
  child { node[align=center]{Conseil\\ \emph{If I were you, I would apologise.}} }
  child { node[align=center]{Impossibilité\\ \emph{If I were a bird, I would fly.}} }
  child { node[align=center]{Politesse\\ \emph{If you came, I might cook eru.}} };
\end{tikzpicture}
\legende{Carte mentale : les quatre grandes fonctions du second conditional.}
\end{popfigure}

\begin{exoresolu}[Mettez les verbes à la bonne forme]
Énoncé. Complétez avec le second conditional :
\begin{enumerate}
\item If I (be) rich, I (buy) a house in Douala.
\item If she (have) a bicycle, she (not / walk) to school.
\item If we (live) near the stadium, we (can) watch every match.
\end{enumerate}
\tcblower
Solution détaillée.
\begin{enumerate}
\item \eng{If I \textbf{were} rich, I \textbf{would buy} a house in Douala.} — \emph{If-clause} au passé d'irréel : \eng{were} (forme soignée avec \eng{I}) ; principale : \eng{would} + base verbale \eng{buy}. « Si j'étais riche, j'achèterais une maison à Douala. »
\item \eng{If she \textbf{had} a bicycle, she \textbf{would not walk} to school.} — Past simple de \eng{have} : \eng{had} ; négation : \eng{would not} + \eng{walk}. « Si elle avait un vélo, elle n'irait pas à l'école à pied. »
\item \eng{If we \textbf{lived} near the stadium, we \textbf{could watch} every match.} — Past simple de \eng{live} : \eng{lived} ; variante avec \eng{could} + base verbale. « Si nous habitions près du stade, nous pourrions regarder chaque match. »
\end{enumerate}
\end{exoresolu}

\begin{attention}[Les pièges des francophones]
\begin{enumerate}
\item \textbf{Jamais de \eng{would} après \eng{if} !} Ne dites pas \emph{If I would have money...}. La règle d'or : \eng{would} vit dans la proposition principale, jamais dans la \emph{if-clause}. On écrit \eng{If I \textbf{had} money, I \textbf{would} help you.}
\item \textbf{\eng{If I was you} est familier.} À l'écrit et à l'examen, écrivez toujours \eng{If I \textbf{were} you}.
\item \textbf{Ne traduisez pas le conditionnel français par un futur.} « Je construirais » = \eng{I would build}, jamais \eng{I will build} dans cette structure.
\item \textbf{Le passé n'est pas du passé.} \eng{If I knew the answer} signifie « si je savais la réponse \emph{maintenant} », pas « si j'avais su ».
\end{enumerate}
\end{attention}

\begin{savaistu}[« If I Were a Boy »]
La chanson \emph{If I Were a Boy} de la chanteuse américaine Beyoncé est un exemple parfait du second conditional : elle imagine ce qu'elle ferait si elle était un garçon — situation évidemment impossible. Remarquez qu'elle emploie \eng{were} et non \eng{was} : même dans la musique pop, la forme soignée s'impose quand on veut bien s'exprimer. Écoutez la chanson et repérez tous les \eng{would} !
\end{savaistu}

\begin{exercice}[À vous de jouer]
\begin{enumerate}
\item Traduisez : « Si j'étais toi, je réviserais mes leçons tous les soirs. »
\item Complétez : If my village (have) electricity, the students (study) better at night.
\item Choisissez entre first et second conditional : If it (rain) tomorrow, the match (be) cancelled. — La météo annonce de la pluie.
\item Écrivez deux phrases sur vos rêves avec \eng{If I had...} et \eng{If I were...}.
\end{enumerate}
\end{exercice}

\begin{corrige}[Exercice « À vous de jouer »]
\begin{enumerate}
\item \eng{If I were you, I would revise my lessons every evening.}
\item \eng{If my village \textbf{had} electricity, the students \textbf{would study} better at night.}
\item First conditional, car la pluie est annoncée (réelle) : \eng{If it \textbf{rains} tomorrow, the match \textbf{will be} cancelled.}
\item Réponses libres, par exemple : \eng{If I had a camera, I would photograph the animals of Waza Park.} / \eng{If I were the mayor, I would clean the streets of my town.}
\end{enumerate}
\end{corrige}

\coursec{Third conditional : le regret du passé}

Dans la section précédente, nous avons étudié le \cle{second conditional}, qui sert à imaginer une situation irréelle dans le présent ou le futur : \eng{If I were rich, I would travel the world.} « Si j'étais riche, je voyagerais à travers le monde. » Mais que se passe-t-il quand la situation irréelle se situe dans le \cle{passé}, un passé révolu qu'on ne peut plus modifier ? C'est exactement le rôle du \cle{third conditional}, la structure du regret, du reproche et du « si seulement... ».

\courssub{Observation : le match perdu}

Lisons d'abord ce court texte. Le lycée de Buea vient de perdre un match de football important. Le lendemain, dans la cour de récréation, les élèves commentent la défaite.

\begin{textebox}[After the match — a conversation in Buea]
The students of Government High School Buea were very disappointed on Monday morning. Their team had lost the final 2--1 against a school from Limbe.

“If Ngong had trained regularly, the team would have won the match,” said Tabi, the captain. “He missed two easy goals.”

“And if the referee had seen the foul, he would have given a penalty,” added Manka. “Everybody saw it, except him!”

“If we had left home earlier, we would have arrived before the rain started,” said Enow. “The pitch was terrible in the second half.”

Their teacher, Mr Efange, smiled. “If I had known you would analyse the match so well in English, I would have organised the debate in class! Listen: the match is over. You cannot change the past. But you can prepare the future. If you train hard this term, you will win the cup next year.”

The students laughed. Tabi concluded: “If we had listened to Mr Efange last term, we wouldn't have lost!”
\end{textebox}

\textbf{Glossaire.} \emph{referee} (n.) : arbitre ; \emph{foul} (n.) : faute (au football) ; \emph{pitch} (n.) : terrain ; \emph{cup} (n.) : coupe ; \emph{term} (n.) : trimestre ; \emph{to miss a goal} (v.) : rater un but.

\textbf{Questions de compréhension.}
\begin{enumerate}
\item Why were the students disappointed?
\item What did Tabi say about Ngong?
\item What would the referee have done if he had seen the foul?
\item What advice does Mr Efange give to the students?
\end{enumerate}

\textbf{Observation grammaticale.} Soulignons les structures répétées :

\eng{If Ngong \textbf{had trained} regularly, the team \textbf{would have won} the match.}

\eng{If the referee \textbf{had seen} the foul, he \textbf{would have given} a penalty.}

\eng{If we \textbf{had listened} to Mr Efange, we \textbf{wouldn't have lost}.}

On observe deux blocs constants : après \eng{if}, on trouve \eng{had} + participe passé ; dans la proposition principale, on trouve \eng{would have} + participe passé. Et surtout : toutes ces hypothèses portent sur un passé terminé, impossible à changer. Le match est fini. Ngong ne s'est pas entraîné. L'arbitre n'a pas vu la faute. C'est trop tard.

\courssub{La règle : forme et emploi}

\begin{propriete}[Third conditional : la structure]
Le \cle{third conditional} exprime une hypothèse \textbf{contraire à la réalité passée}. Sa structure est :

\begin{center}
\textbf{If + sujet + had + participe passé, sujet + would have + participe passé.}
\end{center}

\begin{itemize}
\item Proposition subordonnée : \eng{if} + \cle{past perfect} (\eng{had} + participe passé).
\item Proposition principale : \eng{would have} + participe passé.
\end{itemize}

La condition n'a \textbf{pas été réalisée} dans le passé : elle est donc irréelle, imaginaire, et le résultat décrit ne s'est jamais produit. Comme pour les autres conditionnels, les deux propositions peuvent être inversées ; dans ce cas, on supprime la virgule : \eng{The team would have won if Ngong had trained regularly.}
\end{propriete}

\begin{exemplebox}[Exemples traduits]
\eng{If you had revised, you would have passed.} « Si tu avais révisé, tu aurais réussi. » (Réalité : tu n'as pas révisé, tu as échoué.)

\eng{If we had left earlier, we would have caught the bus.} « Si nous étions partis plus tôt, nous aurions attrapé le car. » (Réalité : nous sommes partis tard, nous avons raté le car.)

\eng{If she had studied medicine, she would have become a doctor.} « Si elle avait fait médecine, elle serait devenue médecin. »

\eng{If they had booked their tickets, they would have travelled to Douala.} « S'ils avaient réservé leurs billets, ils seraient allés à Douala. »
\end{exemplebox}

\textbf{Emplois du third conditional.}
\begin{enumerate}
\item \textbf{Le regret} : \eng{If I had worked harder, I would have succeeded.} « Si j'avais travaillé plus dur, j'aurais réussi. »
\item \textbf{Le reproche} : \eng{If you had told me the truth, I would have forgiven you.} « Si tu m'avais dit la vérité, je t'aurais pardonné. »
\item \textbf{Le soulagement} (avec la négation) : \eng{If we had taken that bendskin, we would have had an accident!} « Si nous avions pris ce bendskin, nous aurions eu un accident ! » (Heureusement, nous ne l'avons pas pris.)
\item \textbf{L'analyse d'un événement passé} : échec à un examen, défaite sportive, histoire hypothétique : \eng{If Cameroon had scored that penalty, the team would have qualified.} « Si le Cameroun avait marqué ce penalty, l'équipe se serait qualifiée. »
\end{enumerate}

\begin{popfigure}
\begin{tikzpicture}[scale=1]
\draw[thick, -{Stealth}, popdark] (0,0) -- (12,0);
\fill[poporange] (2.5,0) circle (0.12);
\node[align=center, above] at (2.5,0.25) {\textbf{Passé révolu}\\ condition non réalisée};
\draw[thick, poporange] (1.7,-0.9) -- (3.3,-0.35);
\draw[thick, poporange] (1.7,-0.35) -- (3.3,-0.9);
\node[align=center, below, poporange] at (2.5,-1.0) {\eng{If I had known...}\\ (je ne savais pas)};
\fill[popblue] (9,0) circle (0.12);
\node[align=center, above] at (9,0.25) {\textbf{Présent}\\ le locuteur regrette};
\node[align=center, below, popblue] at (9,-0.6) {\eng{...I would have told you.}};
\draw[thick, -{Stealth}, poppurple] (3.6,0.9) -- (8.2,0.9);
\node[align=center, above, poppurple] at (5.9,0.95) {impossible à changer};
\end{tikzpicture}
\legende{Le third conditional : la condition appartient à un passé fermé ; seul le regret existe au présent.}
\end{popfigure}

\courssub{Rappel : la formation du past perfect}

Le third conditional exige le \cle{past perfect} après \eng{if}. Rappelons sa formation, très simple car invariable.

\begin{propriete}[Past perfect : had + participe passé]
Le past perfect se forme avec \eng{had} (identique à \textbf{toutes} les personnes) + le participe passé du verbe :

\begin{center}
\eng{I / you / he / she / it / we / they had worked, had gone, had seen...}
\end{center}

Négation : \eng{had not (hadn't) + participe passé}. Interrogation : \eng{Had + sujet + participe passé...?}
\end{propriete}

\begin{center}
\begin{tabular}{|l|l|l|l|}
\hline
\rowcolor{popblueL} Verbe & Base verbale & Past perfect & Traduction \\
\hline
work & work & had worked & avait travaillé \\
\hline
go & go & had gone & était allé(e) \\
\hline
see & see & had seen & avait vu \\
\hline
take & take & had taken & avait pris \\
\hline
know & know & had known & avait su \\
\hline
leave & leave & had left & était parti(e) \\
\hline
catch & catch & had caught & avait attrapé \\
\hline
write & write & had written & avait écrit \\
\hline
\end{tabular}
\end{center}

\begin{attention}[Verbes irréguliers : le participe passé]
La seule difficulté du past perfect est le participe passé des verbes irréguliers. Révisez la liste : \eng{go -- went -- gone}, \eng{see -- saw -- seen}, \eng{take -- took -- taken}, \eng{write -- wrote -- written}, \eng{speak -- spoke -- spoken}, \eng{eat -- ate -- eaten}, \eng{begin -- began -- begun}. Ne dites jamais \eng{had went} ou \eng{had saw} : c'est une faute grave. Le past perfect exige toujours le \textbf{participe passé} (troisième colonne), jamais le prétérit (deuxième colonne).
\end{attention}

\courssub{Variantes : could have et might have}

À la place de \eng{would have} dans la proposition principale, on peut utiliser d'autres modaux pour nuancer le sens.

\begin{center}
\begin{tabular}{|l|l|l|}
\hline
\rowcolor{popblueL} Structure & Sens & Exemple \\
\hline
\eng{would have + p.p.} & résultat certain & \eng{If you had asked, I would have helped.} \\
\hline
\eng{could have + p.p.} & possibilité, capacité & \eng{If you had asked, I could have helped.} \\
\hline
\eng{might have + p.p.} & possibilité incertaine & \eng{If you had asked, I might have helped.} \\
\hline
\end{tabular}
\end{center}

\begin{exemplebox}[Exemples traduits]
\eng{If you had asked, I could have helped.} « Si tu avais demandé, j'aurais pu t'aider. »

\eng{If she had left earlier, she might have caught the train.} « Si elle était partie plus tôt, elle aurait peut-être attrapé le train. »

\eng{If we had played better, we could have won the cup.} « Si nous avions mieux joué, nous aurions pu gagner la coupe. »
\end{exemplebox}

\textbf{La contraction orale.} À l'oral, les anglophones contractent presque toujours : \eng{had} devient \eng{'d} et \eng{would} devient aussi \eng{'d}. Attention au contexte pour les distinguer !

\eng{If I'd known, I'd have told you.} « Si j'avais su, je te l'aurais dit. »

Ici, \eng{I'd known} = \eng{I \textbf{had} known} (suivi d'un participe passé seul), et \eng{I'd have told} = \eng{I \textbf{would} have told} (suivi de \eng{have}). Règle simple : \eng{'d} + participe passé = \eng{had} ; \eng{'d have} = \eng{would have}.

\courssub{Exprimer un reproche : la méthode}

\begin{methode}[Construire un reproche au third conditional]
Pour reprocher à quelqu'un une faute passée, suivez ces étapes :
\begin{enumerate}
\item Identifiez la faute commise (ou l'action non faite) : \emph{you did not listen to me}.
\item Mettez-la au past perfect après \eng{if}, à l'affirmatif si l'action n'a pas été faite : \eng{If you had listened to me...}
\item Ajoutez la conséquence imaginaire avec \eng{would(n't) have} + participe passé : \eng{...you wouldn't have failed.}
\item Vérifiez : jamais de \eng{would} dans la proposition en \eng{if} !
\end{enumerate}
Résultat : \eng{If you had listened to me, you wouldn't have failed.} « Si tu m'avais écouté, tu n'aurais pas échoué. »
\end{methode}

Autres reproches de la vie familiale et scolaire :

\eng{If you had done your homework, the teacher wouldn't have punished you.} « Si tu avais fait tes devoirs, le professeur ne t'aurait pas puni. »

\eng{If you had saved your money, you would have bought that phone.} « Si tu avais économisé ton argent, tu aurais acheté ce téléphone. »

\eng{If you hadn't eaten all the mangoes, you wouldn't have been sick.} « Si tu n'avais pas mangé toutes les mangues, tu n'aurais pas été malade. »

\courssub{Les trois conditionnels : une même situation comparée}

Comparons les trois conditionnels sur une même situation : l'examen de fin d'année d'une élève de Première.

\begin{center}
\begin{tabular}{|l|l|l|}
\hline
\rowcolor{popblueL} Conditionnel & Phrase & Sens \\
\hline
First & \eng{If she works hard, she will pass.} & réel : l'examen est devant elle \\
\hline
Second & \eng{If she worked hard, she would pass.} & irréel du présent : elle ne travaille pas \\
\hline
Third & \eng{If she had worked hard, she would have passed.} & irréel du passé : trop tard, elle a échoué \\
\hline
\end{tabular}
\end{center}

\begin{popfigure}
\begin{tikzpicture}[scale=1]
\node[draw, rounded corners, fill=popgreenL, align=center, minimum width=3.4cm] (first) at (0,0) {\textbf{First}\\ \eng{If she works,}\\ \eng{she will pass.}};
\node[draw, rounded corners, fill=poporangeL, align=center, minimum width=3.4cm] (second) at (0,-2.4) {\textbf{Second}\\ \eng{If she worked,}\\ \eng{she would pass.}};
\node[draw, rounded corners, fill=poppinkL, align=center, minimum width=3.4cm] (third) at (0,-4.8) {\textbf{Third}\\ \eng{If she had worked,}\\ \eng{she would have passed.}};
\node[align=left, anchor=west] at (4.6,0) {futur possible};
\node[align=left, anchor=west] at (4.6,-2.4) {présent imaginaire};
\node[align=left, anchor=west] at (4.6,-4.8) {passé perdu : regret};
\draw[thick, -{Stealth}, popmuted] (2.2,0) -- (4.4,0);
\draw[thick, -{Stealth}, popmuted] (2.2,-2.4) -- (4.4,-2.4);
\draw[thick, -{Stealth}, popmuted] (2.2,-4.8) -- (4.4,-4.8);
\end{tikzpicture}
\legende{La même situation « examen » déclinée en trois conditionnels : plus on descend, plus l'hypothèse s'éloigne de la réalité.}
\end{popfigure}

Comparaison avec le français : la correspondance est ici presque parfaite, ce qui est une bonne nouvelle pour les francophones.

\begin{center}
\begin{tabularx}{\linewidth}{|X|X|X|}
\hline
\rowcolor{popblueL} Français & English & Remarque \\
\hline
Si tu avais révisé (plus-que-parfait) & If you had revised (past perfect) & correspondance directe \\
\hline
tu aurais réussi (conditionnel passé) & you would have passed & \eng{would have} + participe passé \\
\hline
Si j'aurais su (incorrect en français soigné) & If I would have known (impossible !) & jamais \eng{would} après \eng{if} \\
\hline
\end{tabularx}
\end{center}

\begin{attention}[Le piège n° 1 : jamais \eng{would} après \eng{if}]
Beaucoup de francophones, influencés par le « si j'aurais su » du français familier, écrivent : \eng{If I would have known, I would have told you.} C'est une \textbf{faute grave} en anglais. Après \eng{if}, on utilise le past perfect, jamais \eng{would have} :

\eng{If I \textbf{had} known, I would have told you.} « Si j'avais su, je te l'aurais dit. »

Autre piège : ne pas confondre \eng{had} et \eng{would} dans la contraction \eng{'d}. \eng{If I'd studied} = \eng{had} ; \eng{I'd have passed} = \eng{would}.
\end{attention}

\courssub{Entraînement : exercice résolu}

\begin{exoresolu}[Mets les verbes au third conditional]
Énoncé. Complète les phrases avec la forme correcte des verbes entre parenthèses.

1. If Ngong \_\_\_\_\_\_\_\_ (train) regularly, the team \_\_\_\_\_\_\_\_ (win) the match.

2. If they \_\_\_\_\_\_\_\_ (leave) earlier, they \_\_\_\_\_\_\_\_ (not / miss) the bus to Yaoundé.

3. If I \_\_\_\_\_\_\_\_ (know) about the test, I \_\_\_\_\_\_\_\_ (revise) my lessons.

4. If you \_\_\_\_\_\_\_\_ (ask) me, I \_\_\_\_\_\_\_\_ (can / help) you.

\tcblower
Solution détaillée.

1. \eng{If Ngong \textbf{had trained} regularly, the team \textbf{would have won} the match.} Après \eng{if} : past perfect (\eng{had} + participe passé \eng{trained}) ; proposition principale : \eng{would have} + \eng{won} (participe passé irrégulier de \eng{win} : \eng{win -- won -- won}).

2. \eng{If they \textbf{had left} earlier, they \textbf{wouldn't have missed} the bus to Yaoundé.} Participe passé irrégulier : \eng{leave -- left -- left}. La conséquence est négative : \eng{would not have missed}.

3. \eng{If I \textbf{had known} about the test, I \textbf{would have revised} my lessons.} Attention : \eng{know -- knew -- known}. Jamais \eng{If I would have known} !

4. \eng{If you \textbf{had asked} me, I \textbf{could have helped} you.} Variante avec \eng{could have} : la capacité d'aider existait, mais la demande n'a jamais été faite.
\end{exoresolu}

\begin{exercice}[À ton tour]
1. Traduis en anglais : « Si nous avions écouté le professeur, nous aurions réussi l'examen. »

2. Transforme en reproche au third conditional : \emph{You did not save money. You could not buy the book.}

3. Complète : If the referee \_\_\_\_\_\_\_\_ (see) the foul, he \_\_\_\_\_\_\_\_ (give) a penalty.

4. Écris deux phrases au third conditional pour exprimer un regret personnel (examen, match, voyage).
\end{exercice}

\begin{corrige}[Exercice « À ton tour »]
1. \eng{If we had listened to the teacher, we would have passed the exam.}

2. \eng{If you had saved money, you would have bought the book.} (ou \eng{...you could have bought the book.})

3. \eng{If the referee \textbf{had seen} the foul, he \textbf{would have given} a penalty.} Participe passé de \eng{see} : \eng{seen} ; de \eng{give} : \eng{given}.

4. Exemples possibles : \eng{If I had revised mathematics, I would have had a better mark.} / \eng{If we had trained every day, we would have won the final.}
\end{corrige}

\begin{aretenir}[L'essentiel du third conditional]
\begin{itemize}
\item Structure : \eng{If + had + participe passé, ... would have + participe passé.}
\item Emploi : hypothèse contraire à la réalité passée ; regret, reproche, soulagement.
\item Variantes : \eng{could have} (capacité), \eng{might have} (possibilité incertaine).
\item Oral : \eng{If I'd known, I'd have told you.}
\item Interdit absolu : \eng{would have} après \eng{if}. On dit \eng{If I had known}, jamais \eng{If I would have known}.
\end{itemize}
\end{aretenir}

\coursec{Synthèse, comparaison avec le français et pièges}

Nous venons d'étudier séparément les trois conditionnels : le \cle{first conditional} pour l'hypothèse réelle du futur, le \cle{second conditional} pour l'irréel du présent, et le \cle{third conditional} pour le regret d'un passé révolu. Le moment est venu de rassembler tout cela, de comparer l'anglais et le français, et surtout d'apprendre à éviter les pièges qui coûtent des points aux examens.

\courssub{Le tableau récapitulatif des trois conditionnels}

\begin{propriete}[Tableau de synthèse des trois conditionnels]
\begin{center}
\begin{tabularx}{\linewidth}{|l|X|X|X|}
\hline
\rowcolor{popblueL} \textbf{Type} & \textbf{If-clause} & \textbf{Main clause} & \textbf{Sens} \\
\hline
First conditional & if + present simple & will + base verbale & hypothèse réelle, futur possible \\
\hline
Second conditional & if + past simple & would + base verbale & présent imaginaire, irréel \\
\hline
Third conditional & if + past perfect & would have + participe passé & passé révolu, regret \\
\hline
\end{tabularx}
\end{center}
\end{propriete}

\begin{exemplebox}[Un exemple traduit pour chaque type]
\eng{If it rains tomorrow, we will stay at home.} « S'il pleut demain, nous resterons à la maison. » (first : cela peut vraiment arriver)

\eng{If I had a bendskin, I would take you to school every morning.} « Si j'avais une moto-taxi, je t'emmènerais à l'école chaque matin. » (second : je n'ai pas de moto, c'est imaginaire)

\eng{If you had left earlier, you would have caught the bus to Douala.} « Si tu étais parti plus tôt, tu aurais attrapé le bus pour Douala. » (third : tu n'es pas parti tôt, c'est fini)
\end{exemplebox}

\begin{popfigure}
\begin{center}
\begin{tikzpicture}[font=\small]
\draw[fill=popblueL, draw=popblue, rounded corners] (-2.2,0) rectangle (2.2,0.9);
\node[align=center] at (0,0.45) {\textbf{Type} : first};
\draw[fill=popblue!15, draw=popblue, rounded corners] (-2.2,-1.1) rectangle (2.2,-0.2);
\node[align=center] at (0,-0.65) {if + \textbf{present} / \textbf{will} + BV};
\draw[fill=poporangeL, draw=poporange, rounded corners] (2.8,0) rectangle (7.2,0.9);
\node[align=center] at (5,0.45) {\textbf{Type} : second};
\draw[fill=poporange!15, draw=poporange, rounded corners] (2.8,-1.1) rectangle (7.2,-0.2);
\node[align=center] at (5,-0.65) {if + \textbf{past simple} / \textbf{would} + BV};
\draw[fill=popgreenL, draw=popgreen, rounded corners] (7.8,0) rectangle (12.2,0.9);
\node[align=center] at (10,0.45) {\textbf{Type} : third};
\draw[fill=popgreen!15, draw=popgreen, rounded corners] (7.8,-1.1) rectangle (12.2,-0.2);
\node[align=center] at (10,-0.65) {if + \textbf{past perfect} / \textbf{would have} + p.p.};
\draw[->, thick, popmuted] (0,-1.1) -- (0,-1.8);
\draw[->, thick, popmuted] (5,-1.1) -- (5,-1.8);
\draw[->, thick, popmuted] (10,-1.1) -- (10,-1.8);
\node[align=center, text=popdark] at (0,-2.2) {futur possible};
\node[align=center, text=popdark] at (5,-2.2) {présent imaginaire};
\node[align=center, text=popdark] at (10,-2.2) {passé révolu};
\end{tikzpicture}
\legende{Grand tableau récapitulatif : structure et sens des trois conditionnels.}
\end{center}
\end{popfigure}

\courssub{Comparaison avec le français : un transfert positif}

Bonne nouvelle pour les francophones : le français et l'anglais construisent leurs phrases hypothétiques de façon \cle{très proche}. À chaque combinaison française correspond une combinaison anglaise presque symétrique.

\begin{center}
\begin{tabularx}{\linewidth}{|X|X|X|}
\hline
\rowcolor{popblueL} \textbf{Français} & \textbf{English} & \textbf{Remarque} \\
\hline
Si tu viens, je serai là. (si + présent / futur) & If you come, I will be there. & présent des deux côtés après si ; futur dans l'autre proposition \\
\hline
Si tu venais, je serais content. (si + imparfait / conditionnel) & If you came, I would be happy. & imparfait = past simple ; conditionnel = would \\
\hline
Si tu étais venu, j'aurais été content. (si + plus-que-parfait / conditionnel passé) & If you had come, I would have been happy. & plus-que-parfait = past perfect ; conditionnel passé = would have + p.p. \\
\hline
\end{tabularx}
\end{center}

\begin{aretenir}[Le transfert est positif]
La logique est la même dans les deux langues : plus l'hypothèse est irréelle ou lointaine, plus on « recule » dans les temps. Présent pour le possible, passé pour l'imaginaire, plus-que-parfait pour le révolu. Tu peux donc t'appuyer sur ton intuition francophone, à une condition absolue : respecter l'interdiction de \eng{will} et \eng{would} après \eng{if}.
\end{aretenir}

\begin{propriete}[Règle d'or : jamais de will ni de would dans la if-clause]
Dans la proposition introduite par \eng{if}, on n'emploie \textbf{jamais} \eng{will} ni \eng{would}. Le futur et le conditionnel se mettent uniquement dans la \cle{main clause}.

\eng{If it rains, I will stay home.} (correct) — \eng{If it will rain...} (faux)

Seule exception, hors programme : \eng{will} après \eng{if} avec le sens de « vouloir bien » : \eng{If you will wait a moment, the headmaster will see you.} « Si vous voulez bien attendre un instant, le proviseur vous recevra. »
\end{propriete}

\courssub{Comment choisir le bon conditionnel ?}

\begin{methode}[Trois questions pour choisir]
\begin{enumerate}
\item La situation est-elle \textbf{passée et terminée}, irréversible ? Alors j'utilise le \textbf{third conditional} : \eng{If I had known, I would have helped.}
\item La situation est-elle \textbf{imaginaire au présent}, contraire à la réalité actuelle ? Alors j'utilise le \textbf{second conditional} : \eng{If I were rich, I would travel.}
\item La situation est-elle \textbf{possible dans le futur} ? Alors j'utilise le \textbf{first conditional} : \eng{If I pass the exam, I will celebrate.}
\end{enumerate}
\end{methode}

\begin{popfigure}
\begin{center}
\begin{tikzpicture}[font=\small, node distance=0.6cm]
\node[draw=popdark, fill=popcream, rounded corners, align=center, text width=4.5cm] (start) {L'hypothèse porte sur quoi ?};
\node[draw=popgreen, fill=popgreenL, rounded corners, align=center, below left=1cm and -0.5cm of start, text width=3.4cm] (past) {le passé, c'est fini};
\node[draw=poporange, fill=poporangeL, rounded corners, align=center, below=1cm of start, text width=3.2cm] (present) {le présent, mais imaginaire};
\node[draw=popblue, fill=popblueL, rounded corners, align=center, below right=1cm and -0.5cm of start, text width=3.4cm] (future) {le futur, c'est possible};
\node[draw=popgreen, rounded corners, align=center, below=0.5cm of past, text width=3.4cm] (third) {\textbf{THIRD}\\ if + past perfect\\ would have + p.p.};
\node[draw=poporange, rounded corners, align=center, below=0.5cm of present, text width=3.2cm] (second) {\textbf{SECOND}\\ if + past simple\\ would + BV};
\node[draw=popblue, rounded corners, align=center, below=0.5cm of future, text width=3.4cm] (first) {\textbf{FIRST}\\ if + present\\ will + BV};
\draw[->, thick, popmuted] (start) -- (past);
\draw[->, thick, popmuted] (start) -- (present);
\draw[->, thick, popmuted] (start) -- (future);
\draw[->, thick, popgreen] (past) -- (third);
\draw[->, thick, poporange] (present) -- (second);
\draw[->, thick, popblue] (future) -- (first);
\end{tikzpicture}
\legende{Organigramme de choix : futur possible, présent imaginaire ou passé révolu ?}
\end{center}
\end{popfigure}

\courssub{Les quatre pièges à éviter}

\begin{attention}[Piège 1 : will ou would dans la if-clause]
C'est l'erreur la plus fréquente des francophones, car le français met le futur après « si »... non, justement : le français dit « s'il \emph{pleut} » (présent), jamais « s'il pleuvra ». Fais pareil en anglais.

Faux : \eng{If you will come, I will be happy.} Correct : \eng{If you come, I will be happy.}

Faux : \eng{If I would have money, I would buy a phone.} Correct : \eng{If I had money, I would buy a phone.}
\end{attention}

\begin{attention}[Piège 2 : confondre second et third conditional]
Demande-toi toujours : est-ce que je parle du présent ou du passé ?

\eng{If I lived in Buea, I would see the mountain every day.} (second : je n'habite pas à Buea maintenant)

\eng{If I had lived in Buea, I would have seen the mountain every day.} (third : je n'y ai jamais habité, c'est du passé révolu)

Mélanger les deux (\eng{If I had money, I would have bought...}) détruit le sens de la phrase.
\end{attention}

\begin{attention}[Piège 3 : « would have went » est une faute grave]
Après \eng{would have}, il faut \textbf{toujours} le participe passé : \eng{would have gone}, \eng{would have done}, \eng{would have been}, \eng{would have taken}. Jamais le prétérit.

Faux : \eng{She would have went to the market.} Correct : \eng{She would have gone to the market.}

Révise ta liste de verbes irréguliers : c'est ici qu'elle devient indispensable.
\end{attention}

\begin{attention}[Piège 4 : la virgule]
Quand la \eng{if-clause} est en \textbf{tête de phrase}, on met une virgule entre les deux propositions : \eng{If it rains, we will stay home.}

Quand la \eng{if-clause} est en \textbf{fin de phrase}, il n'y a \textbf{pas} de virgule : \eng{We will stay home if it rains.}
\end{attention}

\begin{exoresolu}[Choisis le bon conditionnel]
Énoncé : complète avec la forme correcte du verbe entre parenthèses.

1. If we (leave) now, we (arrive) before dark. (vrai projet pour ce soir)

2. If my father (be) a businessman, we (live) in Bastos. (ce n'est pas le cas)

3. If the students (listen) to the teacher, they (not fail) the test last week.
\tcblower
Solution détaillée :

1. Projet réel pour ce soir, futur possible : first conditional. \eng{If we leave now, we will arrive before dark.} « Si nous partons maintenant, nous arriverons avant la nuit. »

2. Situation contraire à la réalité présente : second conditional. \eng{If my father were a businessman, we would live in Bastos.} « Si mon père était homme d'affaires, nous habiterions à Bastos. » (\eng{were} pour toutes les personnes au second conditional.)

3. « Last week » : passé révolu, regret : third conditional. \eng{If the students had listened to the teacher, they would not have failed the test.} « Si les élèves avaient écouté le professeur, ils n'auraient pas échoué au contrôle. » Attention : \eng{would not have failed}, participe passé \eng{failed}, pas \eng{falled}.
\end{exoresolu}

\begin{exercice}[À toi de jouer]
Mets le verbe entre parenthèses à la forme correcte, puis traduis chaque phrase en français.

1. If you (study) harder, you (pass) the GCE next year.

2. If I (be) you, I (apologize) to my sister.

3. If they (take) the early bus, they (not miss) the match in Yaoundé.

4. We (go) to the beach in Kribi if the weather (be) fine tomorrow.

5. If she (have) a computer, she (do) her homework faster.
\end{exercice}

\begin{corrige}[Exercice « À toi de jouer »]
1. \eng{If you study harder, you will pass the GCE next year.} « Si tu étudies plus sérieusement, tu réussiras le GCE l'année prochaine. » (first : futur possible)

2. \eng{If I were you, I would apologize to my sister.} « Si j'étais toi, je m'excuserais auprès de ma sœur. » (second : conseil, situation imaginaire ; expression figée \eng{If I were you}.)

3. \eng{If they had taken the early bus, they would not have missed the match in Yaoundé.} « S'ils avaient pris le premier bus, ils n'auraient pas raté le match à Yaoundé. » (third : passé révolu ; participe passé \eng{taken}.)

4. \eng{We will go to the beach in Kribi if the weather is fine tomorrow.} « Nous irons à la plage de Kribi s'il fait beau demain. » (first ; pas de virgule car la if-clause est en fin de phrase.)

5. \eng{If she had a computer, she would do her homework faster.} « Si elle avait un ordinateur, elle ferait ses devoirs plus vite. » (second : elle n'a pas d'ordinateur maintenant.)
\end{corrige}

\courssub{Compléments Bac : conditionnels mixtes, I wish et If only}

\begin{definition}[Le conditionnel mixte (complément Bac)]
Un \cle{conditionnel mixte} combine une \textbf{condition passée} (third) avec un \textbf{résultat présent} (second) : la cause est révolue, mais la conséquence se voit encore aujourd'hui.

Forme : \eng{if + past perfect} / \eng{would + base verbale}.

\eng{If I had taken the job, I would live in Bamenda now.} « Si j'avais accepté ce travail, j'habiterais à Bamenda maintenant. »

\eng{If I had worked harder, I would be rich now.} « Si j'avais travaillé plus dur, je serais riche aujourd'hui. »

Au niveau Première, il suffit de savoir le \textbf{reconnaître} et le comprendre.
\end{definition}

\begin{definition}[I wish et If only (complément Bac)]
Pour exprimer un \textbf{souhait} ou un \textbf{regret}, l'anglais utilise \eng{I wish} ou \eng{If only} (« si seulement »), suivis des mêmes reculs de temps que les conditionnels :

\eng{I wish + past simple} : regret sur le présent. \eng{I wish I lived near the school.} « Si seulement j'habitais près de l'école. »

\eng{I wish + past perfect} : regret sur le passé. \eng{If only I had revised my irregular verbs!} « Si seulement j'avais révisé mes verbes irréguliers ! »
\end{definition}

\begin{savaistu}[Les conditionnels au GCE]
Au \cle{GCE Ordinary Level} comme à l'\cle{Advanced Level}, les conditionnels apparaissent chaque année dans la section \emph{Grammar} : exercices de complétion (\eng{If she had come, she \ldots\ the prize}) et de transformation (\eng{I didn't study, so I failed} devient \eng{If I \ldots}). Maîtriser les trois structures et la règle d'or — jamais de \eng{will} ou \eng{would} après \eng{if} — te garantit des points faciles le jour de l'examen.
\end{savaistu}

\begin{aretenir}[L'essentiel de la synthèse]
\begin{itemize}
\item First : \eng{if + present} / \eng{will + BV} — futur possible.
\item Second : \eng{if + past simple} / \eng{would + BV} — présent imaginaire.
\item Third : \eng{if + past perfect} / \eng{would have + p.p.} — passé révolu.
\item Jamais de \eng{will} ni de \eng{would} dans la if-clause.
\item Après \eng{would have}, toujours le participe passé.
\item Virgule seulement quand la if-clause ouvre la phrase.
\end{itemize}
\end{aretenir}

\coursec{Entraînement guidé et production}

Nous avons maintenant toutes les cartes en main~: les trois conditionnels, leurs formes, leurs emplois et les pièges qui guettent les francophones. Il est temps de passer de la théorie à la pratique. Cette dernière partie te propose un parcours d'entraînement progressif, construit comme un escalier~: chaque marche est un peu plus exigeante que la précédente, de la simple identification jusqu'à la production libre sur ta propre vie. C'est exactement la progression que tu retrouveras à l'examen.

\begin{popfigure}
\begin{tikzpicture}[
  stepS/.style={draw=popdark, thick, rounded corners=4pt, align=center, text=white, font=\small\bfseries, minimum height=1cm, minimum width=3.5cm, inner xsep=5pt},
  lab/.style={font=\footnotesize, text=popdark, align=left, anchor=west},
  arrow/.style={very thick, popdark, -{Stealth}}
]
\matrix [row sep=0.4cm, column sep={0.6cm,between origins}, nodes={anchor=west}, ampersand replacement=\&] {
\node[stepS, fill=popblue] (s1) {1. Identifier}; \& \node[lab] {reconnaître le type (1, 2 ou 3)}; \\
\node[stepS, fill=popgreen] (s2) {2. Compléter}; \& \node[lab] {conjuguer le verbe au bon temps}; \\
\node[stepS, fill=popgold] (s3) {3. Transformer}; \& \node[lab] {réécrire avec \emph{If...} (type Bac)}; \\
\node[stepS, fill=poporange] (s4) {4. Traduire}; \& \node[lab] {français $\rightarrow$ anglais sans calquer}; \\
\node[stepS, fill=popink] (s5) {5. Produire}; \& \node[lab] {écrire sur sa vie (projet, rêve, regret)}; \\
};
\draw[arrow] (s1.south west) -- ++(-0.5,0) -| (s5.north east) node[midway, below, font=\footnotesize\itshape, text=popmuted] {difficulté croissante};
\end{tikzpicture}
\legende{L'escalier de difficulté~: cinq marches, de l'identification à la production libre.}
\end{popfigure}

\courssub{Exercice 1~: identifier le type de conditionnel}

Avant de conjuguer quoi que ce soit, il faut savoir \emph{lire} une phrase conditionnelle. La méthode est simple~: regarde le temps de la subordonnée (après \cle{if}), puis celui de la principale, et compare avec le tableau des trois types (rappel du \emph{zero conditional} inclus pour éviter la confusion).

\begin{center}
\begin{tabularx}{\linewidth}{|l|X|X|X|}
\hline
\rowcolor{popblueL} Type & Après \emph{if} & Proposition principale & Sens \\
\hline
Zero & present simple & present simple & vérité générale, fait scientifique \\
\hline
First & present simple & will + base verbale & hypothèse réelle, future \\
\hline
Second & past simple & would + base verbale & irréel du présent, rêve \\
\hline
Third & past perfect & would have + participe passé & regret du passé \\
\hline
\end{tabularx}
\end{center}

\begin{exoresolu}[Exercice 1 — De quel type s'agit-il ?]
Indique si chaque phrase relève du first, second ou third conditional.

\begin{enumerate}
\item \eng{If it rains tomorrow, the match will be cancelled.}
\item \eng{If I were the president, I would build more schools.}
\item \eng{If she had left earlier, she would have caught the bus.}
\item \eng{If you heat water, it boils.}
\item \eng{If we had known about the strike, we would have stayed at home.}
\item \eng{If my father gets a bonus, we will travel to Kribi.}
\end{enumerate}

\tcblower
\textbf{Solution détaillée.}
\begin{enumerate}
\item \eng{rains} (present simple) + \eng{will be cancelled} $\rightarrow$ \textbf{first conditional}. La pluie de demain est une possibilité réelle.
\item \eng{were} (past simple, forme irréelle) + \eng{would build} $\rightarrow$ \textbf{second conditional}. Je ne suis pas président~: c'est un rêve.
\item \eng{had left} (past perfect) + \eng{would have caught} $\rightarrow$ \textbf{third conditional}. Elle n'est pas partie tôt~: regret sur le passé.
\item \eng{heat} + \eng{boils}~: deux present simples $\rightarrow$ c'est le \textbf{zero conditional} (vérité générale), vu en classe de seconde. Piège classique~: ce n'est ni le type 1, ni le 2, ni le 3~!
\item \eng{had known} + \eng{would have stayed} $\rightarrow$ \textbf{third conditional}.
\item \eng{gets} + \eng{will travel} $\rightarrow$ \textbf{first conditional}.
\end{enumerate}
\end{exoresolu}

\courssub{Exercice 2~: compléter avec le verbe au temps correct}

Ici, on te donne le verbe entre parenthèses et tu dois le conjuguer. La clé est la \cle{concordance} des deux propositions~: un temps en appelle un autre, comme deux danseurs qui doivent rester synchronisés.

\begin{attention}[La concordance des deux propositions]
Dans les exercices de complétion, vérifie toujours que les deux moitiés de la phrase «~s'appellent~» l'une l'autre~:
\begin{itemize}
\item \textbf{past perfect} après \eng{if} $\rightarrow$ obligatoirement \eng{would have} + participe passé dans l'autre proposition~: \eng{If I \textbf{had studied}, I \textbf{would have passed}.}
\item \textbf{past simple} après \eng{if} $\rightarrow$ \eng{would} + base verbale~: \eng{If I \textbf{had} money, I \textbf{would buy} a bike.}
\item \textbf{present simple} après \eng{if} $\rightarrow$ \eng{will} + base verbale~: \eng{If it \textbf{rains}, we \textbf{will stay} at home.}
\end{itemize}
Si tu écris \eng{If I had studied, I would pass}, tu mélanges le type 3 et le type 2~: la phrase boite. Garde la même logique du début à la fin.
\end{attention}

\begin{exoresolu}[Exercice 2 — Conjugue les verbes entre parenthèses]
\begin{enumerate}
\item \eng{If you (not hurry), you (miss) the bus to Douala.}
\item \eng{If I (be) rich, I (buy) a house in Buea.}
\item \eng{If they (listen) to the teacher, they (not fail) the test last week.}
\item \eng{If she (pass) her GCE, she (go) to university next year.}
\end{enumerate}

\tcblower
\textbf{Solution détaillée.}
\begin{enumerate}
\item \eng{If you \textbf{don't hurry}, you \textbf{will miss} the bus to Douala.} — First conditional~: present simple après \eng{if}, \eng{will} + base verbale dans la principale. «~Si tu ne te dépêches pas, tu vas rater le bus pour Douala.~»
\item \eng{If I \textbf{were} rich, I \textbf{would buy} a house in Buea.} — Second conditional. Attention~: avec le verbe \emph{be}, on utilise \eng{were} pour toutes les personnes dans le style soigné (\eng{If I were you...}).
\item \eng{If they \textbf{had listened} to the teacher, they \textbf{wouldn't have failed} the test last week.} — Third conditional~: le repère \eng{last week} et le sens (reproche sur le passé) imposent le past perfect, donc \eng{would have} + participe passé de l'autre côté.
\item \eng{If she \textbf{passes} her GCE, she \textbf{will go} to university next year.} — First conditional~: projet réel, daté (\eng{next year}).
\end{enumerate}
\end{exoresolu}

\begin{exercice}[À toi de jouer — complétion]
Conjugue les verbes entre parenthèses.
\begin{enumerate}
\item \eng{If we (win) the match on Saturday, the whole quarter (celebrate).}
\item \eng{If my grandmother (live) in Yaoundé, we (visit) her every week.}
\item \eng{If you (not forget) your identity card, the police officer (not fine) you yesterday.}
\item \eng{If I (know) the answer, I (tell) you, believe me.}
\end{enumerate}
\end{exercice}

\begin{corrige}[Exercice «~À toi de jouer — complétion~»]
\begin{enumerate}
\item \eng{If we \textbf{win} the match on Saturday, the whole quarter \textbf{will celebrate}.} (first)
\item \eng{If my grandmother \textbf{lived} in Yaoundé, we \textbf{would visit} her every week.} (second)
\item \eng{If you \textbf{hadn't forgotten} your identity card, the police officer \textbf{wouldn't have fined} you yesterday.} (third)
\item \eng{If I \textbf{knew} the answer, I \textbf{would tell} you, believe me.} (second)
\end{enumerate}
\end{corrige}

\courssub{Exercice 3~: la transformation, l'épreuve reine du Bac}

L'exercice de transformation est un classique des épreuves du Probatoire et du Baccalauréat~: on te donne une phrase et tu dois la réécrire en commençant par \eng{If...}, sans changer le sens. Voici la méthode infaillible.

\begin{methode}[Réussir les transformations «~If...~» au Bac]
\begin{enumerate}
\item \textbf{Repère le temps de la phrase de départ.} Présent ou futur~? $\rightarrow$ first conditional. Présent irréel~? $\rightarrow$ second. Passé composé ou plus-que-parfait~? $\rightarrow$ third.
\item \textbf{Choisis le conditionnel correspondant} et construis la subordonnée~: \eng{if} + sujet + verbe au bon temps.
\item \textbf{Transforme l'autre partie} en proposition principale~: \eng{will}, \eng{would} ou \eng{would have} + participe passé.
\item \textbf{Vérifie qu'il n'y a ni \eng{will} ni \eng{would} juste après \eng{if}.} C'est l'erreur éliminatoire~: \eng{If you will work...} est faux.
\item \textbf{Relis le sens~:} ta phrase doit dire exactement la même chose que la phrase de départ.
\end{enumerate}
\end{methode}

\begin{exemplebox}[Deux transformations traitées pas à pas]
\textbf{Exemple A.} \eng{Work hard and you will pass.} («~Travaille dur et tu réussiras.~»)
\begin{itemize}
\item Étape 1~: \eng{will pass} annonce un futur réel $\rightarrow$ first conditional.
\item Étape 2~: l'ordre \eng{work hard} devient la condition~: \eng{If you work hard}.
\item Étape 3~: la conséquence garde \eng{will}~: \eng{you will pass}.
\item Résultat~: \eng{\textbf{If you work hard, you will pass.}}
\end{itemize}
\textbf{Exemple B.} \eng{I didn't study, so I failed.} («~Je n'ai pas étudié, donc j'ai échoué.~»)
\begin{itemize}
\item Étape 1~: \eng{didn't study} et \eng{failed} sont au passé $\rightarrow$ third conditional (regret).
\item Étape 2~: la cause devient la condition au past perfect~: \eng{If I had studied}.
\item Étape 3~: la conséquence devient \eng{would have} + participe passé, à la forme négative car l'échec n'aurait pas eu lieu~: \eng{I wouldn't have failed}.
\item Résultat~: \eng{\textbf{If I had studied, I wouldn't have failed.}}
\end{itemize}
\end{exemplebox}

\begin{exercice}[Exercice 3 — Transforme en commençant par «~If...~»]
\begin{enumerate}
\item \eng{Save your money and you will buy a new phone.}
\item \eng{He didn't wear a helmet, so he got injured in the bendskin accident.}
\item \eng{Take an umbrella or you will get wet.}
\item \eng{She didn't listen to her parents, so she made a mistake.}
\end{enumerate}
\end{exercice}

\begin{corrige}[Exercice 3]
\begin{enumerate}
\item \eng{\textbf{If you save your money, you will buy a new phone.}} (first~: futur réel)
\item \eng{\textbf{If he had worn a helmet, he wouldn't have got injured in the bendskin accident.}} (third~: passé, regret~; participe passé de \emph{wear}~: \eng{worn})
\item \eng{\textbf{If you don't take an umbrella, you will get wet.}} (first~; attention, \eng{or} introduit ici une condition \emph{négative})
\item \eng{\textbf{If she had listened to her parents, she wouldn't have made a mistake.}} (third)
\end{enumerate}
\end{corrige}

\courssub{Exercice 4~: traduction du français vers l'anglais}

Traduire, c'est résister à la tentation du mot à mot. En français, on dit «~si j'aurais su~» dans le langage familier~: en anglais, jamais de conditionnel après \eng{if}. Traduis d'abord la \emph{logique} (réel, irréel, regret), puis choisis le type.

\begin{exoresolu}[Exercice 4 — Traduis en anglais]
\begin{enumerate}
\item Si tu étudies régulièrement, tu réussiras au Probatoire.
\item Si j'étais toi, je parlerais au directeur.
\item Si nous avions pris le car de 6~heures, nous serions arrivés à temps à Bamenda.
\item Si j'aurais su, je serais venu plus tôt. (phrase entendue dans la cour de récréation...)
\end{enumerate}

\tcblower
\textbf{Solution détaillée.}
\begin{enumerate}
\item \eng{If you study regularly, you will pass the Probatoire.} — Hypothèse réelle~: first conditional. Ne traduis pas «~réussiras~» par un futur après \eng{if}~!
\item \eng{If I were you, I would speak to the headmaster.} — Expression figée~: \eng{If I were you} (jamais \eng{If I was you} dans le style soigné).
\item \eng{If we had taken the six o'clock coach, we would have arrived in Bamenda on time.} — Regret du passé~: past perfect + \eng{would have} + participe passé.
\item \eng{If I \textbf{had known}, I \textbf{would have come} earlier.} — C'est le piège annoncé~: le français familier «~si j'aurais su~» est déjà fautif en français (il faut «~si j'avais su~»), et sa traduction mot à mot \eng{If I would have known} est une double faute en anglais. Après \eng{if}, c'est le past perfect, point final.
\end{enumerate}
\end{exoresolu}

\courssub{Exercice 5~: corrige les phrases fautives}

Voici des phrases réellement produites par des élèves francophones. Trouve l'erreur, corrige-la et explique la règle.

\begin{exercice}[Exercice 5 — Chasse aux erreurs]
\begin{enumerate}
\item \eng{If I will have money, I will buy a computer.}
\item \eng{If she would come, I would be happy.}
\item \eng{If they had went to school, they would have succeeded.}
\item \eng{If I was you, I will apologize.}
\item \eng{If he had studied, he would passed the exam.}
\end{enumerate}
\end{exercice}

\begin{corrige}[Exercice 5]
\begin{enumerate}
\item \eng{If I \textbf{have} money, I will buy a computer.} — Jamais \eng{will} après \eng{if}~: present simple dans la subordonnée.
\item \eng{If she \textbf{came}, I would be happy.} — Jamais \eng{would} après \eng{if}~: past simple pour l'irréel du présent.
\item \eng{If they had \textbf{gone} to school, they would have succeeded.} — Après \eng{had}, il faut le \emph{participe passé}~: \eng{go -- went -- gone}. Erreur de verbe irrégulier.
\item \eng{If I \textbf{were} you, I \textbf{would} apologize.} — Deux fautes~: \eng{were} (pas \eng{was}) dans l'expression figée, et \eng{would} (pas \eng{will}) pour l'irréel.
\item \eng{If he had studied, he would \textbf{have passed} the exam.} — Après \eng{would}, on ne met jamais un participe passé seul~: il faut \eng{would have} + participe passé.
\end{enumerate}
\end{corrige}

\courssub{Production libre guidée~: parle de ta vie}

Dernière marche de l'escalier~: la production. Tu vas écrire \textbf{trois phrases} sur ta propre existence, une par conditionnel~:

\begin{itemize}
\item \textbf{Un projet réel} (first conditional)~: quelque chose qui peut vraiment arriver cette année. Modèle~: \eng{If I pass my exams, my parents will buy me a bicycle.}
\item \textbf{Un rêve} (second conditional)~: quelque chose d'impossible ou d'improbable aujourd'hui. Modèle~: \eng{If I lived in London, I would watch Arsenal play every weekend.}
\item \textbf{Un regret} (third conditional)~: quelque chose que tu aurais voulu faire différemment. Modèle~: \eng{If I had worked harder in Form Five, I would have had better marks.}
\end{itemize}

\begin{experience}[Atelier d'écriture en binômes]
\begin{enumerate}
\item Écris tes trois phrases au brouillon (5~minutes), en silence.
\item Échange ta feuille avec ton voisin~: chacun corrige la copie de l'autre en utilisant la grille d'auto-évaluation ci-dessous.
\item Lis ta meilleure phrase à la classe~; les autres élèves doivent dire de quel type de conditionnel il s'agit.
\end{enumerate}
\end{experience}

\begin{textebox}[Grille d'auto-évaluation — My conditional checklist]
Before I give my work to the teacher, I check my three sentences one by one.

\textbf{Sentence 1 — my real project (first conditional).} I did not put “will” after “if”. I used the present simple in the if-clause and “will” + verb in the main clause.

\textbf{Sentence 2 — my dream (second conditional).} I did not put “would” after “if”. I used the past simple in the if-clause and “would” + verb in the main clause. With the verb “to be”, I wrote “If I were...”.

\textbf{Sentence 3 — my regret (third conditional).} I used the past perfect after “if” and “would have” + past participle in the main clause. My past participle is correct — I checked my irregular verbs list.

\textbf{General check.} I did not put “will” or “would” after “if”. I used the right tense in each clause. My past participle is correct. My two clauses match~: they belong to the same type of conditional. I read my sentences aloud and they sound natural.
\end{textebox}

\begin{aretenir}[L'essentiel de l'entraînement]
\begin{itemize}
\item Identifier avant de conjuguer~: le temps après \eng{if} révèle le type.
\item Les deux propositions dansent ensemble~: past perfect $\leftrightarrow$ \eng{would have} + participe passé.
\item Transformation type Bac~: temps de départ $\rightarrow$ type de conditionnel $\rightarrow$ vérification finale.
\item Jamais \eng{will} ni \eng{would} juste après \eng{if}~: c'est la faute signature du francophone.
\item En production, une phrase par type~: un projet, un rêve, un regret.
\end{itemize}
\end{aretenir}

Tu as gravé les cinq marches de l'escalier. Les conditionnels n'ont désormais plus de secret pour toi~: à toi de les faire vivre, à l'écrit comme à l'oral, chaque fois que tu parles de tes projets, de tes rêves et de tes regrets.

\newpage

\coursec{Activité d'intégration}

\coursec{Lesson 10 — Grammar: Conditionals (1, 2, 3)}

\courssub{Observation et découverte : Le conseil de famille des Nkeng}

\begin{textebox}[A family meeting in Bafoussam]
The Nkeng family lives in Bafoussam, in the West Region of Cameroon. Last year was difficult: the harvest was poor, school fees were expensive, and the old television broke down. This Sunday evening, after dinner, the whole family sits in the living room for a family meeting. Father wants to plan the future. Mother has advice for everybody. Their son Alain, who is in \eng{Première} (Year 12), dreams of a different life. And Grandmother remembers the past and thinks about the mistakes the family made.

Each person speaks. Each person uses “if”. Listen to the beginning of their meeting:

\textbf{Father:} “If we save two thousand francs every week, we will buy a new television before Christmas.”

\textbf{Mother:} “Alain, if I were you, I would spend less time on WhatsApp and more time on your books.”

\textbf{Alain:} “If we lived in Douala, I would join a football academy. But we are here, so…”

\textbf{Grandmother:} “Ah, my children… If we had planted the maize earlier last year, we would have had a better harvest. But the rains came and the fields were empty.”
\end{textebox}

\begin{savaistu}[The family meeting : une tradition vivante]
Dans de nombreuses familles camerounaises, le \eng{family meeting} ou conseil de famille est un moment important : on y discute de l'argent, de l'école, des travaux des champs et des conflits. C'est un lieu naturel pour les hypothèses : on y parle de projets (\eng{if we save...}), de conseils (\eng{if I were you...}) et de regrets (\eng{if we had known...}). Les conditionnels sont donc la grammaire vivante de la vie familiale !
\end{savaistu}

\begin{exercice}[Tâche 1 — Compréhension et observation (10 min)]
Lis le dialogue de la famille Nkeng.
\begin{enumerate}
\item Souligne les quatre phrases conditionnelles.
\item Indique pour chacune son type : \cle{first}, \cle{second} ou \cle{third conditional}.
\item Recopie dans ton cahier la structure de chaque phrase : \eng{if} + temps, proposition principale + temps.
\end{enumerate}
\end{exercice}

\begin{corrige}[Exercice 1 — Observation]
\begin{itemize}
\item \textbf{Father:} “If we \textbf{save} (present simple), we \textbf{will buy} (future simple)…” $\rightarrow$ \cle{First conditional}.
\item \textbf{Mother:} “If I \textbf{were} (past simple \textit{were}), I \textbf{would spend} (would + BV)…” $\rightarrow$ \cle{Second conditional}.
\item \textbf{Alain:} “If we \textbf{lived} (past simple), we \textbf{would join} (would + BV)…” $\rightarrow$ \cle{Second conditional}.
\item \textbf{Grandmother:} “If we \textbf{had planted} (past perfect), we \textbf{would have had} (would have + p.p.)…” $\rightarrow$ \cle{Third conditional}.
\end{itemize}
\end{corrige}

\courssub{First Conditional : L'hypothèse réelle (Real Possibility)}

\begin{propriete}[Forme et emploi du First Conditional]
\textbf{Structure :} \eng{If + Present Simple, ... will + Base Verbale (BV)} \\
\textbf{Emploi :} Exprime une \cle{hypothèse réelle}, un \cle{projet} ou une \cle{menace} possible dans le futur. La condition est réalisable.
\begin{itemize}
\item La proposition \eng{if} (subordonnée) est au \textbf{present simple} (pas de \eng{will} !).
\item La proposition principale est au \textbf{futur simple} (\eng{will} + BV).
\item L'ordre peut s'inverser : \eng{We will buy a TV if we save money.} (pas de virgule).
\end{itemize}
\end{propriete}

\begin{exemplebox}[Exemples traduits — First Conditional]
\begin{itemize}
\item \eng{If it rains tomorrow, we will stay at home.} « S'il pleut demain, nous resterons à la maison. »
\item \eng{If you work hard, you will pass your exams.} « Si tu travailles dur, tu réussiras tes examens. »
\item \eng{We will miss the bus if we don't leave now.} « Nous raterons le bus si nous ne partons pas maintenant. »
\item \eng{If the cocoa price rises, farmers will earn more.} « Si le prix du cacao augmente, les planteurs gagneront plus. »
\end{itemize}
\end{exemplebox}

\begin{attention}[Piège n°1 : \eng{will} après \eng{if} ? JAMAIS !]
\begin{center}
\begin{tabularx}{\linewidth}{|X|X|}\hline
\rowcolor{poporangeL} \textbf{Incorrect (Francophonie)} & \textbf{Correct (Anglais)} \\ \hline
\eng{If we \textbf{will save} money, we will buy a TV.} & \eng{If we \textbf{save} money, we \textbf{will buy} a TV.} \\ \hline
\eng{If it \textbf{will rain}, I \textbf{will take} an umbrella.} & \eng{If it \textbf{rains}, I \textbf{will take} an umbrella.} \\ \hline
\end{tabularx}
\end{center}
\eng{If} appelle le présent pour exprimer le futur. Règle d'or : \textbf{No \eng{will} after \eng{if}}.
\end{attention}

\courssub{Second Conditional : L'irréel du présent (Unreal Present / Advice)}

\begin{propriete}[Forme et emploi du Second Conditional]
\textbf{Structure :} \eng{If + Past Simple, ... would + Base Verbale (BV)} \\
\textbf{Emploi :} Exprime une \cle{hypothèse imaginaire}, un \cle{rêve}, un \cle{conseil} ou une situation \cle{contraire à la réalité présente}. La condition est peu probable ou impossible \textbf{maintenant}.
\begin{itemize}
\item La proposition \eng{if} est au \textbf{past simple} (souvent \eng{were} pour tous les sujets dans le style formel/conseil).
\item La proposition principale utilise \textbf{\eng{would} + BV}.
\item \eng{Would} = \eng{'d} à l'oral (\eng{I'd, you'd, he'd, we'd, they'd}).
\end{itemize}
\end{propriete}

\begin{exemplebox}[Exemples traduits — Second Conditional]
\begin{itemize}
\item \eng{If I were you, I would apologise to her.} « Si j'étais toi, je m'excuserais auprès d'elle. » (Conseil standard : \eng{were} pour \eng{I}).
\item \eng{If we lived in Yaoundé, we would visit the National Museum.} « Si nous habitions Yaoundé, nous visiterions le Musée national. » (Réalité : nous n'habitons pas Yaoundé).
\item \eng{If Alain had a smartphone, he would use educational apps.} « Si Alain avait un smartphone, il utiliserait des applis éducatives. » (Réalité : il n'en a pas).
\item \eng{Would you help me if I asked you?} « M'aiderais-tu si je te le demandais ? »
\end{itemize}
\end{exemplebox}

\begin{attention}[Piège n°2 : \eng{If I was you} vs \eng{If I were you}]
Dans le conseil (\eng{If I were you...}), la grammaire standard exige \textbf{\eng{were}} pour toutes les personnes (subjonctif anglais survivant).
\begin{itemize}
\item \textbf{Correct (Standard/Examen) :} \eng{If I \textbf{were} you, I would study harder.}
\item \textbf{Toléré (Oral familier) :} \eng{If I \textbf{was} you, I would study harder.}
\item \textbf{Faux (Interférence français) :} \eng{If I \textbf{would be} you...} $\leftarrow$ \textbf{Jamais de \eng{would} dans la proposition \eng{if} !}
\end{itemize}
\end{attention}

\courssub{Third Conditional : Le regret du passé (Unreal Past / Regret)}

\begin{propriete}[Forme et emploi du Third Conditional]
\textbf{Structure :} \eng{If + Past Perfect (had + Participe Passé), ... would have + Participe Passé} \\
\textbf{Emploi :} Exprime un \cle{regret}, une \cle{reproche} ou une \cle{hypothèse impossible} car le passé ne peut pas changer. La condition n'a \textbf{pas été réalisée}.
\begin{itemize}
\item Proposition \eng{if} : \textbf{\eng{had} + Participe Passé} (Past Perfect).
\item Proposition principale : \textbf{\eng{would have} + Participe Passé}.
\item Contractions orales : \eng{would have} $\rightarrow$ \eng{would've} (se prononce « \textit{ou-deuve} ») ; \eng{had} $\rightarrow$ \eng{'d} (\eng{If I'd known...}).
\end{itemize}
\end{propriete}

\begin{exemplebox}[Exemples traduits — Third Conditional]
\begin{itemize}
\item \eng{If we had planted earlier, we would have had a better harvest.} « Si nous avions planté plus tôt, nous aurions eu une meilleure récolte. »
\item \eng{If you had told me, I would have helped you.} « Si tu me l'avais dit, je t'aurais aidé. »
\item \eng{She wouldn't have failed if she had revised.} « Elle n'aurait pas échoué si elle avait révisé. »
\item \eng{If Grandfather had listened to Grandmother, they would have bought the field.} « Si Grand-père avait écouté Grand-mère, ils auraient acheté le champ. »
\end{itemize}
\end{exemplebox}

\begin{attention}[Piège n°3 : Confusion \eng{had} (auxiliaire) et \eng{have} (participe) / \eng{would have}]
\begin{center}
\begin{tabularx}{\linewidth}{|X|X|}\hline
\rowcolor{popredL} \textbf{Erreur fréquente} & \textbf{Correction \& Explication} \\ \hline
\eng{If we \textbf{would have} planted...} & \eng{If we \textbf{had planted}...} (Past Perfect dans le \eng{if}) \\ \hline
\eng{...we \textbf{would had} a better harvest.} & \eng{...we \textbf{would have had} a better harvest.} (\eng{have} = auxiliaire, \eng{had} = p.p. de \eng{have}) \\ \hline
\eng{If I \textbf{would have known}...} & \eng{If I \textbf{had known}... / If I'd known...} \\ \hline
\end{tabularx}
\end{center}
Rappel verbes irréguliers clés : \eng{buy/bought/bought}, \eng{sell/sold/sold}, \eng{pay/paid/paid}, \eng{hear/heard/heard}, \eng{tell/told/told}.
\end{attention}

\courssub{Synthèse, Comparaison avec le français et Pièges}

\begin{aretenir}[Les trois conditionnels — Tableau récapitulatif]
\begin{center}
\begin{tabularx}{\linewidth}{|l|X|X|X|}\hline
\rowcolor{popblueL} \textbf{Type} & \textbf{Structure} & \textbf{Degré de réalité} & \textbf{Emploi principal} \\ \hline
\cle{First} & \eng{If + Present Simple, will + BV} & \textbf{Réel / Probable} (Futur) & Projet, promesse, avertissement, prédiction \\ \hline
\cle{Second} & \eng{If + Past Simple, would + BV} & \textbf{Irréel / Imaginaire} (Présent) & Rêve, conseil (\eng{If I were you}), hypothèse théorique \\ \hline
\cle{Third} & \eng{If + Past Perfect, would have + P.P.} & \textbf{Impossible} (Passé) & Regret, reproche, critique, situation contrefactuelle \\ \hline
\end{tabularx}
\end{center}
\textbf{Règle d'or commune :} \textbf{Jamais de \eng{will} ni de \eng{would} dans la proposition introduite par \eng{if}.}
\end{aretenir}

\begin{center}
\begin{tabularx}{\linewidth}{|X|X|X|}\hline
\rowcolor{poppurpleL} \textbf{Français (Si + Indicatif / Conditionnel)} & \textbf{English (If + ...)} & \textbf{Remarque} \\ \hline
Si je \textbf{vais} à Douala, je \textbf{verrai} mon oncle. & \eng{If I \textbf{go} to Douala, I \textbf{will see} my uncle.} & Présent $\rightarrow$ \textbf{First} \\ \hline
Si j'\textbf{allais} à Douala, je \textbf{verrais} mon oncle. & \eng{If I \textbf{went} to Douala, I \textbf{would see} my uncle.} & Imparfait $\rightarrow$ \textbf{Second} \\ \hline
Si j'\textbf{étais allé} à Douala, j'\textbf{aurais vu} mon oncle. & \eng{If I \textbf{had gone} to Douala, I \textbf{would have seen} my uncle.} & Plus-que-parfait $\rightarrow$ \textbf{Third} \\ \hline
\end{tabularx}
\end{center}

\begin{attention}[Autres pièges à éviter pour les francophones]
\begin{itemize}
\item \textbf{Interrogation indirecte :} \eng{I don't know \textbf{if he will come}.} (Pas un conditionnel ! \eng{will} est correct ici car c'est une question indirecte sur le futur).
\item \textbf{\eng{Unless} = \eng{If not} :} \eng{Unless we save money, we won't buy the TV.} = \eng{If we don't save money...} (First conditional). \eng{Unless} appelle l'affirmatif.
\item \textbf{\eng{In case} (précaution) $\neq$ \eng{If} (condition) :} \eng{Take an umbrella \textbf{in case} it rains.} (Précaution : on le prend avant). \eng{Take an umbrella \textbf{if} it rains.} (Condition : on le prend seulement s'il pleut).
\item \textbf{Traduction de « même si » :} \eng{Even if} + même structure conditionnelle. Ex: \eng{Even if I were rich, I would work.} (Second).
\item \textbf{Prononciation :} \eng{would have} $\rightarrow$ /wʊdəv/ « \textit{ou-deuve} » (\eng{would've}). Ne pas prononcer le \eng{h} de \eng{have}.
\end{itemize}
\end{attention}

\courssub{Entraînement guidé et Production : Le conseil de famille (Jeu de rôle)}

\begin{experience}[Tâche 2 — Préparation du conseil de famille (20 min)]
Par groupes de 4, distribuez les rôles : \eng{Father}, \eng{Mother}, \eng{Child (Son/Daughter)}, \eng{Grandparent}.
Chaque membre prépare \textbf{au moins trois phrases} pour la réunion :
\begin{enumerate}
\item Un \cle{projet réel} pour la famille $\rightarrow$ \textbf{First Conditional} : \eng{If we + present simple, ... will + BV}.
\item Un \cle{rêve} ou un \cle{conseil} $\rightarrow$ \textbf{Second Conditional} : \eng{If I were you, I would...} / \eng{If we + past simple, ... would + BV}.
\item Un \cle{regret} sur le passé $\rightarrow$ \textbf{Third Conditional} : \eng{If we had + P.P., ... would have + P.P.}
\end{enumerate}
\textbf{Thèmes suggérés (contexte camerounais) :} l'argent (économies, \eng{school fees}, \eng{cocoa sales}), l'école (résultats, \eng{smartphone}, applications), les champs (\eng{harvest}, \eng{plant maize}), la maison (\eng{TV}, \eng{generator}, \eng{repairs}), le football (\eng{academy}, \eng{boots}), les voyages (\eng{Yaoundé}, \eng{Douala}, \eng{Kribi}), les disputes (\eng{quarrel}, \eng{forgive}, \eng{apologise}).
\end{experience}

\begin{experience}[Tâche 3 — Jeu de rôle et Écoute active (15 min)]
\begin{enumerate}
\item Chaque groupe joue sa scène devant la classe (2 à 3 minutes).
\item \textbf{Écoute active :} Les autres groupes notent sur une fiche le nombre de conditionnels corrects entendus par type (First / Second / Third).
\item \textbf{Focus prononciation :} Vérifier les contractions : \eng{I'll} /aɪl/, \eng{we'd} /wiːd/ (pour \eng{would}), \eng{would've} /ˈwʊdəv/ (« \textit{ou-deuve} »), \eng{hadn't} /ˈhædənt/.
\end{enumerate}
\end{experience}

\begin{exercice}[Tâche 4 — Production écrite : \eng{Our family's ifs} (10 min)]
Individuellement, rédigez un court paragraphe (6 à 8 lignes) intitulé \eng{Our family's ifs}.
\begin{itemize}
\item Présentez brièvement la famille (noms, lieu, contexte).
\item Rapportez (discours indirect ou citation) \textbf{au moins une phrase de chaque type conditionnel} prononcée lors de votre jeu de rôle.
\item Concluez par une décision prise par la famille (futur simple : \eng{will} / \eng{going to}).
\end{itemize}
\end{exercice}

\begin{exoresolu}[Production modèle : Scène et Paragraphe]
Énoncé : Joue la scène du conseil de famille, puis rédige le paragraphe \eng{Our family's ifs}.
\tcblower
\textbf{Scène modèle (extrait) :}

\eng{\textbf{Father:} If we sell two bags of cocoa this season, we will pay Alain's school fees without borrowing money.}

\eng{\textbf{Mother:} And if you stopped drinking with your friends every evening, we would save even more! Alain, if I were you, I would apologise to your sister.}

\eng{\textbf{Alain:} OK, Mum. If I had a smartphone, I would study with educational applications... but I know we cannot afford it.}

\eng{\textbf{Grandmother:} If your grandfather had listened to me, we would have bought the big field near the river. Today we would be rich!}

\textbf{Paragraphe modèle :}

\eng{Our family's ifs}

\eng{Last Sunday, the Nkeng family held a meeting in their living room in Bafoussam. Everybody spoke about the future, the present and the past. Father made a real project: “If we sell two bags of cocoa, we will pay the school fees.” Mother gave advice: “If I were you, I would apologise to your sister.” Alain expressed a dream: “If I had a smartphone, I would study more.” Finally, Grandmother expressed a regret: “If my husband had listened to me, we would have bought the big field.” The meeting finished with a decision: the family will save money every week.}

« Dimanche dernier, la famille Nkeng a tenu une réunion dans son salon à Bafoussam. Chacun a parlé du futur, du présent et du passé. Le père a fait un projet réel : "Si nous vendons deux sacs de cacao, nous paierons les frais de scolarité." La mère a donné un conseil : "Si j'étais toi, je m'excuserais auprès de ta sœur." Alain a exprimé un rêve : "Si j'avais un smartphone, j'étudierais davantage." Enfin, la grand-mère a exprimé un regret : "Si mon mari m'avait écoutée, nous aurions acheté le grand champ." La réunion s'est terminée par une décision : la famille économisera de l'argent chaque semaine. »

\textbf{Commentaire des choix de langue :}
\begin{itemize}
\item \eng{If we sell..., we will pay...} : First conditional, projet réaliste (vente possible) ; présent après \eng{if}, futur dans la principale.
\item \eng{If I were you, I would apologise...} : Second conditional, conseil classique ; \eng{were} pour toutes les personnes dans cette expression figée.
\item \eng{If I had a smartphone, I would study...} : Second conditional, rêve (Alain n'a pas de smartphone) ; past simple \eng{had}, puis \eng{would} + BV.
\item \eng{If my husband had listened..., we would have bought...} : Third conditional, regret irréversible ; past perfect (\eng{had listened}) puis \eng{would have} + participe passé (\eng{bought}, irrégulier de \eng{buy}).
\item \eng{Today we would be rich!} (dans la scène) : \textbf{Mixed Conditional} (3rd condition $\rightarrow$ 2nd result) : passé non réalisé $\rightarrow$ conséquence présente imaginaire. Niveau excellence.
\item \eng{The family will save money} : Futur simple hors structure conditionnelle, pour annoncer la décision finale.
\end{itemize}
\end{exoresolu}

\courssub{Bilan de la leçon}

Cette leçon a permis de mobiliser les trois conditionnels dans une situation de communication authentique et ancrée dans la réalité camerounaise : le conseil de famille. Les compétences travaillées sont multiples :
\begin{itemize}
\item \textbf{Reading / Listening :} Compréhension du dialogue support (texte / écoute simulée).
\item \textbf{Grammar :} Identification, analyse, construction des trois structures conditionnelles.
\item \textbf{Speaking :} Jeu de rôle (interaction, prononciation des contractions \eng{'ll, 'd, would've}).
\item \textbf{Writing :} Rédaction d'un paragraphe récapitulatif structuré (\eng{Our family's ifs}).
\item \textbf{Vocabulary :} Lexique de la vie familiale, sociale, agricole et scolaire.
\end{itemize}

\begin{aretenir}[Ce que je dois retenir pour le Bac et la vie]
Le choix du conditionnel dépend du \cle{degré de réalité} de l'hypothèse :
\begin{itemize}
\item \textbf{Possible / Projet} $\rightarrow$ \cle{First Conditional} (\eng{If + Present, Will + BV}).
\item \textbf{Imaginaire / Conseil / Rêve (Présent)} $\rightarrow$ \cle{Second Conditional} (\eng{If + Past, Would + BV}).
\item \textbf{Impossible / Regret (Passé)} $\rightarrow$ \cle{Third Conditional} (\eng{If + Had + P.P., Would Have + P.P.}).
\end{itemize}
\textbf{Check-list avant de rendre ma copie :}
\begin{enumerate}
\item Pas de \eng{will} / \eng{would} après \eng{if}.
\item \eng{If I were you} (pas \eng{was} à l'écrit, pas \eng{would be}).
\item Third conditional : \eng{had} (auxiliaire) + \eng{P.P.} / \eng{would have} + \eng{P.P.} (deux participes passés !).
\item Verbes irréguliers vérifiés (\eng{buy/bought}, \eng{sell/sold}, \eng{pay/paid}, \eng{go/went/gone}).
\end{enumerate}
\end{aretenir}

\begin{exercice}[Devoir / Approfondissement : \eng{My own ifs}]
Rédigez chez vous un second paragraphe intitulé \eng{My own ifs} (8 à 10 lignes) sur \textbf{votre propre vie} :
\begin{enumerate}
\item Une phrase \textbf{First Conditional} : un projet réel pour cette année scolaire (ex: \eng{If I pass my Probatoire, I will...}).
\item Une phrase \textbf{Second Conditional} : un rêve pour votre vie future (ex: \eng{If I were a doctor, I would...}).
\item Une phrase \textbf{Third Conditional} : un regret sur l'année dernière (ex: \eng{If I had worked harder last year, I would have...}).
\end{enumerate}
Vous le présenterez à votre voisin au début du prochain cours (Speaking : \eng{Peer correction}).
\end{exercice}

\newpage

\coursec{Méthodes et exercices résolus}

\coursec{Lesson 10 — Grammar: Conditionals (1, 2, 3)}

\courssub{Observation et découverte}

\begin{textebox}[Dialogue : « Family Dreams and Regrets » — Contexte camerounais]
\eng{Setting: A Sunday afternoon at the Mvondo family house in Yaoundé. Father (Pierre), Mother (Marie), and their son (Jean, 17) are chatting after lunch.}

\eng{Father:} “Jean, \textbf{if you pass} your Probatoire this year, \textbf{we will organise} a big celebration. Your uncle in Douala \textbf{has promised} to pay for your driving lessons.”

\eng{Jean:} “Thanks, Dad. \textbf{If I were} you, \textbf{I would invest} that money in my future business instead. But \textbf{if I had} the choice, \textbf{I would go} to university in Buea.”

\eng{Mother:} “\textbf{If your father had listened} to me ten years ago, \textbf{we would have bought} that plot of land in Bastos. Look at the prices now! \textbf{Unless we save} more, \textbf{we won't afford} it.”

\eng{Father:} “No regrets, Marie. \textbf{If we had not spent} that money on Jean's health, \textbf{he would not be} here with us today. That is what matters.”
\end{textebox}

\begin{experience}[Observation guidée — Travail en binôme]
Lisez le dialogue et répondez aux questions ci-dessous. Repérez les phrases en \textbf{gras}.
\begin{enumerate}
\item Combien de types de phrases conditionnelles différentes trouvez-vous ? Classez-les selon le temps de la situation (futur probable, présent imaginaire, passé révolu).
\item Dans la phrase du père \eng{“If you pass..., we will organise...”}, quels temps verbaux sont utilisés dans la proposition introduite par \eng{if} et dans la principale ?
\item Dans la phrase de Jean \eng{“If I were you, I would invest...”}, le verbe \eng{were} est-il un passé réel ? Que signifie \eng{would} ici ?
\item Dans la phrase de la mère \eng{“If your father had listened..., we would have bought...”}, identifiez la forme verbale après \eng{if} et après \eng{would}. Quel sentiment exprime cette phrase ?
\item Trouvez un exemple avec \eng{unless}. Par quoi peut-on le remplacer ?
\end{enumerate}
\end{experience}

\begin{aretenir}[Bilan de l'observation]
Le dialogue présente les \textbf{trois conditionnels standards} :
\begin{itemize}
\item \textbf{Type 1 (Réel / Futur probable) :} \eng{If you pass, we will organise.} $\rightarrow$ Condition possible, conséquence future certaine.
\item \textbf{Type 2 (Irréel / Présent imaginaire) :} \eng{If I were you, I would invest.} $\rightarrow$ Situation imaginaire \textbf{maintenant}, contraire à la réalité (je ne suis pas toi).
\item \textbf{Type 3 (Irréel / Passé révolu) :} \eng{If he had listened, we would have bought.} $\rightarrow$ Regret ou reproche sur le \textbf{passé} ; la condition ne s'est pas réalisée.
\end{itemize}
\cle{Unless} = \eng{if... not} (négation de la condition). Ex. : \eng{Unless we save} = \eng{If we do not save}.
\end{aretenir}

\courssub{First conditional : l'hypothèse réelle}

\begin{propriete}[Structure et emploi du First Conditional]
\begin{center}
\begin{tabularx}{\linewidth}{|X|X|}
\hline
\rowcolor{popblueL} \textbf{If-clause (Condition)} & \textbf{Main clause (Résultat)} \\
\hline
\eng{If} + \textbf{Present Simple} & \textbf{Will / Can / May / Must / Shall / Imperative} + \textbf{Base Verbale} \\
\hline
\eng{If it rains,} & \eng{we will stay at home.} \\
\eng{If you study hard,} & \eng{you can pass the Bac.} \\
\eng{If you see Paul,} & \eng{tell him to call me.} \\
\hline
\end{tabularx}
\end{center}

\textbf{Emploi :} Hypothèse \textbf{réelle, possible, probable} dans un futur proche ou lointain. La réalisation de la condition entraîne logiquement le résultat.

\textbf{Règles clés :}
\begin{enumerate}
\item \textbf{Jamais de \eng{will} / \eng{would} après \eng{if}.} Le présent simple suffit à marquer le futur dans la subordonnée.
\item \textbf{Accord 3\textsuperscript{e} personne du singulier :} \eng{If he comes...}, \eng{If she has...}, \eng{If the government builds...}
\item \textbf{Ordre des propositions :} 
\begin{itemize}
\item \eng{If} en début de phrase $\rightarrow$ \textbf{virgule} obligatoire. \eng{If it rains, we will stay.}
\item \eng{If} au milieu $\rightarrow$ \textbf{pas de virgule}. \eng{We will stay if it rains.}
\end{itemize}
\item \textbf{Modaux alternatifs :} \eng{can} (possibilité/permission), \eng{may} (probabilité), \eng{must} (obligation), \eng{imperative} (ordre/conseil).
\end{enumerate}
\end{propriete}

\begin{methode}[Construire le First Conditional pas à pas]
\begin{enumerate}
\item \textbf{Analyser le contexte :} La condition est-elle envisageable, réaliste, planifiée ? (Ex. : météo, emploi du temps, promesse, menace, conseil).
\item \textbf{Écrire la subordonnée :} \eng{If} + Sujet + \textbf{Present Simple} (forme affirmative, négative \eng{do/does not}, ou interrogative).
\item \textbf{Écrire la principale :} Sujet + \textbf{Will + Base Verbale} (ou \eng{can/may/must/imperative} + BV).
\item \textbf{Contrôler :} Virgule si \eng{If} en tête ; pas de \eng{will} après \eng{if} ; \textbf{-s} à la 3\textsuperscript{e} personne du singulier au présent.
\end{enumerate}
\end{methode}

\begin{exemplebox}[Exemples variés — Contexte camerounais]
\begin{itemize}
\item \eng{If the rain stops, the match will start on time.} « Si la pluie s'arrête, le match commencera à l'heure. »
\item \eng{If you don't revise your lessons, you will fail the exam.} « Si tu ne révises pas tes leçons, tu échoueras à l'examen. »
\item \eng{If you go to Marché Central, buy some plantains.} « Si tu vas au Marché Central, achète des plantains. » (Impératif)
\item \eng{My father will be angry if I come home late.} « Mon père sera fâché si je rentre tard. » (Pas de virgule)
\item \eng{Unless you hurry, you will miss the bendskin.} « À moins que tu ne te dépêches, tu rateras le bendskin. » (\eng{Unless} = \eng{If you don't})
\end{itemize}
\end{exemplebox}

\begin{exoresolu}[Compléter avec le First Conditional]
\textbf{Énoncé.} Complete the sentences with the correct form of the verbs in brackets (First Conditional).
\begin{enumerate}
\item If the government \ldots\ldots\ldots (build) more classrooms, students \ldots\ldots\ldots (learn) in better conditions.
\item My sister \ldots\ldots\ldots (be) very happy if she \ldots\ldots\ldots (get) a scholarship to study in Buea.
\item If you \ldots\ldots\ldots (not / hurry), you \ldots\ldots\ldots (miss) the bus to Douala.
\item \ldots\ldots\ldots (call) me if you \ldots\ldots\ldots (need) any help with your homework.
\item Unless the cocoa prices \ldots\ldots\ldots (rise), the farmers \ldots\ldots\ldots (not / earn) enough money.
\end{enumerate}
\tcblower
\textbf{Solution détaillée.}
\begin{enumerate}
\item Contexte : projet réel. Subordonnée : \eng{builds} (3\textsuperscript{e} pers. sg). Principale : \eng{will learn}. \eng{If the government \textbf{builds} more classrooms, students \textbf{will learn} in better conditions.} « Si le gouvernement construit plus de salles de classe, les élèves apprendront dans de meilleures conditions. »
\item Principale d'abord : \eng{will be} (base verbale \eng{be} après \eng{will}). Subordonnée : \eng{gets} (3\textsuperscript{e} pers. sg). \eng{My sister \textbf{will be} very happy if she \textbf{gets} a scholarship to study in Buea.} « Ma sœur sera très heureuse si elle obtient une bourse pour étudier à Buea. »
\item Subordonnée négative : \eng{do not hurry} (sujet \eng{you}). Principale : \eng{will miss}. \eng{If you \textbf{do not hurry}, you \textbf{will miss} the bus to Douala.} « Si tu ne te dépêches pas, tu rateras le car pour Douala. »
\item Impératif en principale (conseil) + Type 1 en subordonnée. \eng{\textbf{Call} me if you \textbf{need} any help with your homework.} « Appelle-moi si tu as besoin d'aide pour tes devoirs. »
\item \eng{Unless} = \eng{If ... not}. Subordonnée : \eng{rise} (pluriel \eng{prices}). Principale négative : \eng{will not earn} (\eng{won't earn}). \eng{Unless the cocoa prices \textbf{rise}, the farmers \textbf{won't earn} enough money.} « À moins que les prix du cacao n'augmentent, les agriculteurs ne gagneront pas assez d'argent. »
\end{enumerate}
\textbf{Conclusion.} Dans tous les cas, la condition est jugée possible : on applique \eng{if + present simple} et \eng{will + base verbale} (ou équivalent), sans jamais mettre \eng{will} dans la proposition introduite par \eng{if}.
\end{exoresolu}

\courssub{Second conditional : l'irréel du présent}

\begin{propriete}[Structure et emploi du Second Conditional]
\begin{center}
\begin{tabularx}{\linewidth}{|X|X|}
\hline
\rowcolor{poporangeL} \textbf{If-clause (Condition)} & \textbf{Main clause (Résultat)} \\
\hline
\eng{If} + \textbf{Past Simple} & \textbf{Would / Could / Might} + \textbf{Base Verbale} \\
\hline
\eng{If I had more money,} & \eng{I would buy a car.} \\
\eng{If she were taller,} & \eng{she could play basketball.} \\
\eng{If they knew the truth,} & \eng{they might be angry.} \\
\hline
\end{tabularx}
\end{center}

\textbf{Emploi :} Hypothèse \textbf{irréelle, imaginaire, contraire à la réalité présente} ou peu probable. Le \textbf{Past Simple} ne marque pas le passé, mais la \textbf{distance par rapport à la réalité} (distanciation).

\textbf{Règles clés :}
\begin{enumerate}
\item \textbf{Verbe \eng{be} :} On utilise \textbf{\eng{were}} à \textbf{toutes les personnes} (formel et standard au Bac). \eng{If I were you...}, \eng{If he were here...}. (\eng{was} possible à l'oral familier pour I/he/she/it, mais à éviter à l'écrit).
\item \textbf{Jamais de \eng{would} / \eng{would have} après \eng{if}.}
\item \textbf{Modaux alternatifs :} \eng{could} (capacité/possibilité), \eng{might} (possibilité plus faible).
\item \textbf{Expression figée :} \eng{If I were you, I would...} (Donner un conseil).
\end{enumerate}
\end{propriete}

\begin{methode}[Construire le Second Conditional pas à pas]
\begin{enumerate}
\item \textbf{Repérer l'irréel :} La situation est-elle un rêve, un souhait, un conseil, ou contraire aux faits actuels ? (Ex. : \eng{If I were rich...} — Je ne suis pas riche).
\item \textbf{Construire la subordonnée :} \eng{If} + Sujet + \textbf{Past Simple}. Verbes irréguliers : forme prétérit (colonne 2). \eng{be} $\rightarrow$ \eng{were}.
\item \textbf{Construire la principale :} Sujet + \textbf{Would + Base Verbale} (ou \eng{could/might} + BV).
\item \textbf{Vérifier :} Pas de \eng{would} dans le \eng{if}-clause. Accord du passé simple (pas de \textbf{-s} à la 3\textsuperscript{e} personne).
\end{enumerate}
\end{methode}

\begin{exemplebox}[Exemples variés — Rêves et conseils]
\begin{itemize}
\item \eng{If Cameroon won the World Cup, the whole country would celebrate for weeks.} « Si le Cameroun gagnait la Coupe du monde, tout le pays ferait la fête pendant des semaines. » (Peu probable)
\item \eng{If I were the Minister of Education, I would reduce class sizes.} « Si j'étais ministre de l'Éducation, je réduirais les effectifs des classes. » (Imaginaire / Conseil)
\item \eng{If my parents had more time, they would travel with us.} « Si mes parents avaient plus de temps, ils voyageraient avec nous. » (Contraire à la réalité présente)
\item \eng{If you spoke English fluently, you could work in a call centre in Douala.} « Si tu parlais anglais couramment, tu pourrais travailler dans un centre d'appels à Douala. »
\end{itemize}
\end{exemplebox}

\begin{exoresolu}[Traduire et compléter avec le Second Conditional]
\textbf{Énoncé.}
\begin{enumerate}
\item Complete : If Cameroon \ldots\ldots\ldots (win) the World Cup, the whole country \ldots\ldots\ldots (celebrate) for weeks.
\item Complete : If I \ldots\ldots\ldots (be) the Minister of Education, I \ldots\ldots\ldots (reduce) class sizes.
\item Translate into English : « Si mes parents avaient plus de temps, ils voyageraient avec nous. »
\item Complete : If we \ldots\ldots\ldots (live) in a bigger house, my grandmother \ldots\ldots\ldots (stay) with us.
\end{enumerate}
\tcblower
\textbf{Solution détaillée.}
\begin{enumerate}
\item Hypothèse imaginaire (peu probable) $\rightarrow$ Type 2. \eng{win} (irrégulier : win -- won -- won) $\rightarrow$ \eng{won}. \eng{would celebrate}. \eng{If Cameroon \textbf{won} the World Cup, the whole country \textbf{would celebrate} for weeks.}
\item Conseil / Imaginaire $\rightarrow$ Type 2. \eng{be} $\rightarrow$ \eng{were} (toutes personnes). \eng{would reduce}. \eng{If I \textbf{were} the Minister of Education, I \textbf{would reduce} class sizes.}
\item Français : Imparfait (\eng{avaient}) + Conditionnel (\eng{voyageraient}) = Irréel présent $\rightarrow$ Type 2. \eng{have} $\rightarrow$ \eng{had}. \eng{would travel}. \eng{If my parents \textbf{had} more time, they \textbf{would travel} with us.}
\item \eng{live} $\rightarrow$ \eng{lived} (régulier). \eng{would stay}. \eng{If we \textbf{lived} in a bigger house, my grandmother \textbf{would stay} with us.} « Si nous habitions une plus grande maison, notre grand-mère resterait avec nous. »
\end{enumerate}
\textbf{Conclusion.} Le signal de l'irréel est l'écart avec la réalité : on recule d'un temps (Present $\rightarrow$ Past Simple) dans la subordonnée et on utilise \eng{would/could/might} + BV dans la principale. \textbf{Ne jamais écrire \eng{If I would be...}} (faute grave).
\end{exoresolu}

\courssub{Third conditional : le regret du passé}

\begin{propriete}[Structure et emploi du Third Conditional]
\begin{center}
\begin{tabularx}{\linewidth}{|X|X|}
\hline
\rowcolor{popgreenL} \textbf{If-clause (Condition)} & \textbf{Main clause (Résultat)} \\
\hline
\eng{If} + \textbf{Past Perfect} (\eng{Had} + \textbf{Participe Passé}) & \textbf{Would / Could / Might} + \textbf{Have} + \textbf{Participe Passé} \\
\hline
\eng{If you had revised,} & \eng{you would have passed.} \\
\eng{If they had left earlier,} & \eng{they could have caught the train.} \\
\eng{If it hadn't rained,} & \eng{we might have played.} \\
\hline
\end{tabularx}
\end{center}

\textbf{Emploi :} Hypothèse \textbf{irréelle dans le passé}. L'événement est \textbf{terminé}, on ne peut plus le changer. Exprime le \textbf{regret}, le \textbf{reproche}, la \textbf{critique} ou la \textbf{spéculation} sur un passé différent.

\textbf{Règles clés :}
\begin{enumerate}
\item \textbf{Subordonnée :} \eng{If} + Sujet + \textbf{Had} + \textbf{Participe Passé (V3)}. Négatif : \eng{Had not / Hadn't} + V3.
\item \textbf{Principale :} Sujet + \textbf{Would Have} + \textbf{Participe Passé (V3)}. Négatif : \eng{Would not have / Wouldn't have} + V3.
\item \textbf{Jamais de \eng{would have} après \eng{if}.}
\item \textbf{Prononciation :} \eng{would have} $\rightarrow$ \eng{would've} /wʊdəv/ (noté « oude-av » en transcription française). \textbf{Interdiction absolue d'écrire \eng{would of}.}
\item \textbf{Verbes irréguliers :} Maîtriser la colonne 3 (Participe Passé) : \eng{go -- gone}, \eng{do -- done}, \eng{see -- seen}, \eng{write -- written}, \eng{wake -- woken}, \eng{be -- been}.
\end{enumerate}
\end{propriete}

\begin{methode}[Construire le Third Conditional pas à pas]
\begin{enumerate}
\item \textbf{Repérer l'irréel du passé :} L'événement est passé, la condition ne s'est pas réalisée. (Ex. : \eng{If you had come...} — Tu n'es pas venu).
\item \textbf{Inverser la réalité :} Si la réalité est négative (\eng{didn't study}), la condition est affirmative (\eng{had studied}). Si la réalité est affirmative (\eng{woke up late}), la condition est négative (\eng{had not woken up}).
\item \textbf{Construire la subordonnée :} \eng{If} + Sujet + \textbf{Had + V3}.
\item \textbf{Construire la principale :} Sujet + \textbf{Would Have + V3} (ou \eng{Could Have / Might Have} + V3).
\item \textbf{Vérifier les participes passés} (réguliers \textbf{-ed} / irréguliers colonne 3).
\end{enumerate}
\end{methode}

\begin{exemplebox}[Exemples variés — Regrets et reproches]
\begin{itemize}
\item \eng{If Ngono had studied for the test, she would have passed.} « Si Ngono avait révisé pour le contrôle, elle aurait réussi. » (Regret)
\item \eng{If we had not woken up late, we would not have missed the bus.} « Si nous ne nous étions pas réveillés tard, nous n'aurions pas raté le car. » (Reproche / Analyse)
\item \eng{If the farmers had planted early, the harvest would have been good.} « Si les agriculteurs avaient planté tôt, la récolte aurait été bonne. » (\eng{be} $\rightarrow$ \eng{been})
\item \eng{If you had told me the truth, I could have helped you.} « Si tu m'avais dit la vérité, j'aurais pu t'aider. » (\eng{Could have} = capacité passée non réalisée)
\end{itemize}
\end{exemplebox}

\begin{exoresolu}[Former le Third Conditional à partir d'une situation réelle]
\textbf{Énoncé.} Rewrite each situation as a Third Conditional sentence.
\begin{enumerate}
\item Ngono did not study for the test, so she failed. (If Ngono \ldots)
\item We missed the bus because we woke up late. (If we \ldots)
\item The farmers did not plant early, so the harvest was poor. (If the farmers \ldots)
\item He didn't apply for the job, so he didn't get it. (If he \ldots)
\end{enumerate}
\tcblower
\textbf{Solution détaillée.}
\begin{enumerate}
\item Réalité : \eng{did not study} (négatif) $\rightarrow$ \eng{failed}. Condition contraire : affirmative. \eng{had studied} (régulier). Résultat contraire : \eng{would have passed}. \eng{If Ngono \textbf{had studied} for the test, she \textbf{would have passed}.}
\item Réalité : \eng{woke up late} (affirmatif, irrégulier \eng{wake -- woke -- woken}) $\rightarrow$ \eng{missed}. Condition contraire : négative. \eng{had not woken up} (ou \eng{hadn't woken up}). Résultat contraire : \eng{would not have missed} (\eng{wouldn't have missed}). \eng{If we \textbf{hadn't woken up} late, we \textbf{wouldn't have missed} the bus.}
\item Réalité : \eng{did not plant} (négatif) $\rightarrow$ \eng{was poor}. Condition contraire : \eng{had planted} (régulier). Résultat contraire : \eng{would have been good} (\eng{be} $\rightarrow$ \eng{been}). \eng{If the farmers \textbf{had planted} early, the harvest \textbf{would have been} good.}
\item Réalité : \eng{didn't apply} $\rightarrow$ \eng{didn't get}. Condition : \eng{had applied}. Résultat : \eng{would have got} (UK) / \eng{gotten} (US). Privilégier \eng{got} (UK). \eng{If he \textbf{had applied} for the job, he \textbf{would have got} it.}
\end{enumerate}
\textbf{Conclusion.} Pour transformer une cause réelle passée en regret : on inverse la valeur (affirmatif $\leftrightarrow$ négatif) et on applique \eng{if + had + V3} / \eng{would have + V3}. \textbf{Vérifiez toujours le participe passé des verbes irréguliers.}
\end{exoresolu}

\courssub{Synthèse, comparaison avec le français et pièges}

\begin{aretenir}[Tableau récapitulatif des trois conditionnels]
\begin{center}
\begin{tabularx}{\linewidth}{|X|X|X|X|}
\hline
\rowcolor{popblueL} \textbf{Type} & \textbf{Sens / Temps de la situation} & \textbf{If-clause} & \textbf{Main clause} \\
\hline
\textbf{1 — Réel} & Futur probable / Possible & \eng{If + Present Simple} & \eng{Will / Can / May / Imperative + BV} \\
\hline
\textbf{2 — Irréel Présent} & Présent imaginaire / Rêve / Conseil / Peu probable & \eng{If + Past Simple} (\eng{were} $\forall$ pers.) & \eng{Would / Could / Might + BV} \\
\hline
\textbf{3 — Irréel Passé} & Passé révolu / Regret / Reproche / Contre-factuel & \eng{If + Past Perfect} (\eng{Had + V3}) & \eng{Would / Could / Might + Have + V3} \\
\hline
\end{tabularx}
\end{center}

\textbf{Repères temporels (Mots-clés) :}
\begin{itemize}
\item Type 1 : \eng{tomorrow, next week, if, when, as soon as, unless, in case}.
\item Type 2 : \eng{now, today, if I were you, imagine, suppose, what if}.
\item Type 3 : \eng{yesterday, last year, ago, in 2020, unfortunately, regret}.
\end{itemize}
\end{aretenir}

\begin{attention}[Les 5 erreurs qui coûtent des points au Bac]
\begin{enumerate}
\item \textbf{\eng{Will} ou \eng{Would} après \eng{If}.} 
\eng{If I \eng{will have} money...} $\rightarrow$ \textbf{INCORRECT}. 
\eng{If I \textbf{have} money, I \textbf{will buy} a phone.} (Type 1) \\
\eng{If I \eng{would have} money...} $\rightarrow$ \textbf{INCORRECT}. 
\eng{If I \textbf{had} money, I \textbf{would buy} a phone.} (Type 2)
\item \textbf{Calque du français « si j'aurais » / « si j'aurais su ».} 
En français familier : « Si j'aurais su, j'aurais venu. » 
En anglais : \textbf{JAMAIS} de \eng{would} ou \eng{would have} après \eng{If}. \\
\eng{If I \eng{would have known}, I \eng{would have come}.} $\rightarrow$ \textbf{INCORRECT}. \\
\eng{If I \textbf{had known}, I \textbf{would have come}.} (Type 3)
\item \textbf{Oublier le \textbf{-s} de la 3\textsuperscript{e} personne au Type 1.} 
\eng{If she \eng{come}, we will start.} $\rightarrow$ \textbf{INCORRECT}. 
\eng{If she \textbf{comes}, we will start.} Le Present Simple garde toutes ses règles après \eng{If}.
\item \textbf{Mélanger les types (Hybrides interdits en 1\textsuperscript{re}).} 
\eng{If I \eng{had} money, I \eng{will} buy...} $\rightarrow$ \textbf{INCORRECT}. \\
Soit Type 1 : \eng{If I \textbf{have} money, I \textbf{will buy}...} \\
Soit Type 2 : \eng{If I \textbf{had} money, I \textbf{would buy}...} \\
\textit{Note : Les conditionnels mixtes (Type 2/3) ne sont pas au programme de 1\textsuperscript{re} A.}
\item \textbf{\eng{Would of} au lieu de \eng{Would have}.} 
À l'oral \eng{would've} se prononce « oude-av ». 
\eng{I \eng{would of} done it.} $\rightarrow$ \textbf{FAUTE D'ORTHOGRAPHE GRAVE}. 
\eng{I \textbf{would have done} it.} / \eng{I \textbf{would've done} it.}
\end{enumerate}
\end{attention}

\begin{savaistu}[Le conditionnel en anglais camerounais (CamE) vs Standard British English]
Au Cameroun anglophone (Northwest/Southwest), on entend souvent des structures influencées par le pidgin ou le français :
\begin{itemize}
\item \eng{« If I would have known... »} (Calque du français / Pidgin \eng{If I for know...}) $\rightarrow$ Standard : \eng{If I had known...}
\item \eng{« Had it been I knew... »} (Traduction littérale de « Si j'avais su... ») $\rightarrow$ Standard : \eng{If I had known...}
\item \eng{« If I was you... »} (Au lieu de \eng{were}) $\rightarrow$ Accepté à l'oral familier, mais \textbf{\eng{were} exigé à l'écrit et au Bac}.
\end{itemize}
L'anglais standard (britannique) reste la norme pour les examens officiels (Probatoire, Baccalauréat).
\end{savaistu}

\courssub{Entraînement guidé et production}

\begin{methode}[Identifier le type de conditionnel à partir du contexte]
Face à une phrase à trous ou à une production écrite, suivez cette démarche :
\begin{enumerate}
\item \textbf{Lisez toute la phrase} et posez la question clé : La condition est-elle (a) possible/probable (futur), (b) imaginaire/contre-réalité (présent), (c) contraire à un fait passé (regret) ?
\item \textbf{Repérez les indices temporels :} \eng{next year, tomorrow} $\rightarrow$ Type 1 ; \eng{now, if I were you} $\rightarrow$ Type 2 ; \eng{yesterday, last year, unfortunately} $\rightarrow$ Type 3.
\item \textbf{Appliquez la grille de correspondance} (voir Tableau récapitulatif ci-dessus).
\item \textbf{Vérifiez la cohérence interne :} Les deux moitiés doivent appartenir au même type (sauf conditionnels mixtes, hors programme).
\end{enumerate}
\end{methode}

\begin{exoresolu}[Choix du conditionnel dans un mini-texte]
\textbf{Énoncé.} Complete the text with the correct conditional form of the verbs in brackets.

\emph{Every family has dreams. If my uncle's cocoa farm \textbf{(1)} (produce) more next year, he \textbf{(2)} (send) his son to university. My cousin says: “If I \textbf{(3)} (be) richer, I \textbf{(4)} (buy) my father a new house.” Last year, the family had problems: if the rains \textbf{(5)} (not / destroy) the crops, they \textbf{(6)} (earn) much more money.}
\tcblower
\textbf{Solution détaillée.}
\begin{itemize}
\item \textbf{(1) et (2)} : Projet réel pour l'année prochaine (\eng{next year}) $\rightarrow$ \textbf{Type 1}. Subordonnée : Present Simple, 3\textsuperscript{e} pers. $\rightarrow$ \eng{produces}. Principale : \eng{will send}. \eng{If my uncle's cocoa farm \textbf{produces} more next year, he \textbf{will send} his son to university.}
\item \textbf{(3) et (4)} : Un rêve, contraire à la réalité présente (il n'est pas riche) $\rightarrow$ \textbf{Type 2}. \eng{be} $\rightarrow$ \eng{were}. Principale : \eng{would buy}. \eng{“If I \textbf{were} richer, I \textbf{would buy} my father a new house.”}
\item \textbf{(5) et (6)} : Événement passé et terminé (\eng{last year}), regret $\rightarrow$ \textbf{Type 3}. Subordonnée négative au Past Perfect : \eng{had not destroyed} (\eng{hadn't destroyed}). Principale : \eng{would have earned} (V3 de \eng{earn} : earned). \eng{If the rains \textbf{had not destroyed} the crops, they \textbf{would have earned} much more money.}
\end{itemize}
\textbf{Conclusion.} Le même texte peut mélanger les trois types : c'est le \textbf{sens} (réel futur / rêve présent / regret passé) qui commande le choix, jamais la forme du verbe français seule.
\end{exoresolu}

\begin{methode}[Rédiger un paragraphe avec les trois conditionnels]
Au Bac, on te demandera souvent d'écrire un paragraphe ou une lettre contenant des hypothèses. Démarche :
\begin{enumerate}
\item \textbf{Planifie} : Une phrase de chaque type suffit pour montrer ta maîtrise (une intention réelle, un rêve, un regret).
\item \textbf{Varie les connecteurs} : \eng{if}, mais aussi \eng{unless} (= if... not), \eng{as soon as}, \eng{when} (Type 1 uniquement), \eng{provided that}.
\item \textbf{Soigne les marqueurs logiques} : \eng{first}, \eng{then}, \eng{finally} pour organiser ; \eng{unfortunately}, \eng{regrettably} pour introduire un regret (Type 3).
\item \textbf{Relis} en vérifiant 3 points critiques : 
\begin{itemize}
\item Pas de \eng{will} / \eng{would} / \eng{would have} après \eng{if}.
\item Participes passés corrects (V3) pour Type 3.
\item Virgule si la phrase commence par \eng{If}.
\end{itemize}
\end{enumerate}
\textbf{Structures utiles à réutiliser (Patterns) :}
\begin{itemize}
\item \eng{If I pass my exams, I will...} (Projet réel)
\item \eng{If I had the opportunity, I would...} (Rêve)
\item \eng{If I had worked harder last term, I would have...} (Regret)
\item \eng{Unless you hurry, you will be late.} « À moins que tu ne te dépêches, tu seras en retard. »
\end{itemize}
\end{methode}

\begin{exoresolu}[Production guidée : « My plans and my regrets »]
\textbf{Énoncé.} Write a short paragraph (60--80 words) about your family life, using at least one First Conditional, one Second Conditional and one Third Conditional.
\tcblower
\textbf{Solution détaillée — Modèle de rédaction.}

\emph{Plan :} (1) Projet réel pour l'année ; (2) Rêve familial ; (3) Regret sur le passé ; (4) Conclusion / Bonne résolution.

\emph{Rédaction :}

\eng{This year, if I work hard at school, my parents will be very proud of me. Our family is not rich, but if we had more money, we would build a bigger house in Yaoundé and my mother would stop selling at the market. Unfortunately, last year I was lazy: if I had listened to my teachers, I would have passed all my exams. Now I have decided to change, and unless I am ill, I will revise every evening.}

\emph{Traduction :} « Cette année, si je travaille dur à l'école, mes parents seront très fiers de moi. Notre famille n'est pas riche, mais si nous avions plus d'argent, nous construirions une maison plus grande à Yaoundé et ma mère cesserait de vendre au marché. Malheureusement, l'année dernière j'ai été paresseux : si j'avais écouté mes professeurs, j'aurais réussi tous mes examens. Maintenant j'ai décidé de changer, et à moins d'être malade, je réviserai chaque soir. »

\emph{Justification grammaticale :}
\begin{itemize}
\item \textbf{Type 1 :} \eng{if I work... will be} (Present Simple / Will) ; \eng{unless I am ill... will revise} (\eng{Unless} + Present / Will).
\item \textbf{Type 2 :} \eng{if we had... would build / would stop} (Past Simple / Would + BV). \eng{were} non utilisé ici car \eng{have} est verbe lexical.
\item \textbf{Type 3 :} \eng{if I had listened... would have passed} (Past Perfect / Would Have + V3). \eng{listen} régulier $\rightarrow$ \eng{listened}.
\item \textbf{Ponctuation :} Chaque proposition introduite par \eng{If/Unless} en début de phrase est suivie d'une virgule.
\end{itemize}
\end{exoresolu}

\begin{exercice}[Entraînement 1 : Complétion guidée]
Complete the sentences with the correct form of the verbs in brackets. Identify the conditional type (1, 2 or 3) for each.
\begin{enumerate}
\item If it \ldots\ldots\ldots (rain) tomorrow, we \ldots\ldots\ldots (cancel) the picnic. [Type: \ldots]
\item If I \ldots\ldots\ldots (be) a bird, I \ldots\ldots\ldots (fly) to Bamenda. [Type: \ldots]
\item She \ldots\ldots\ldots (not / fail) the exam if she \ldots\ldots\ldots (study) harder last term. [Type: \ldots]
\item \ldots\ldots\ldots (call) the police if you \ldots\ldots\ldots (see) a thief. [Type: \ldots]
\item If my father \ldots\ldots\ldots (not / lose) his job, we \ldots\ldots\ldots (buy) a new car last month. [Type: \ldots]
\end{enumerate}
\end{exercice}

\begin{exercice}[Entraînement 2 : Transformation]
Rewrite the sentences using the word in brackets. Do not change the meaning.
\begin{enumerate}
\item We didn't go to the beach because it rained. (If / would have)
\item I don't have a laptop, so I can't work from home. (If / would)
\item You must hurry, otherwise you will miss the train. (Unless)
\item He didn't take his umbrella, so he got wet. (If / would not have)
\item She is not here, so she cannot help us. (If / could)
\end{enumerate}
\end{exercice}

\begin{exercice}[Entraînement 3 : Traduction (Français $\rightarrow$ Anglais)]
Translate into English.
\begin{enumerate}
\item Si tu viens me voir samedi, nous irons au stade ensemble.
\item Si j'étais toi, j'accepterais cette offre d'emploi à Douala.
\item Si nous avions pris le petit-déjeuner, nous n'aurions pas faim maintenant. (Attention : conditionnel mixte non au programme, mais ici passé $\rightarrow$ présent. \textit{Indication : utiliser Type 3 pour la condition, Type 2 pour le résultat si autorisé, sinon reformuler. Pour le programme strict 1ère A, on évite les mixtes. Proposer plutôt :} « Si nous avions pris le petit-déjeuner, nous n'aurions pas eu faim à ce moment-là. »)
\item À moins qu'il ne pleuve, le match aura lieu.
\item Si les élèves avaient révisé, ils auraient réussi le Probatoire.
\end{enumerate}
\end{exercice}

\begin{corrige}[Exercice 1]
\begin{enumerate}
\item \eng{rains / will cancel} (Type 1). \eng{If it \textbf{rains} tomorrow, we \textbf{will cancel} the picnic.}
\item \eng{were / would fly} (Type 2). \eng{If I \textbf{were} a bird, I \textbf{would fly} to Bamenda.}
\item \eng{would not have failed / had studied} (Type 3). \eng{She \textbf{would not have failed} the exam if she \textbf{had studied} harder last term.}
\item \eng{Call / see} (Type 1 - Impératif). \eng{\textbf{Call} the police if you \textbf{see} a thief.}
\item \eng{had not lost / would have bought} (Type 3). \eng{If my father \textbf{hadn't lost} his job, we \textbf{would have bought} a new car last month.}
\end{enumerate}
\end{corrige}

\begin{corrige}[Exercice 2]
\begin{enumerate}
\item \eng{If it \textbf{hadn't rained}, we \textbf{would have gone} to the beach.}
\item \eng{If I \textbf{had} a laptop, I \textbf{would} work from home.} (ou \eng{could work})
\item \eng{\textbf{Unless you hurry}, you \textbf{will miss} the train.}
\item \eng{If he \textbf{had taken} his umbrella, he \textbf{wouldn't have got} wet.}
\item \eng{If she \textbf{were} here, she \textbf{could help} us.}
\end{enumerate}
\end{corrige}

\begin{corrige}[Exercice 3]
\begin{enumerate}
\item \eng{If you \textbf{come} to see me on Saturday, we \textbf{will go} to the stadium together.}
\item \eng{If I \textbf{were} you, I \textbf{would accept} this job offer in Douala.}
\item \eng{If we \textbf{had had} breakfast, we \textbf{would not have been} hungry then.} (Version passée cohérente Type 3).
\item \eng{\textbf{Unless it rains}, the match \textbf{will take place}.} (ou \eng{will be played}).
\item \eng{If the students \textbf{had revised}, they \textbf{would have passed} the Probatoire.}
\end{enumerate}
\end{corrige}

\newpage

\coursec{Exercices}

\courssub{Je vérifie mes connaissances}

\begin{exercice}[QCM : la bonne forme]
Entoure la forme correcte dans chaque phrase.
\begin{enumerate}
\item If I were you, I (will / would / would have) apologise to your sister.
\item If it rains tomorrow, we (stay / will stay / would stay) at home.
\item If she had studied harder, she (will pass / would pass / would have passed) her examination.
\item If my father (has / will have / would have) time, he will take us to Douala.
\item If they had left earlier, they (would not miss / will not miss / would not have missed) the bus to Bamenda.
\end{enumerate}
\end{exercice}

\begin{exercice}[Vrai ou faux ?]
Dis si chaque affirmation est vraie ou fausse, puis justifie ta réponse en une phrase.
\begin{enumerate}
\item In the sentence “If it rains, we will stay at home”, the condition is imaginary.
\item The second conditional expresses an unreal situation in the present or the future.
\item In the third conditional, we use \eng{would have} + past participle in the if-clause.
\item After \eng{if}, we can use \eng{will} in the first conditional: “If it will rain, we will stay home.”
\item “If I were you” is correct, even with the pronoun \eng{I}.
\end{enumerate}
\end{exercice}

\begin{exercice}[Vocabulaire : les mots de l'hypothèse]
Complète chaque phrase avec le mot qui convient : \eng{unless — provided — in case — wish — regret}.
\begin{enumerate}
\item I will lend you my bicycle \ldots\ldots\ldots you return it before Friday.
\item Take your umbrella \ldots\ldots\ldots it rains this afternoon.
\item \ldots\ldots\ldots you work hard, you will not pass the Probatoire.
\item I \ldots\ldots\ldots I were taller so I could play basketball for the school team.
\item Many students \ldots\ldots\ldots not listening to their teachers' advice last year.
\end{enumerate}
\end{exercice}

\begin{exercice}[Conjugaison à compléter]
Mets le verbe entre parenthèses à la forme correcte (first, second ou third conditional).
\begin{enumerate}
\item If she (work) harder, she (pass) her examination last year.
\item If we (win) the match tomorrow, the whole school (celebrate).
\item If I (be) rich, I (build) a big house for my grandmother in the village.
\item If you (not hurry), you (miss) the bendskin.
\item If they (listen) to the radio yesterday, they (hear) the good news.
\end{enumerate}
\end{exercice}

\courssub{Je m'entraîne}

\begin{exercice}[Traduction]
Traduis les phrases suivantes en anglais.
\begin{enumerate}
\item Si mes parents avaient eu de l'argent, ils m'auraient envoyé à l'université.
\item Si tu viens à la fête de ma sœur, tu rencontreras toute ma famille.
\item Si j'étais toi, je parlerais à mon oncle avant de prendre une décision.
\item S'il avait plu hier, nous ne serions pas allés au marché de Mokolo.
\item Si mon frère travaille dur, il réussira son examen en juin.
\end{enumerate}
\end{exercice}

\begin{exercice}[Transformation]
Réécris chaque phrase en commençant par \eng{If}, sans changer le sens.
\begin{enumerate}
\item I didn't listen to my parents, so I made a mistake. \par
\emph{If I \ldots}
\item She was ill, so she didn't come to the family meeting. \par
\emph{If she \ldots}
\item You don't save money, so you can't buy a new phone. \par
\emph{If you \ldots}
\item He didn't wear a helmet, so he was injured in the bendskin accident. \par
\emph{If he \ldots}
\item Study hard and you will succeed. \par
\emph{If you \ldots}
\end{enumerate}
\end{exercice}

\begin{exercice}[Texte à trous]
Complète le texte suivant avec la forme correcte des verbes entre parenthèses.
\begin{textebox}[A letter to my cousin]
Dear Nadia,

Thank you for your letter. Life in Yaoundé is not easy at the moment. If I (have) \ldots\ldots\ldots more money, I (buy) \ldots\ldots\ldots all my school books before the end of the month. Last term, if I (revise) \ldots\ldots\ldots my lessons every evening, I (get) \ldots\ldots\ldots better marks in English. But I have decided to change. If I (work) \ldots\ldots\ldots seriously this year, I am sure I (pass) \ldots\ldots\ldots the Probatoire. If only I (be) \ldots\ldots\ldots as disciplined as you!

Write soon.

Your cousin, Alain
\end{textebox}
\end{exercice}

\begin{exercice}[Appariement]
Relie chaque début de phrase (1 à 6) à sa fin logique (a à f) en respectant la concordance des temps.
\begin{center}
\begin{tabularx}{\linewidth}{|X|X|}
\hline
\rowcolor{popblueL} Début & Fin \\
\hline
1. If my uncle comes from Buea next week, & a. she would have helped you with your homework. \\
\hline
2. If I were the head of my family, & b. we will visit our grandmother in the village. \\
\hline
3. If you had asked your sister, & c. they will miss the family reunion. \\
\hline
4. If we had more free time, & d. I would listen to everybody's problems. \\
\hline
5. If they don't leave Douala early, & e. he would understand our traditions better. \\
\hline
6. If my cousin spoke Ewondo, & f. we would visit our relatives more often. \\
\hline
\end{tabularx}
\end{center}
\end{exercice}

\courssub{Je me prépare au Bac}

\begin{exercice}[Compréhension et langue — type Bac]
Lis le texte suivant, puis réponds aux questions.
\begin{textebox}[A family decision]
When Ekane finished his GCE Advanced Level in Buea, his family held a long meeting. His father wanted him to study law in Yaoundé, but Ekane dreamed of becoming an engineer. “If you study law, you will find a good job quickly,” his father said. His mother disagreed: “If our son were free to choose, he would be much happier.” Ekane himself had regrets. “If I had worked harder in physics last year, I would have obtained a better grade,” he admitted. His elder sister, who had refused her parents' advice five years earlier, added: “If I had listened to you, I would not have wasted two years at home. Ekane must follow his dream, but he must also work hard.” In the end, the family agreed: if Ekane passes the entrance examination, he will study engineering in Douala. If he fails, he will accept his father's plan without complaining.
\end{textebox}
\begin{enumerate}
\item Answer these questions in your own words:
\begin{enumerate}
\item What did Ekane's father want him to study?
\item Why did Ekane's sister regret her past decision?
\item What will happen if Ekane passes the entrance examination?
\end{enumerate}
\item Language work:
\begin{enumerate}
\item Find in the text one example of each conditional (first, second, third) and copy it.
\item In “If our son were free to choose, he would be much happier”, explain why the verb \eng{were} is used with \eng{our son}.
\item Rewrite: “Ekane didn't work hard in physics, so he didn't get a good grade.” Begin with: \emph{If Ekane \ldots}
\item Correct the mistake: “If Ekane will fail, he will accept his father's plan.”
\end{enumerate}
\end{enumerate}
\end{exercice}

\begin{exercice}[Traduction — type Bac]
Traduis le passage suivant en anglais.
\begin{textebox}[Texte à traduire]
Si mon père avait accepté le travail à Douala, toute la famille aurait déménagé. Nous aurions quitté notre village et j'aurais perdu tous mes amis. Heureusement, il a refusé. Si je termine mes études, j'aiderai mes parents et je construirai une maison pour eux. Si j'étais riche, je paierais aussi les études de ma petite sœur, car elle est très intelligente.
\end{textebox}
\end{exercice}

\begin{exercice}[Composition — type Bac]
Write a paragraph of about 80 to 100 words entitled \emph{“If I were the head of my family”}. In your paragraph:
\begin{itemize}
\item use at least \textbf{two second conditionals} to describe the changes you would make;
\item use at least \textbf{one first conditional} to announce a real plan for the future;
\item use at least \textbf{one third conditional} to express a regret about your family's past;
\item underline each conditional sentence you use.
\end{itemize}
\end{exercice}

\newpage

\coursec{Corrigés des exercices}

\courssub{Corrigés des exercices}

\begin{corrige}[Exercice 1 — QCM : la bonne forme]
\textbf{Réponses correctes :}
\begin{enumerate}
\item \eng{If I were you, I \textbf{would} apologise to your sister.} — Second conditional : \eng{if} + past simple, \eng{would} + base verbale. La structure figée \eng{If I were you} (conseil) exige toujours \eng{would}.
\item \eng{If it rains tomorrow, we \textbf{will stay} at home.} — First conditional : \eng{if} + present simple, \eng{will} + base verbale. La condition est réelle (demain).
\item \eng{If she had studied harder, she \textbf{would have passed} her examination.} — Third conditional : \eng{if} + past perfect, \eng{would have} + participe passé.
\item \eng{If my father \textbf{has} time, he will take us to Douala.} — First conditional : après \eng{if}, on utilise le present simple, jamais \eng{will}.
\item \eng{If they had left earlier, they \textbf{would not have missed} the bus to Bamenda.} — Third conditional à la forme négative : \eng{would not have} + participe passé (\eng{missed}).
\end{enumerate}
\textbf{Méthode :} repère d'abord le temps de la subordonnée en \eng{if} (present, past simple ou past perfect) ; il détermine automatiquement la forme de la proposition principale.
\end{corrige}

\begin{corrige}[Exercice 2 — Vrai ou faux ?]
\begin{enumerate}
\item \textbf{Faux.} \eng{In the first conditional, the condition is real or possible: it may really rain tomorrow.} (Au first conditional, l'hypothèse est réelle ; c'est au second conditional qu'elle est imaginaire.)
\item \textbf{Vrai.} \eng{The second conditional (if + past simple, would + verb) describes an unreal or imaginary situation in the present or the future.} Exemple : \eng{If I had a car, I would drive to Kribi.}
\item \textbf{Faux.} \eng{In the third conditional, “would have + past participle” is used in the main clause; the if-clause takes the past perfect (“had + past participle”).} On dit \eng{If she \textbf{had studied}, she \textbf{would have passed}}.
\item \textbf{Faux.} \eng{After “if”, we never use “will” to express a condition; we use the present simple: “If it rains, we will stay home.”} C'est l'erreur la plus fréquente des francophones, par calque de « s'il pleuvra ».
\item \textbf{Vrai.} \eng{In the second conditional, “were” is used with all persons in formal English: “If I were you”, “If she were here”.} (\eng{Was} se tolère à l'oral familier, mais \eng{were} est la forme correcte et exigée à l'examen.)
\end{enumerate}
\end{corrige}

\begin{corrige}[Exercice 3 — Vocabulaire : les mots de l'hypothèse]
\begin{enumerate}
\item \eng{I will lend you my bicycle \textbf{provided} you return it before Friday.} — \eng{provided (that)} = « à condition que » ; suivi du present simple.
\item \eng{Take your umbrella \textbf{in case} it rains this afternoon.} — \eng{in case} = « au cas où » ; suivi du present simple, jamais de \eng{will}.
\item \eng{\textbf{Unless} you work hard, you will not pass the Probatoire.} — \eng{unless} = « à moins que » (= \eng{if \ldots not}) ; attention, la phrase est déjà négative dans la principale.
\item \eng{I \textbf{wish} I were taller so I could play basketball for the school team.} — \eng{wish} + past simple exprime un regret sur le présent ; d'où \eng{were}.
\item \eng{Many students \textbf{regret} not listening to their teachers' advice last year.} — \eng{regret} + \emph{-ing} = « regretter d'avoir fait ».
\end{enumerate}
\textbf{Point de vigilance :} ne confonds pas \eng{in case} (« au cas où », précaution) et \eng{if} (condition) : \eng{Take an umbrella in case it rains} ≠ \eng{I will stay home if it rains}.
\end{corrige}

\begin{corrige}[Exercice 4 — Conjugaison à compléter]
\begin{enumerate}
\item \eng{If she \textbf{had worked} harder, she \textbf{would have passed} her examination last year.} — Third conditional (regret sur le passé, indiqué par \eng{last year}).
\item \eng{If we \textbf{win} the match tomorrow, the whole school \textbf{will celebrate}.} — First conditional (hypothèse réelle, \eng{tomorrow}).
\item \eng{If I \textbf{were} rich, I \textbf{would build} a big house for my grandmother in the village.} — Second conditional (irréel du présent ; \eng{were} pour toutes les personnes).
\item \eng{If you \textbf{do not hurry}, you \textbf{will miss} the bendskin.} — First conditional à la forme négative : present simple \eng{don't hurry} après \eng{if}.
\item \eng{If they \textbf{had listened} to the radio yesterday, they \textbf{would have heard} the good news.} — Third conditional ; participe passé irrégulier : \eng{hear — heard — heard}.
\end{enumerate}
\textbf{Rappel des trois structures :}
\begin{center}
\begin{tabular}{|l|l|l|}
\hline
\rowcolor{popblueL} Type & If-clause & Main clause \\
\hline
First & present simple & will + verb \\
\hline
Second & past simple & would + verb \\
\hline
Third & past perfect & would have + past participle \\
\hline
\end{tabular}
\end{center}
\end{corrige}

\begin{corrige}[Exercice 5 — Traduction]
\begin{enumerate}
\item \eng{If my parents had had money, they would have sent me to university.} — Third conditional. Note : \eng{had had} (auxiliaire \eng{had} + participe \eng{had}) est correct, bien qu'il surprenne les francophones.
\item \eng{If you come to my sister's party, you will meet all my family.} — First conditional ; pas de \eng{will} après \eng{if}.
\item \eng{If I were you, I would talk to my uncle before making a decision.} — Structure figée de conseil ; \eng{before} + \emph{-ing} (« avant de »).
\item \eng{If it had rained yesterday, we would not have gone to the Mokolo market.} — Third conditional négatif ; participe passé irrégulier : \eng{go — went — gone}.
\item \eng{If my brother works hard, he will pass his examination in June.} — First conditional ; « réussir un examen » se dit \eng{pass an examination} (faux ami : \eng{succeed} ne prend pas de complément direct).
\end{enumerate}
\textbf{Erreurs à éviter :} \eng{*If my parents would have had money\ldots} (jamais \eng{would} dans la subordonnée) ; \eng{*If you will come\ldots} (jamais \eng{will} après \eng{if}).
\end{corrige}

\begin{corrige}[Exercice 6 — Transformation]
\begin{enumerate}
\item \eng{If I \textbf{had listened to my parents, I would not have made} a mistake.} — Third conditional : les deux faits sont passés et réels, donc l'hypothèse est contraire à la réalité.
\item \eng{If she \textbf{had not been ill, she would have come} to the family meeting.} — Third conditional négatif dans la subordonnée ; participe passé : \eng{come — came — come}.
\item \eng{If you \textbf{saved money, you could buy} a new phone.} — Second conditional (situation présente irréelle) ; \eng{could} = « pourrais ».
\item \eng{If he \textbf{had worn} a helmet, he \textbf{would not have been injured} in the bendskin accident.} — Third conditional ; participe passé irrégulier : \eng{wear — wore — worn}.
\item \eng{If you \textbf{study hard, you will succeed}.} — First conditional : l'impératif + \eng{and} équivaut à une condition réelle.
\end{enumerate}
\textbf{Méthode :} demande-toi si les faits de la phrase de départ sont passés (→ third conditional) ou présents (→ second conditional), puis inverse la polarité : fait réel affirmatif → hypothèse négative, et inversement.
\end{corrige}

\begin{corrige}[Exercice 7 — Texte à trous]
Texte complété :
\begin{enumerate}
\item \eng{If I \textbf{had} more money, I \textbf{would buy} all my school books before the end of the month.} — Second conditional : Alain n'a pas assez d'argent maintenant (irréel du présent).
\item \eng{Last term, if I \textbf{had revised} my lessons every evening, I \textbf{would have got} better marks in English.} — Third conditional : regret sur le trimestre passé (\eng{last term}).
\item \eng{If I \textbf{work} seriously this year, I am sure I \textbf{will pass} the Probatoire.} — First conditional : projet réel pour cette année.
\item \eng{If only I \textbf{were} as disciplined as you!} — \eng{If only} + past simple exprime un souhait irréel au présent, comme \eng{wish} ; \eng{were} pour la première personne.
\end{enumerate}
\textbf{Point de vigilance :} les indices temporels (\eng{at the moment}, \eng{last term}, \eng{this year}) sont les clés qui permettent de choisir le bon conditional. Souligne-les avant de conjuguer.
\end{corrige}

\begin{corrige}[Exercice 8 — Appariement]
\textbf{Correspondances :} 1 — b ; 2 — d ; 3 — a ; 4 — f ; 5 — c ; 6 — e.
\begin{enumerate}
\item \eng{If my uncle comes from Buea next week, \textbf{we will visit our grandmother in the village.}} — First conditional : present + \eng{will}.
\item \eng{If I were the head of my family, \textbf{I would listen to everybody's problems.}} — Second conditional : past simple + \eng{would}.
\item \eng{If you had asked your sister, \textbf{she would have helped you with your homework.}} — Third conditional : past perfect + \eng{would have helped}.
\item \eng{If we had more free time, \textbf{we would visit our relatives more often.}} — Second conditional.
\item \eng{If they don't leave Douala early, \textbf{they will miss the family reunion.}} — First conditional négatif.
\item \eng{If my cousin spoke Ewondo, \textbf{he would understand our traditions better.}} — Second conditional.
\end{enumerate}
\textbf{Vérification :} seules ces associations respectent la concordance des temps ; par exemple, 1 (present) ne peut pas aller avec a (\eng{would have helped}), qui exige un past perfect dans la subordonnée.
\end{corrige}

\begin{corrige}[Exercice 9 — Compréhension et langue — type Bac]
\textbf{1. Comprehension questions (réponses modèles) :}
\begin{enumerate}
\item \eng{Ekane's father wanted him to study law in Yaoundé, because he believed it would help him find a good job quickly.}
\item \eng{She regretted her decision because she had refused her parents' advice and, as a result, she wasted two years at home.}
\item \eng{If Ekane passes the entrance examination, he will study engineering in Douala.}
\end{enumerate}
\textbf{2. Language work :}
\begin{enumerate}
\item Exemples attendus :
\begin{itemize}
\item First conditional : \eng{“If you study law, you will find a good job quickly.”} (ou \eng{“If Ekane passes the entrance examination, he will study engineering in Douala”}).
\item Second conditional : \eng{“If our son were free to choose, he would be much happier.”}
\item Third conditional : \eng{“If I had worked harder in physics last year, I would have obtained a better grade.”} (ou la phrase de la sœur).
\end{itemize}
\item \eng{“Were” is used because, in the second conditional, the verb “to be” takes the form “were” with all persons in careful English. The situation is unreal: the son is not free to choose.}
\item \eng{If Ekane \textbf{had worked hard in physics, he would have got a good grade.}} — Third conditional, polarité inversée.
\item Correction : \eng{“If Ekane \textbf{fails}, he will accept his father's plan.”} — Après \eng{if}, on n'emploie jamais \eng{will} pour exprimer la condition ; le present simple remplace le futur.
\end{enumerate}
\textbf{Barème indicatif (sur 20) :} compréhension 6 pts (2 par réponse correcte et grammaticale) ; relevé des trois conditionnels 6 pts ; explication de \eng{were} 2 pts ; transformation 3 pts ; correction de l'erreur 3 pts.
\textbf{Points de vigilance :} répondre en phrases complètes et dans ses propres mots (ne pas recopier le texte) ; vérifier la concordance des temps dans chaque phrase copiée ; soigner la forme négative \eng{would not have wasted}.
\end{corrige}

\begin{corrige}[Exercice 10 — Traduction — type Bac]
\textbf{Traduction modèle :}

\eng{If my father had accepted the job in Douala, the whole family would have moved. We would have left our village and I would have lost all my friends. Fortunately, he refused. If I finish my studies, I will help my parents and I will build a house for them. If I were rich, I would also pay for my little sister's education, because she is very intelligent.}

\textbf{Justifications grammaticales :}
\begin{itemize}
\item « Si mon père avait accepté\ldots\ aurait déménagé » → third conditional : \eng{had accepted} / \eng{would have moved}.
\item « j'aurais perdu » → \eng{I would have lost} ; participe passé irrégulier : \eng{lose — lost — lost}.
\item « Si je termine\ldots\ j'aiderai » → first conditional : \eng{If I finish\ldots\ I will help} ; pas de \eng{will} après \eng{if}.
\item « Si j'étais riche » → \eng{If I were rich} (second conditional, \eng{were} obligatoire à l'écrit soigné).
\item « payer les études de quelqu'un » → \eng{pay for somebody's education} (préposition \eng{for} obligatoire).
\end{itemize}
\textbf{Barème indicatif (sur 10) :} 2 pts par phrase correctement traduite (structure du conditional 1 pt, lexique et formes verbales 1 pt).
\textbf{Points de vigilance :} ne pas écrire \eng{*If my father would have accepted} ; ne pas oublier l'article dans \eng{the whole family} ; « heureusement » = \eng{fortunately} (faux ami : \eng{happily} = « joyeusement »).
\end{corrige}

\begin{corrige}[Exercice 11 — Composition — type Bac]
\textbf{Proposition de rédaction modèle (environ 95 mots) :}

\begin{textebox}[If I were the head of my family]
\underline{If I were the head of my family, I would listen to everybody before making any decision.} \underline{I would also encourage my younger brothers and sisters to work hard at school, because education is the key to success.} I know that my parents had difficulties in the past: \underline{if they had received more help from their relatives, they would have finished their studies.} For the future, I have a real plan: \underline{if I pass the Probatoire next year, I will organise a big celebration for the whole family.} Being the head of a family is a serious responsibility, and I would try to be fair with everyone.
\end{textebox}

\textbf{Analyse du modèle :}
\begin{itemize}
\item Un second conditional avec deux conséquences coordonnées : \eng{If I were\ldots, I would listen\ldots\ and I would also encourage\ldots} (irréel du présent, \eng{were} + \eng{would}).
\item Un first conditional : \eng{if I pass\ldots\ I will organise} (projet réel, present + \eng{will}).
\item Un third conditional : \eng{if they had received\ldots\ they would have finished} (regret sur le passé).
\end{itemize}
\textbf{Barème indicatif (sur 20) :} respect du sujet et de la longueur 4 pts ; présence et exactitude des quatre formes conditionnelles exigées 8 pts (2 pts chacune) ; correction grammaticale générale 4 pts ; vocabulaire et cohérence 4 pts.
\textbf{Points de vigilance :} souligner chaque phrase conditionnelle comme demandé ; ne pas écrire \eng{*If I would be} ; varier les sujets (\eng{I}, \eng{they}) ; vérifier chaque participe passé ; conclure le paragraphe par une phrase de bilan.
\end{corrige}

\newpage

\coursec{Fiche bilan}

\begin{popfigure}
\begin{tikzpicture}[mindmap, grow cyclic, every node/.style={concept}, text=white,
root concept/.append style={concept color=popdark, font=\bfseries, minimum size=3.2cm},
level 1 concept/.append style={sibling angle=60, level distance=3.6cm, font=\small, minimum size=2.2cm}]
\node[root concept]{Conditionals 1, 2, 3}
child[concept color=popgreen]{ node[align=center]{First : réel futur\\ If + present, will + BV} }
child[concept color=popblue]{ node[align=center]{Second : irréel présent\\ If + past, would + BV} }
child[concept color=poporange]{ node[align=center]{Third : regret passé\\ If + past perfect, would have + p.p.} }
child[concept color=poppurple]{ node[align=center]{Jamais will/would\\ dans la if-clause} }
child[concept color=poppink]{ node[align=center]{Variantes : unless,\\ provided, in case} }
child[concept color=popgold]{ node[align=center]{Emplois : conseil,\\ souhait, reproche} };
\end{tikzpicture}
\legende{Carte mentale : les trois conditionnels et leurs pièges.}
\end{popfigure}

\begin{bilan}
\textbf{Lexique clé (English / Français)}
\begin{itemize}
\item \eng{unless} = à moins que / sauf si ; \eng{provided (that)} = à condition que ; \eng{in case} = au cas où
\item \eng{if I were you} = si j'étais toi ; \eng{I wish...} = si seulement... ; \eng{as long as} = tant que
\item \eng{to regret} = regretter ; \eng{a piece of advice} = un conseil ; \eng{otherwise} = sinon
\end{itemize}

\textbf{Règles à connaître par cœur}
\begin{center}
\begin{tabularx}{\linewidth}{|l|X|X|}\hline
\rowcolor{popblueL} Type & Structure & Exemple \\ \hline
First & If + present simple, will + BV & \eng{If it rains, we will stay home.} \\ \hline
Second & If + past simple, would + BV & \eng{If I had money, I would travel.} \\ \hline
Third & If + past perfect, would have + p.p. & \eng{If she had studied, she would have passed.} \\ \hline
\end{tabularx}
\end{center}

\begin{itemize}
\item \textbf{Jamais} \eng{will} ni \eng{would} dans la proposition introduite par \eng{if}.
\item \eng{Were} pour toutes les personnes au second conditional : \eng{If I were rich...}
\item La virgule se met quand la if-clause est en tête de phrase.
\item \eng{Unless} = \eng{if... not} : \eng{Unless you work, you will fail.} = \eng{If you do not work, you will fail.}
\item \textbf{Méthode express} : 1) repérer le temps après \eng{if} ; 2) en déduire le type ; 3) construire la proposition principale avec l'auxiliaire correspondant.
\end{itemize}
\end{bilan}

\begin{aretenir}[Checklist avant le Bac]
\begin{itemize}
\item[$\square$] Je connais la structure du first conditional (if + present, will + BV).
\item[$\square$] Je connais la structure du second conditional (if + past, would + BV).
\item[$\square$] Je connais la structure du third conditional (if + past perfect, would have + participe passé).
\item[$\square$] Je ne mets jamais \eng{will} ou \eng{would} dans la if-clause.
\item[$\square$] J'utilise \eng{were} avec \eng{I/he/she} au second conditional.
\item[$\square$] Je sais transformer une phrase avec \eng{unless} en phrase avec \eng{if... not}.
\item[$\square$] Je sais exprimer un conseil : \eng{If I were you, I would...}
\item[$\square$] Je sais exprimer un regret : \eng{If I had worked harder, I would have succeeded.}
\item[$\square$] Je place correctement la virgule quand la if-clause ouvre la phrase.
\item[$\square$] Je distingue hypothèse réelle (first) et hypothèse imaginaire (second) selon le contexte.
\item[$\square$] Je connais les verbes irréguliers fréquents au participe passé (\eng{gone, been, done, written, taken}).
\item[$\square$] Je sais rédiger un court paragraphe mêlant les trois conditionnels sans faute d'accord des temps.
\end{itemize}
\end{aretenir}

\findecours

\end{document}
